\documentclass[11pt]{amsart}

\usepackage[margin=1.1in]{geometry}
\usepackage{amsmath,amssymb,amsthm,mathtools}
\usepackage{enumitem}
\usepackage{aliascnt}
\usepackage{microtype}
\usepackage{hyperref}
\usepackage[nameinlink,capitalize]{cleveref}

\hypersetup{colorlinks=true,linkcolor=blue,citecolor=blue,urlcolor=blue}
\raggedbottom

\theoremstyle{plain}
\newtheorem{theorem}{Theorem}[section]

\newaliascnt{lemma}{theorem}
\newtheorem{lemma}[lemma]{Lemma}
\aliascntresetthe{lemma}

\newaliascnt{proposition}{theorem}
\newtheorem{proposition}[proposition]{Proposition}
\aliascntresetthe{proposition}

\newaliascnt{corollary}{theorem}
\newtheorem{corollary}[corollary]{Corollary}
\aliascntresetthe{corollary}

\newaliascnt{hypothesis}{theorem}
\newtheorem{hypothesis}[hypothesis]{Hypothesis}
\aliascntresetthe{hypothesis}

\theoremstyle{definition}
\newaliascnt{definition}{theorem}
\newtheorem{definition}[definition]{Definition}
\aliascntresetthe{definition}

\theoremstyle{remark}
\newaliascnt{remark}{theorem}
\newtheorem{remark}[remark]{Remark}
\aliascntresetthe{remark}

\crefname{theorem}{Theorem}{Theorems}
\Crefname{theorem}{Theorem}{Theorems}
\crefname{lemma}{Lemma}{Lemmas}
\Crefname{lemma}{Lemma}{Lemmas}
\crefname{proposition}{Proposition}{Propositions}
\Crefname{proposition}{Proposition}{Propositions}
\crefname{corollary}{Corollary}{Corollaries}
\Crefname{corollary}{Corollary}{Corollaries}
\crefname{hypothesis}{Hypothesis}{Hypotheses}
\Crefname{hypothesis}{Hypothesis}{Hypotheses}
\crefname{definition}{Definition}{Definitions}
\Crefname{definition}{Definition}{Definitions}
\crefname{remark}{Remark}{Remarks}
\Crefname{remark}{Remark}{Remarks}

\numberwithin{equation}{section}

\newcommand{\R}{\mathbb R}
\newcommand{\N}{\mathbb N}
\newcommand{\eps}{\varepsilon}
\newcommand{\loc}{\mathrm{loc}}
\newcommand{\supp}{\operatorname{supp}}
\newcommand{\divergence}{\operatorname{div}}
\newcommand{\dist}{\operatorname{dist}}
\newcommand{\esssup}{\operatorname*{ess\,sup}}
\newcommand{\calS}{\mathcal S}
\newcommand{\calR}{\mathcal R}
\newcommand{\calC}{\mathcal C}
\newcommand{\calU}{\mathcal U}
\newcommand{\calH}{\mathcal H}
\DeclareMathOperator{\diver}{div}

\title[Conditional Local Type I Navier--Stokes Reduction]
{Conditional Local Type I Navier--Stokes Reduction: Concentration, Seregin Limits, and Liouville-Class Closures}

\author{Guillem Duran Ballester}
\email{guillem@fragile.tech}
\date{}

\subjclass[2020]{35Q30, 35B44, 35B40, 76D05}
\keywords{Navier--Stokes equations, Type I singularities, ancient solutions, Seregin limits, local concentration, Liouville theorems}

\begin{document}

\begin{abstract}
We prove the concentration, compactness, and Liouville-class reductions used in the
conditional local Type I Navier--Stokes reduction.  The paper contains the local
Caffarelli--Kohn--Nirenberg concentration theorem and Type I/Type II dichotomy,
the admissible Type I Seregin blow-up extraction, the decomposition of
normalized centered ancient limits into standard Liouville classes and a
residual class, the endpoint exclusion of the uniformly \(L^3\)-tight class,
and the structured ancient-solution exclusions covered by classical and
perturbative Liouville theorems.  The residual class is kept as the explicit
residual-class hypothesis \Cref{p1:hyp:no-remainder}.
\end{abstract}

\maketitle
\tableofcontents

\section*{Introduction and Structure of the Argument}

This paper presents a conditional local Type I reduction argument.
It reduces a local pointwise Type I singularity to a normalized, nonzero,
bounded centered ancient solution and excludes the non-residual Liouville
classes available from standard local compactness and known rigidity theorems.
The residual class is not closed in this paper; it is isolated as the residual-class
hypothesis \Cref{p1:hyp:no-remainder}.

The proof is organized in four steps.  First, the local concentration theorem
shows that a singular point carries positive Caffarelli--Kohn--Nirenberg
critical mass at arbitrarily small scales and gives the local Type I/local Type
II dichotomy.  It also identifies the endpoint weak-Serrin estimate used to
obtain terminal velocity no-escape for local pointwise Type I rescalings.
Second, the Type I blow-up analysis proves the admissible Seregin extraction
theorem: an admissible finite-energy Type I singularity produces a smooth
bounded ancient solution of the centered self-similar equation with positive
compact local \(L^3\) velocity mass.  It then defines the Liouville classes and
the residual class.

Third, the uniformly tight analysis proves that no nonzero normalized Seregin
limit can be uniformly \(L^3\)-tight.  The key point is that boundedness plus
uniform \(L^3\)-tightness gives a backward sequence of finite critical norm
after physical pullback, so the Albritton--Barker endpoint ancient Liouville
theorem applies.  Fourth, the structured part proves the axisymmetric and
perturbative weighted-vorticity exclusions used in the Liouville-class analysis.
These results turn the corresponding structured hypotheses into the zero
ancient solution, contradicting the local nonvanishing obtained from the
Seregin extraction.

Consequently, within this paper a normalized Seregin limit generated by an
admissible Type I singularity must either belong to a resolved Liouville class or to
the residual class.  The Type I contradiction is therefore conditional on
the residual-class hypothesis \Cref{p1:hyp:no-remainder}.

\clearpage
\section*{Part I. Local Concentration and the Type I/Type II Dichotomy}

\subsection*{Overview}

We prove a local analytic reduction for possible finite-time singularities of
the three-dimensional Navier--Stokes equations.  For a suitable weak solution,
every singular point at time \(T\) has arbitrarily small backward parabolic
cylinders on which both the standard mean-subtracted
Caffarelli--Kohn--Nirenberg scale-invariant quantity \(C+D\) and the
scale-invariant velocity quantity \(C\) are bounded below by positive
constants.  If these local scale-invariant quantities tend to zero at each
point of the singular set, the solution is regular at time \(T\) by the
\(\varepsilon\)-regularity theorem.  Along any such positive concentration
sequence, the solution either satisfies a local Type I scale-invariant bound or
belongs to the local Type II alternative.  For local pointwise Type I terminal
singularities of suitable weak solutions, the endpoint weak-Serrin argument gives
the velocity no-escape condition
used in the Type I blow-up analysis.

\section{Introduction}

Let \((u,p)\) be a suitable weak solution of the three-dimensional
incompressible Navier--Stokes equations
\[
  \partial_t u+u\cdot\nabla u+\nabla p=\Delta u,
  \qquad \nabla\cdot u=0.
\]
We use the standard normalization in which the viscosity is one.  The weak
solution theory goes back to Leray \cite{Leray1934}.  Suitable weak solutions
and their partial regularity theory were introduced and developed by Scheffer
\cite{Scheffer1976} and by Caffarelli, Kohn, and Nirenberg \cite{CKN1982}; see
also Lin \cite{Lin1998}.  We establish a local reduction from a
possible singularity at time \(T\) to the Type I/Type II blow-up rate
dichotomy.

Scale-invariant regularity and blow-up criteria provide the standard background
for this formulation; see Serrin \cite{Serrin1962}, Giga \cite{Giga1986},
Escauriaza--Seregin--{\v S}verak \cite{ESS2003}, and Seregin
\cite{Seregin2012}.  The principal regularity theorem used below is the
Caffarelli--Kohn--Nirenberg \(\varepsilon\)-regularity theorem.

The proof proceeds in four steps.  First, the local scale-invariant
quantities \(C\) and \(D\) are finite for suitable weak solutions and invariant
under the usual parabolic scaling.  Second, the Caffarelli--Kohn--Nirenberg
\(\varepsilon\)-regularity theorem implies that vanishing of \(C+D\) at a
candidate singular point gives local regularity, while the velocity
\(\varepsilon\)-regularity criterion gives positive concentration in the
scale-invariant \(L^3_{x,t}\) quantity.  Third, a finite-time singularity
therefore produces sequences of shrinking cylinders on which \(C+D\), and
separately \(C\), are bounded below by positive constants.  Finally, once such
a sequence is fixed, it either satisfies a Type I scale-invariant bound or
belongs to the local Type II alternative.

The concentration and Type I/Type II dichotomy proofs are local in
spacetime.  Apart from the standard Caffarelli--Kohn--Nirenberg
\(\varepsilon\)-regularity theorem, all quantities and alternatives used in
that local analysis are defined below.  The local Type I reduction section
shows how the endpoint weak-Serrin estimate upgrades the local pointwise Type I
envelope to terminal velocity no-escape through the local endpoint weak-Serrin
Type I estimate.

\section{Suitable Weak Solutions and Singular Points}

\begin{definition}[Suitable weak solution]
Let \(I\subset\mathbb R\) be an open interval.  A pair \((u,p)\) is a suitable
weak solution on \(\mathbb R^3\times I\) if
\[
  u\in L^\infty_tL^2_{x,\mathrm{loc}}
  \cap L^2_tH^1_{x,\mathrm{loc}},
  \qquad
  p\in L^{3/2}_{\mathrm{loc}},
  \qquad
  \nabla\cdot u=0,
\]
the Navier--Stokes equations hold in the sense of distributions, and the local
energy inequality holds: for almost every \(t_1<t_2\) in \(I\) and every
nonnegative \(\phi\in C_c^\infty(\mathbb R^3\times I)\),
\[
\begin{aligned}
&\int_{\mathbb R^3}\frac{|u(x,t_2)|^2}{2}\phi(x,t_2)\,dx
 +\int_{t_1}^{t_2}\int_{\mathbb R^3}|\nabla u|^2\phi\,dx\,dt \\
&\le
 \int_{\mathbb R^3}\frac{|u(x,t_1)|^2}{2}\phi(x,t_1)\,dx
 +\int_{t_1}^{t_2}\int_{\mathbb R^3}
   \frac{|u|^2}{2}(\partial_t\phi+\Delta\phi)\,dx\,dt \\
&\quad
 +\int_{t_1}^{t_2}\int_{\mathbb R^3}
   \left(\frac{|u|^2}{2}+p\right)u\cdot\nabla\phi\,dx\,dt .
\end{aligned}
\]
The pressure is understood modulo functions of time.
\end{definition}

For \(z_0=(x_0,t_0)\) and \(r>0\), write
\[
  Q_r(z_0)=B_r(x_0)\times(t_0-r^2,t_0).
\]
We also write \(Q_r(x_0,t_0)\) for \(Q_r((x_0,t_0))\).

\begin{definition}[Singular set at a fixed time]
For a time \(T\), define
\[
  \Sigma(T)=
  \{x\in\mathbb R^3:\ u\text{ is not locally bounded in any }
  Q_r(x,T)\}.
\]
Equivalently, \(x\notin\Sigma(T)\) if there are \(r>0\) and \(M<\infty\) such
that \(\|u\|_{L^\infty(Q_r(x,T))}\le M\).
\end{definition}

If a finite-time singularity occurs at time \(T\), then
\(\Sigma(T)\ne\emptyset\) in this local sense.

\section{Parabolic Rescaling and Scale-Invariant Quantities}

Fix \(z_0=(x_0,T)\).  For \(r>0\), define the parabolic rescaling
\[
  u^{(r)}(y,s)=r\,u(x_0+ry,T+r^2s),
  \qquad
  p^{(r)}(y,s)=r^2p(x_0+ry,T+r^2s).
\]
The rescaled variables are defined on the set of \((y,s)\) for which
\((x_0+ry,T+r^2s)\) belongs to the original spacetime domain.

\begin{proposition}[Invariance under parabolic rescaling]
If \((u,p)\) is a suitable weak solution near \((x_0,T)\), then
\((u^{(r)},p^{(r)})\) is a suitable weak solution on the rescaled domain.  The
local energy inequality is preserved, and the pressure remains defined modulo
functions of the rescaled time.
\end{proposition}

\begin{proof}
The distributional equations and the divergence condition follow from the
change of variables \(x=x_0+ry\), \(t=T+r^2s\).  We include the verification of
the local energy inequality to fix the normalization.

Let \(\psi\in C_c^\infty\) be nonnegative and compactly supported in the
rescaled domain.  Define
\[
  \phi(x,t)=
  \psi\left(\frac{x-x_0}{r},\frac{t-T}{r^2}\right).
\]
If \(y=(x-x_0)/r\) and \(s=(t-T)/r^2\), then
\[
  \partial_t\phi=r^{-2}(\partial_s\psi)(y,s),\qquad
  \nabla_x\phi=r^{-1}(\nabla_y\psi)(y,s),\qquad
  \Delta_x\phi=r^{-2}(\Delta_y\psi)(y,s).
\]
For almost every \(s_1<s_2\), the corresponding times
\(T+r^2s_1<T+r^2s_2\) are admissible in the original local energy inequality.
Apply that inequality for \((u,p)\) to \(\phi\) on
\([T+r^2s_1,T+r^2s_2]\).  After the change of variables
\(x=x_0+ry\), \(t=T+r^2s\), using
\[
  u(x_0+ry,T+r^2s)=r^{-1}u^{(r)}(y,s),\qquad
  p(x_0+ry,T+r^2s)=r^{-2}p^{(r)}(y,s),
\]
each term contains the same common factor \(r\).  Dividing by this factor gives
the rescaled inequality
\[
\begin{aligned}
&\int_{\mathbb R^3}\frac{|u^{(r)}(y,s_2)|^2}{2}\psi(y,s_2)\,dy
 +\int_{s_1}^{s_2}\int_{\mathbb R^3}|\nabla_y u^{(r)}|^2\psi\,dy\,ds \\
&\le
 \int_{\mathbb R^3}\frac{|u^{(r)}(y,s_1)|^2}{2}\psi(y,s_1)\,dy
 +\int_{s_1}^{s_2}\int_{\mathbb R^3}
   \frac{|u^{(r)}|^2}{2}(\partial_s\psi+\Delta_y\psi)\,dy\,ds \\
&\quad
 +\int_{s_1}^{s_2}\int_{\mathbb R^3}
   \left(\frac{|u^{(r)}|^2}{2}+p^{(r)}\right)
   u^{(r)}\cdot\nabla_y\psi\,dy\,ds .
\end{aligned}
\]
Thus \((u^{(r)},p^{(r)})\) is suitable on the rescaled domain.  The pressure
transforms according to the Navier--Stokes scaling, and addition of a function
of \(t\) becomes addition of a function of \(s\).
\end{proof}

\begin{definition}[Scale-invariant local quantities]
For \(z_0=(x_0,t_0)\) and \(r>0\), set
\[
  C(z_0,r)=r^{-2}\int_{Q_r(z_0)} |u|^3\,dx\,dt
\]
and
\[
  D(z_0,r)=r^{-2}
  \int_{t_0-r^2}^{t_0}\int_{B_r(x_0)}
  |p(x,t)-(p)_{B_r(x_0)}(t)|^{3/2}\,dx\,dt.
\]
Here \((p)_{B_r(x_0)}(t)\) denotes the spatial mean of \(p(\cdot,t)\) on
\(B_r(x_0)\).
\end{definition}

The subtraction of the spatial mean removes the pressure ambiguity in the form
used by the Caffarelli--Kohn--Nirenberg criterion.

\begin{proposition}[Local finiteness and scaling]\label{p0:prop:critical-scaling}
Let \((u,p)\) be a suitable weak solution near \(z_0=(x_0,T)\).  Then
\(C(z_0,r)\) and \(D(z_0,r)\) are finite for every sufficiently small \(r\).
Moreover, for every fixed \(R<\infty\),
\[
\begin{aligned}
&\int_{-R^2}^{0}\int_{B_R}
\left(
|u^{(r)}|^3+
|p^{(r)}-(p^{(r)})_{B_R}(s)|^{3/2}
\right)\,dy\,ds \\
&\qquad
=R^2\bigl(C(z_0,rR)+D(z_0,rR)\bigr).
\end{aligned}
\]
\end{proposition}

\begin{proof}
Local finiteness of the velocity term follows from
\[
  u\in L^\infty_tL^2_{x,\mathrm{loc}}
  \cap L^2_tH^1_{x,\mathrm{loc}}
  \subset L^3_{\mathrm{loc}}
\]
on finite cylinders.  Indeed, if \(B\Subset\mathbb R^3\) and
\((a,b)\Subset I\), Sobolev and interpolation give
\[
  \int_a^b \|u(t)\|_{L^3(B)}^3\,dt
  \le
  C
  \left(\operatorname*{ess\,sup}_{a<t<b}\|u(t)\|_{L^2(B)}\right)^{3/2}
  \left(\int_a^b \|u(t)\|_{H^1(B)}^2\,dt\right)^{3/4}
  (b-a)^{1/4}.
\]
Thus \(u\in L^3_{\mathrm{loc}}\).  The pressure term is finite because
\(p\in L^{3/2}_{\mathrm{loc}}\), and subtracting a spatial mean over the same
ball preserves local \(L^{3/2}\) integrability.  More explicitly, for a ball
\(B\),
\[
  |p-(p)_B(t)|^{3/2}\le C\left(|p|^{3/2}+|(p)_B(t)|^{3/2}\right),
\]
while Jensen's inequality gives
\[
  |(p)_B(t)|^{3/2}
  \le \frac1{|B|}\int_B |p(x,t)|^{3/2}\,dx.
\]
After integration over \(B\) and over time, the mean-subtracted pressure term is
controlled by the local \(L^{3/2}\) norm of \(p\).

For the scaling identity, use \(x=x_0+ry\), \(t=T+r^2s\).  The velocity term
transforms as
\[
  \int_{-R^2}^{0}\int_{B_R}|u^{(r)}|^3\,dy\,ds
  =r^{-2}\int_{T-(rR)^2}^{T}\int_{B_{rR}(x_0)}|u|^3\,dx\,dt.
\]
The pressure means transform according to
\[
  (p^{(r)})_{B_R}(s)=r^2(p)_{B_{rR}(x_0)}(T+r^2s),
\]
because
\[
  \frac1{|B_R|}\int_{B_R} r^2p(x_0+ry,T+r^2s)\,dy
  =
  \frac{r^2}{|B_{rR}|}\int_{B_{rR}(x_0)}p(x,T+r^2s)\,dx .
\]
Therefore
\[
\begin{aligned}
&\int_{-R^2}^{0}\int_{B_R}
|p^{(r)}-(p^{(r)})_{B_R}(s)|^{3/2}\,dy\,ds \\
&\qquad =
r^{-2}\int_{T-(rR)^2}^{T}\int_{B_{rR}(x_0)}
|p-(p)_{B_{rR}(x_0)}(t)|^{3/2}\,dx\,dt .
\end{aligned}
\]
Adding the velocity and pressure identities gives the displayed formula.
\end{proof}

\section{\texorpdfstring{\(\varepsilon\)-Regularity}{Epsilon-Regularity} and the Vanishing Case}

\begin{theorem}[Caffarelli--Kohn--Nirenberg \(\varepsilon\)-regularity]
\label{p0:thm:ckn}
There exists a universal constant \(\varepsilon_0>0\) with the
following property.  Let \((u,p)\) be a suitable weak solution in \(Q_r(z_0)\).
If
\[
  C(z_0,r)+D(z_0,r)<\varepsilon_0,
\]
then \(u\) is locally bounded in a smaller cylinder, for instance in
\(Q_{r/2}(z_0)\).  In particular, if \(z_0=(x_0,T)\), then
\(x_0\notin\Sigma(T)\).
\end{theorem}

\begin{proof}
This is the local \(\varepsilon\)-regularity theorem of
Caffarelli--Kohn--Nirenberg \cite{CKN1982}; see also Lin \cite{Lin1998}.  It is
stated here with the spatial mean subtraction used in \(D\).  Equivalent
representatives of the pressure differ by functions of time, and the spatial
mean subtraction on \(B_r(x_0)\) removes this ambiguity.
\end{proof}

\begin{theorem}[Velocity \(\varepsilon\)-regularity]\label{p0:thm:velocity-ckn}
There exists a universal constant \(\varepsilon_v>0\) with the following
property.  Let \((u,p)\) be a suitable weak solution in \(Q_r(z_0)\).  If
\[
  C(z_0,r)<\varepsilon_v,
\]
then \(u\) is locally bounded in \(Q_{r/2}(z_0)\).  In particular, if
\(z_0=(x_0,T)\), then \(x_0\notin\Sigma(T)\).
\end{theorem}

\begin{proof}
This is the standard \(L^3_{x,t}\) form of the Caffarelli--Kohn--Nirenberg
\(\varepsilon\)-regularity criterion.  It follows from the standard local
pressure decay lemma together with \cref{p0:thm:ckn}; equivalently, it is the
\(L^3_{x,t}\)-smallness criterion
proved in the Caffarelli--Kohn--Nirenberg theory and in Lin's direct proof
\cite{CKN1982,Lin1998}.  The statement is scale invariant because \(C\) is
invariant under the parabolic Navier--Stokes scaling.
\end{proof}

\begin{definition}[Vanishing local scale-invariant quantities at time \(T\)]
We say that the local scale-invariant quantities vanish at time \(T\) if, for
every \(x_0\in\Sigma(T)\),
\[
  \limsup_{r\downarrow0}
  \bigl(C((x_0,T),r)+D((x_0,T),r)\bigr)=0.
\]
This condition is vacuous when \(\Sigma(T)=\emptyset\).
\end{definition}

\begin{proposition}[Equivalent rescaled formulation]
\label{p0:prop:vanishing-rescaled}
The preceding condition is equivalent to the following statement: for every
\(x_0\in\Sigma(T)\) and every fixed \(R<\infty\),
\[
  \lim_{r\downarrow0}
  \int_{-R^2}^{0}\int_{B_R}
  \left(
  |u^{(r)}|^3+
  |p^{(r)}-(p^{(r)})_{B_R}(s)|^{3/2}
  \right)\,dy\,ds=0.
\]
\end{proposition}

\begin{proof}
By \cref{p0:prop:critical-scaling}, the rescaled integral over
\(B_R\times(-R^2,0)\) equals
\[
  R^2\bigl(C((x_0,T),rR)+D((x_0,T),rR)\bigr).
\]
Thus vanishing of \(C+D\) implies vanishing on every fixed rescaled cylinder.
Conversely, taking \(R=1\) gives the original pointwise condition.
\end{proof}

\begin{theorem}[Regularity under vanishing local quantities]\label{p0:thm:no-concentration}
Let \((u,p)\) be a suitable weak solution on
\(\mathbb R^3\times(T-\delta,T)\) for some \(\delta>0\).  If the local
scale-invariant quantities vanish at time \(T\), then
\(\Sigma(T)=\emptyset\).  Thus no finite-time singularity occurs at time \(T\)
in the no-concentration case.
\end{theorem}

\begin{proof}
Suppose, toward a contradiction, that \(x_0\in\Sigma(T)\).  By the vanishing
hypothesis, choose \(r>0\) such that
\[
  C((x_0,T),r)+D((x_0,T),r)<\varepsilon_0.
\]
\Cref{p0:thm:ckn} implies that \(u\) is bounded in \(Q_{r/2}(x_0,T)\), hence
\(x_0\notin\Sigma(T)\).  This contradicts the choice of \(x_0\).  Hence
\(\Sigma(T)=\emptyset\).
\end{proof}

\begin{corollary}[Vanishing local quantities exclude a terminal singularity]
\label{p0:cor:escape}
Let \((u,p)\) be suitable near \((x_0,T)\).  If
\[
  \limsup_{r\downarrow0}
  \bigl(C((x_0,T),r)+D((x_0,T),r)\bigr)=0,
\]
then \(u\) is bounded in some backward cylinder \(Q_\rho(x_0,T)\), and hence
\(x_0\notin\Sigma(T)\).  Equivalently, membership in \(\Sigma(T)\) requires
that \(C+D\) have a positive lower bound along a sequence of shrinking backward
parabolic cylinders.
\end{corollary}

\begin{proof}
The hypothesis gives a scale at which \(C+D\) is below the universal threshold
in \cref{p0:thm:ckn}.  The \(\varepsilon\)-regularity theorem gives local boundedness in
a backward cylinder ending at \(T\), so \(x_0\notin\Sigma(T)\).
\end{proof}

\section{Positive Local Concentration}

\begin{lemma}[Failure of vanishing gives a positive concentration sequence]
\label{p0:lem:positive-concentration}
Suppose the local scale-invariant quantities do not vanish at time \(T\).  Then
there exist \(x_0\in\Sigma(T)\), a constant \(\eta>0\), and radii
\(r_n\downarrow0\) such that
\[
  C((x_0,T),r_n)+D((x_0,T),r_n)\ge\eta
  \qquad\text{for all }n.
\]
Equivalently, the rescaled sequence around \((x_0,T)\) has a positive lower
bound for the corresponding local scale-invariant quantities on a fixed
bounded cylinder.
\end{lemma}

\begin{proof}
Negating the definition of vanishing gives a point \(x_0\in\Sigma(T)\) with
positive limsup of
\[
  C((x_0,T),r)+D((x_0,T),r)
\]
as \(r\downarrow0\).  Choose \(\eta>0\) below that limsup and choose a sequence
\(r_n\downarrow0\) along which the displayed lower bound holds.  The rescaled
formulation follows from \cref{p0:prop:critical-scaling}.
\end{proof}

\begin{definition}[Positive local concentration sequence]
\label{p0:def:positive-concentration-sequence}
Let \(x_0\in\Sigma(T)\).  A sequence \(r_n\downarrow0\) is a positive local
concentration sequence at \((x_0,T)\) if there is \(\eta>0\) such that
\[
  C((x_0,T),r_n)+D((x_0,T),r_n)\ge\eta
  \qquad\text{for all }n.
\]
The corresponding parabolic rescalings are those defined by
\[
  u_n(y,s)=r_nu(x_0+r_ny,T+r_n^2s),
  \qquad
  p_n(y,s)=r_n^2p(x_0+r_ny,T+r_n^2s).
\]
\end{definition}

\begin{proposition}[Persistence of positive concentration]
\label{p0:prop:compactness-persistence}
Let \((u_n,p_n)\) be suitable weak solutions on a fixed cylinder
\(B_R\times(-R^2,0)\).  Suppose that
\[
  \int_{-R^2}^{0}\int_{B_R}
  \left(|u_n|^3+|p_n-(p_n)_{B_R}(s)|^{3/2}\right)\,dy\,ds
  \ge\eta>0
\]
for all \(n\).  Suppose further that, after passing to a subsequence,
\[
  u_n\to U\quad\text{in }L^3(B_R\times(-R^2,0))
\]
and
\[
  p_n-(p_n)_{B_R}(s)\to \Pi
  \quad\text{in }L^{3/2}(B_R\times(-R^2,0)).
\]
Then
\[
  \int_{-R^2}^{0}\int_{B_R}
  \left(|U|^3+|\Pi|^{3/2}\right)\,dy\,ds\ge\eta.
\]
In particular, the limiting pair retains the same positive lower bound for the
mean-subtracted local scale-invariant quantity.  If
\(\Pi=P-(P)_{B_R}(s)\), this is the corresponding
Caffarelli--Kohn--Nirenberg quantity of \((U,P)\) on
\(B_R\times(-R^2,0)\).
\end{proposition}

\begin{proof}
Strong convergence in \(L^3\) gives
\[
  \int_{-R^2}^{0}\int_{B_R}|u_n|^3\,dy\,ds
  \longrightarrow
  \int_{-R^2}^{0}\int_{B_R}|U|^3\,dy\,ds.
\]
Strong convergence in \(L^{3/2}\) gives the analogous convergence of the
mean-subtracted pressure integrals.  Passing to the limit in the lower bound
yields the asserted inequality.
\end{proof}

\begin{remark}
\Cref{p0:prop:compactness-persistence} is an auxiliary compactness observation,
independent of the proof of the dichotomy below.
\end{remark}

\begin{proposition}[Local alternative]\label{p0:prop:local-alternative}
Let \((u,p)\) be a suitable weak solution on
\(\mathbb R^3\times(T-\delta,T)\) for some \(\delta>0\).  Exactly one of the
following alternatives holds.
\begin{enumerate}[label=\textup{(\roman*)}]
\item For every \(x_0\in\Sigma(T)\),
\[
  \limsup_{r\downarrow0}
  \bigl(C((x_0,T),r)+D((x_0,T),r)\bigr)=0.
\]
\item There exist \(x_0\in\Sigma(T)\), \(\eta>0\), and \(r_n\downarrow0\) such
that
\[
  C((x_0,T),r_n)+D((x_0,T),r_n)\ge\eta
  \qquad\text{for all }n.
\]
\end{enumerate}
\end{proposition}

\begin{proof}
The first alternative is a universal assertion over the singular set, and the
second is its failure expressed by a sequence.  If the first alternative fails,
then for some \(x_0\in\Sigma(T)\) the corresponding limsup is positive;
choosing a positive number below this limsup and a sequence along which the
lower bound holds gives the second alternative.  Conversely, the second
alternative contradicts the first.  Hence exactly one alternative holds.
\end{proof}

\begin{theorem}[Finite-time singularity yields positive local scale-invariant concentration]
\label{p0:thm:singular-positive}
Let \((u,p)\) be a suitable weak solution on
\(\mathbb R^3\times(T-\delta,T)\).  If \(\Sigma(T)\ne\emptyset\), then there
are \(x_0\in\Sigma(T)\), \(\eta>0\), and \(r_n\downarrow0\) such that
\[
  C((x_0,T),r_n)+D((x_0,T),r_n)\ge\eta
  \qquad\text{for all }n.
\]
\end{theorem}

\begin{proof}
If the first alternative in \cref{p0:prop:local-alternative} held, then
\cref{p0:thm:no-concentration} would imply \(\Sigma(T)=\emptyset\), contrary to
the hypothesis.  Therefore the second alternative in
\cref{p0:prop:local-alternative} holds.
\end{proof}

\begin{corollary}[Singular points carry positive velocity concentration]
\label{p0:cor:velocity-concentration}
Let \((u,p)\) be a suitable weak solution on
\(\mathbb R^3\times(T-\delta,T)\).  If \(x_0\in\Sigma(T)\), then there are
\(\eta_v>0\) and \(r_n\downarrow0\) such that
\[
  C((x_0,T),r_n)\ge\eta_v
  \qquad\text{for all }n.
\]
\end{corollary}

\begin{proof}
If the conclusion failed, then
\[
  \limsup_{r\downarrow0} C((x_0,T),r)=0.
\]
Choose \(r>0\) so small that
\[
  C((x_0,T),r)<\varepsilon_v,
\]
where \(\varepsilon_v\) is the constant in
\cref{p0:thm:velocity-ckn}.  The velocity \(\varepsilon\)-regularity theorem gives
boundedness of \(u\) in \(Q_{r/2}(x_0,T)\), contradicting
\(x_0\in\Sigma(T)\).  Hence the limsup is positive; taking any
\(\eta_v\) below it and passing to a sequence with
\(C((x_0,T),r_n)\ge\eta_v\) gives the result.
\end{proof}

\section{Automatic No-Escape for Finite-Energy Type I Singularities}
\label{p0:sec:auto-no-escape}

The preceding concentration result gives the positive velocity concentration
sequence used in the Type I blow-up analysis.  In the Seregin extraction
theorem \Cref{p1:thm:extraction}, the
remaining compactness hypothesis is a no-escape condition excluding the possibility
that all of the selected scale-invariant velocity mass collapses into the
final rescaled time slice.  We prove here that this no-escape condition is
automatic for finite-energy Type I singularities.

For this section we use, in addition to \(C\) and \(D\), the standard local
kinetic-energy and dissipation quantities
\[
  A(z_0,r)
  =r^{-1}\operatorname*{ess\,sup}_{t_0-r^2<t<t_0}
    \int_{B_r(x_0)} |u(x,t)|^2\,dx,
\]
and
\[
  E(z_0,r)
  =r^{-1}\int_{t_0-r^2}^{t_0}\int_{B_r(x_0)}
    |\nabla u(x,t)|^2\,dx\,dt .
\]
These quantities are invariant under the Navier--Stokes scaling.

\begin{definition}[Finite-energy Type I terminal singularity]
\label{p0:def:finite-energy-typeI-paper0}
Let \((u,p)\) be a suitable weak solution on \(\mathbb R^3\times(0,T)\).
We say that \(z_*=(x_*,T)\) is a finite-energy Type I terminal singularity if
\(x_*\in\Sigma(T)\),
\[
  u\in L^\infty(0,T;L^2(\mathbb R^3))
  \cap L^2(0,T;\dot H^1(\mathbb R^3)),
\]
and there are \(T_0<T\) and \(M<\infty\) such that
\begin{equation}\label{p0:eq:paper0-global-typeI}
  \|u(t)\|_{L^\infty(\mathbb R^3)}
  \le \frac{M}{\sqrt{T-t}},
  \qquad T_0<t<T .
\end{equation}
\end{definition}

The key point is a uniformly local energy propagation estimate.  It is stated
at unit scale, because the application is obtained by parabolic rescaling from
\((x_*,T)\).  The estimate is a local differential inequality: the unknown is the
uniformly-local kinetic energy, and the Type I envelope turns the nonlinear
pressure contribution into an integrable linear coefficient.  No bound below
depends on the size of the global \(L^2\) energy or the global dissipation;
finite energy is used only to place the pressure in the standard whole-space
Leray gauge and to make fixed-scale terminal \(L^3\) absolute continuity
immediate.

We use the standard Leray pressure gauge throughout this section.  At almost
every terminal time on which the Type I bound holds,
\(u(t)\in L^2(\mathbb R^3)\cap L^\infty(\mathbb R^3)\), hence
\(u_i(t)u_j(t)\in L^{3/2}(\mathbb R^3)\).  The whole-space pressure
\[
  p^L(t)=\mathcal R_i\mathcal R_j(u_i(t)u_j(t))
\]
is then defined by the Calderon--Zygmund theory.  The distributional pressure
equation gives
\(-\Delta p=\partial_i\partial_j(u_i u_j)=-\Delta p^L\), and the whole-space
de Rham argument applied to the momentum equation shows that \(p-p^L\) is a
function of time.  Thus replacing \(p\) by \(p^L\) is only a pressure-gauge
change and does not alter suitability.  This gauge choice uses the
finite-energy setting to define the representative; the uniform constants
below, in particular \(B_{\mathrm{uloc}}\) and \(B_A\), do not contain
\(\|u\|_{L^\infty(0,T;L^2(\mathbb R^3))}\) or
\(\int_0^T\|\nabla u(t)\|_{L^2(\mathbb R^3)}^2\,dt\).

\begin{lemma}[Uniformly-local energy propagation under a Type I envelope]
\label{p0:lem:uloc-typeI-propagation}
Let \((v,\pi)\) be a suitable weak solution on
\(\mathbb R^3\times(-4,0)\) with \(v(s)\in L^2(\mathbb R^3)\) for a.e.
\(s\), and take \(\pi\) to be the whole-space Leray pressure representative,
modulo functions of time.  Assume that
\begin{equation}\label{p0:eq:unit-typeI-envelope}
  |v(x,s)|\le M(-s)^{-1/2},
  \qquad \text{for a.e. }(x,s)\in\mathbb R^3\times(-4,0),
\end{equation}
and that for some admissible initial time \(s_0\in(-4,-3)\),
\begin{equation}\label{p0:eq:unit-uloc-initial}
  \alpha_0
  :=\sup_{a\in\mathbb R^3}\int_{B_2(a)} |v(x,s_0)|^2\,dx
  \le K .
\end{equation}
Then there is a constant \(B_{\mathrm{uloc}}=B_{\mathrm{uloc}}(M,K)<\infty\)
such that
\begin{equation}\label{p0:eq:unit-uloc-bound}
  \operatorname*{ess\,sup}_{-1<s<0}\sup_{a\in\mathbb R^3}
  \int_{B_1(a)} |v(x,s)|^2\,dx
  +
  \sup_{a\in\mathbb R^3}\int_{-1}^{0}\int_{B_1(a)}
  |\nabla v|^2\,dx\,ds
  \le B_{\mathrm{uloc}} .
\end{equation}
\end{lemma}

\begin{proof}
We give the proof because the estimate is the mechanism that closes the
terminal no-escape gap.  Let
\[
  \alpha(s)=\sup_{a\in\mathbb R^3}\int_{B_1(a)} |v(x,s)|^2\,dx .
\]
By \eqref{p0:eq:unit-typeI-envelope},
\begin{equation}\label{p0:eq:alpha-typeI-trivial}
  \alpha(s)^{1/2}\le C M(-s)^{-1/2},
  \qquad -4<s<0 .
\end{equation}
Fix \(a\in\mathbb R^3\) and choose
\(\phi_a\in C_c^\infty(B_2(a))\) with
\(0\le\phi_a\le1\), \(\phi_a\equiv1\) on \(B_1(a)\), and
\(|\nabla\phi_a|+|\Delta\phi_a|\le C\), uniformly in \(a\).  Applying the
local energy inequality on \((s_0,t)\), using the standard monotone
approximation of temporal cutoffs and with the pressure gauge chosen below,
gives for a.e. \(t\in(s_0,0)\)
\begin{align}
 &\int |v(t)|^2\phi_a\,dx
   +2\int_{s_0}^{t}\int |\nabla v|^2\phi_a\,dx\,ds       \notag \\
 &\quad\le
   \int |v(s_0)|^2\phi_a\,dx
   +C\int_{s_0}^{t}\int_{B_2(a)} |v|^2\,dx\,ds
   +C\int_{s_0}^{t}\int_{B_2(a)} |v|^3\,dx\,ds          \notag \\
 &\qquad
   +C\int_{s_0}^{t}\int_{B_2(a)}
       |\pi-c_a(s)|\,|v|\,dx\,ds .
\label{p0:eq:lei-unit-estimate}
\end{align}
Here \(c_a(s)\) is an arbitrary function of time; adding or subtracting such a
function does not change the local energy inequality, since
\(\nabla\cdot v=0\) and \(\phi_a\) is compactly supported.

The velocity terms are controlled by the state.  Since \(B_2(a)\) is covered by
finitely many unit balls,
\[
  \int_{B_2(a)} |v|^2\,dx\le C\alpha(s),
\]
and by \eqref{p0:eq:unit-typeI-envelope},
\begin{equation}\label{p0:eq:cubic-unit-bound}
  \int_{B_2(a)} |v|^3\,dx
  \le M(-s)^{-1/2}\int_{B_2(a)} |v|^2\,dx
  \le C M(-s)^{-1/2}\alpha(s).
\end{equation}
It remains to estimate the pressure term uniformly in \(a\).

Use the whole-space pressure representative
\(\pi=\mathcal R_i\mathcal R_j(v_i v_j)\), modulo functions of time.  Let
\(\chi_a\in C_c^\infty(B_8(a))\), \(\chi_a\equiv1\) on \(B_4(a)\), and write
inside \(B_4(a)\)
\[
  \pi=q_a+h_a,
  \qquad
  q_a=\mathcal R_i\mathcal R_j(\chi_a v_i v_j).
\]
Then \(h_a\) is harmonic in \(B_4(a)\).  Calderon--Zygmund boundedness
\cite{Stein1970} and \eqref{p0:eq:unit-typeI-envelope} give, for a.e. \(s\),
\begin{align*}
  \|q_a(s)\|_{L^2(B_4(a))}
  &\le C\|\chi_a v(s)\otimes v(s)\|_{L^2(\mathbb R^3)}  \\
  &\le C M(-s)^{-1/2}\|v(s)\|_{L^2(B_8(a))}             \\
  &\le C M(-s)^{-1/2}\alpha(s)^{1/2},
\end{align*}
where the last step again uses a finite covering by unit balls.  Hence
\begin{equation}\label{p0:eq:local-pressure-unit}
  \int_{B_2(a)} |q_a|\,|v|\,dx
  \le C M(-s)^{-1/2}\alpha(s).
\end{equation}

For the harmonic part, set \(c_a(s)=h_a(a,s)\).  The kernel estimate
\[
  |K_{ij}(x-z)-K_{ij}(a-z)|
  \le C |x-a|\,|z-a|^{-4},
  \qquad x\in B_2(a),\quad |z-a|>4,
\]
where \(K_{ij}\) is the kernel of \(\mathcal R_i\mathcal R_j\), gives
\[
  |h_a(x,s)-c_a(s)|
  \le
  C\sum_{m=2}^{\infty}2^{-4m}
    \int_{2^m<|z-a|<2^{m+1}} |v(z,s)|^2\,dz .
\]
The annulus \(\{2^m<|z-a|<2^{m+1}\}\) is covered by at most \(C2^{3m}\)
unit balls.  Therefore
\begin{equation}\label{p0:eq:harmonic-pressure-unit}
  \|h_a(s)-c_a(s)\|_{L^\infty(B_2(a))}
  \le C\sum_{m=2}^{\infty}2^{-m}\alpha(s)
  \le C\alpha(s).
\end{equation}
Consequently,
\begin{align}
  \int_{B_2(a)} |h_a-c_a|\,|v|\,dx
  &\le C\alpha(s)\int_{B_2(a)} |v|\,dx                    \notag\\
  &\le C\alpha(s)^{3/2}                                    \notag\\
  &\le C M(-s)^{-1/2}\alpha(s),
\label{p0:eq:harmonic-pressure-flux-unit}
\end{align}
where the last inequality is exactly the Type I absorbing estimate
\eqref{p0:eq:alpha-typeI-trivial}.
Combining \eqref{p0:eq:local-pressure-unit} and
\eqref{p0:eq:harmonic-pressure-flux-unit},
\begin{equation}\label{p0:eq:pressure-flux-unit}
  \int_{B_2(a)} |\pi-c_a(s)|\,|v|\,dx
  \le C M(-s)^{-1/2}\alpha(s).
\end{equation}

Insert \eqref{p0:eq:cubic-unit-bound} and \eqref{p0:eq:pressure-flux-unit} into
\eqref{p0:eq:lei-unit-estimate}, then take the supremum over \(a\).  Since
\((-s)^{-1/2}\in L^1(-4,0)\), Gronwall's inequality yields
\[
  \operatorname*{ess\,sup}_{s_0<s<0}\alpha(s)
  \le C K
     \exp\left(C\int_{s_0}^{0}\bigl(1+M(-s)^{-1/2}\bigr)\,ds\right)
  \le B_1(M,K).
\]
Returning to \eqref{p0:eq:lei-unit-estimate}, using this bound, and then letting
\(t\uparrow0\) through admissible times gives the corresponding
uniformly-local dissipation estimate on \((-1,0)\).  This proves
\eqref{p0:eq:unit-uloc-bound}.
\end{proof}

\begin{remark}[Local energy estimate as an absorbing bound]
\label{p0:rem:state-space-stratification}
The preceding proof can be read as a three-term estimate for the local quantity
\(\alpha(s)\).  The diffusion and transport terms are linear in
\(\alpha\), with coefficient \(1+M(-s)^{-1/2}\).  The local pressure term is
also linear after Calderon--Zygmund and the Type I envelope.  The remaining
nonlinear term is the harmonic pressure flux, which is
\(O(\alpha^{3/2})\); the pointwise Type I envelope implies
\(\alpha^{1/2}\le C M(-s)^{-1/2}\), reducing it to the same integrable linear
coefficient.  Thus large values of \(\alpha\) cannot persist: the resulting
differential inequality traps \(\alpha\) in a bounded interval depending
only on \(M\) and the initial uniformly local energy level.
\end{remark}

\begin{proposition}[Automatic terminal kinetic tightness]
\label{p0:prop:auto-terminal-A}
Let \((u,p,z_*)\), \(z_*=(x_*,T)\), be a finite-energy Type I terminal
singularity in the sense of \cref{p0:def:finite-energy-typeI-paper0}.  Then there
exist \(r_0>0\) and \(B_A<\infty\), with \(B_A\) depending only on the Type I
constant \(M\), such that
\begin{equation}\label{p0:eq:auto-A-bound}
  A(z_*,r)\le B_A,
  \qquad 0<r<r_0 .
\end{equation}
Moreover,
\begin{equation}\label{p0:eq:auto-E-bound}
  r^{-1}\int_{T-r^2}^{T}\int_{B_r(x_*)}|\nabla u|^2\,dx\,dt
  \le B_A,
  \qquad 0<r<r_0 .
\end{equation}
\end{proposition}

\begin{proof}
Choose \(r_0>0\) so small that \(T-4r^2>T_0\) for \(0<r<r_0\).  For such
\(r\), use the whole-space Leray pressure gauge fixed above and define the
rescaled solution
\[
  v^{(r)}(y,s)=r u(x_*+ry,T+r^2s),
  \qquad
  \pi^{(r)}(y,s)=r^2 p(x_*+ry,T+r^2s),
\]
on \(\mathbb R^3\times(-4,0)\).  The Type I bound
\eqref{p0:eq:paper0-global-typeI} gives
\[
  |v^{(r)}(y,s)|\le M(-s)^{-1/2},
  \qquad -4<s<0 .
\]
For any admissible \(s_0\in(-4,-3)\) at which the Type I bound holds, and any
\(a\in\mathbb R^3\),
\begin{align*}
  \int_{B_2(a)} |v^{(r)}(y,s_0)|^2\,dy
  &= r^{-1}\int_{B_{2r}(x_*+ra)}
       |u(x,T+r^2s_0)|^2\,dx                                    \\
  &\le r^{-1}|B_{2r}|\frac{M^2}{r^2|s_0|}
   \le C M^2 .
\end{align*}
Such times form a full-measure subset of \((-4,-3)\).  The preceding estimate
is uniform in this choice, so \cref{p0:lem:uloc-typeI-propagation} applies with
\(K=CM^2\), uniformly in \(r\).  Hence
\[
  \operatorname*{ess\,sup}_{-1<s<0}\int_{B_1}|v^{(r)}(y,s)|^2\,dy
  +\int_{-1}^{0}\int_{B_1}|\nabla v^{(r)}|^2\,dy\,ds
  \le B_A .
\]
Undoing the parabolic change of variables gives exactly
\eqref{p0:eq:auto-A-bound} and \eqref{p0:eq:auto-E-bound}.  The constant \(B_A\)
depends only on \(M\); no global energy norm of \(u\) has been used in the
estimate.
\end{proof}

\begin{theorem}[Automatic velocity no-escape]
\label{p0:thm:auto-velocity-no-escape}
Let \((u,p,z_*)\), \(z_*=(x_*,T)\), be a finite-energy Type I terminal
singularity.  Let \(\lambda_k\downarrow0\) be any sequence of scales, and set
\[
  u_k(x,t)=\lambda_k u(x_*+\lambda_k x,T+\lambda_k^2 t).
\]
Then
\begin{equation}\label{p0:eq:auto-no-escape}
  \lim_{\sigma\downarrow0}
  \sup_k\int_{-\sigma}^{0}\int_{B_1}|u_k(x,t)|^3\,dx\,dt=0 .
\end{equation}
\end{theorem}

\begin{proof}
By \cref{p0:prop:auto-terminal-A}, after discarding at most finitely many indices,
\[
  \operatorname*{ess\,sup}_{-1<t<0}\int_{B_1}|u_k(x,t)|^2\,dx\le B_A .
\]
The Type I bound gives, for the same tail of the sequence,
\[
  \|u_k(t)\|_{L^\infty(\mathbb R^3)}
  \le M(-t)^{-1/2},
  \qquad -1<t<0 .
\]
Therefore, for \(0<\sigma<1\),
\[
\begin{aligned}
  \int_{-\sigma}^{0}\int_{B_1}|u_k|^3\,dx\,dt
  &\le
  \int_{-\sigma}^{0}\|u_k(t)\|_{L^\infty(B_1)}
       \int_{B_1}|u_k(x,t)|^2\,dx\,dt  \\
  &\le M B_A\int_{-\sigma}^{0}(-t)^{-1/2}\,dt
   =2M B_A\sigma^{1/2} .
\end{aligned}
\]
This bound is uniform over the tail.  For each of the finitely many discarded
indices \(k\), choose \(\sigma_k>0\) so that
\(T+\lambda_k^2t>T_0\) for \(-\sigma_k<t<0\).  On this terminal subinterval
the Type I envelope and the finite-energy assumption imply
\[
  \int_{-\sigma_k}^{0}\int_{B_1}|u_k|^3\,dx\,dt
  \le
  \int_{-\sigma_k}^{0}\|u_k(t)\|_{L^\infty(B_1)}
       \int_{B_1}|u_k(x,t)|^2\,dx\,dt
  <\infty .
\]
Indeed, on \((-\sigma_k,0)\),
\[
  \|u_k(t)\|_{L^\infty(B_1)}\le M(-t)^{-1/2},
  \qquad
  \int_{B_1}|u_k(x,t)|^2\,dx
  \le \lambda_k^{-1}\|u(T+\lambda_k^2t)\|_{L^2(\mathbb R^3)}^2,
\]
and the right-hand side is finite for this fixed \(k\) by finite energy.
Thus each discarded index has absolute continuity of its \(L^3\)-integral at
the terminal slice.  Taking the maximum over this finite set gives the same
limit for the full supremum over \(k\).  This proves
\eqref{p0:eq:auto-no-escape}.
\end{proof}

\begin{corollary}[Automatic admissibility of Type I concentration sequences]
\label{p0:cor:auto-paperI-admissibility}
Let \((u,p,z_*)\), \(z_*=(x_*,T)\), be a finite-energy Type I terminal
singularity.  Then there exist \(\eta_v>0\) and \(\lambda_k\downarrow0\) such
that
\[
  C(z_*,\lambda_k)\ge\eta_v
  \qquad\text{for every }k,
\]
and the corresponding rescaled velocities satisfy the no-escape condition
\eqref{p0:eq:auto-no-escape}.  Hence the concentration and no-escape
admissibility condition used in the Type I blow-up-limit theorem
\Cref{p1:thm:extraction} is automatic for finite-energy Type I singularities.
\end{corollary}

\begin{proof}
Since \(x_*\in\Sigma(T)\), \cref{p0:cor:velocity-concentration} gives
\(\eta_v>0\) and \(\lambda_k\downarrow0\) with
\(C(z_*,\lambda_k)\ge\eta_v\) for all \(k\).  The no-escape condition for this
same sequence follows from \cref{p0:thm:auto-velocity-no-escape}.
\end{proof}

\section{The Type I/Type II Blow-Up Rate Dichotomy}

Let \(x_0\in\Sigma(T)\), \(\eta>0\), and \(r_n\downarrow0\) be given by
\cref{p0:thm:singular-positive}.  Define
\[
  u_n(y,s)=r_nu(x_0+r_ny,T+r_n^2s),
  \qquad
  p_n(y,s)=r_n^2p(x_0+r_ny,T+r_n^2s).
\]
Then \((u_n,p_n)\) are suitable weak solutions on expanding backward
cylinders, and by \cref{p0:prop:critical-scaling} they retain a positive lower
bound for the mean-subtracted local scale-invariant quantity on a fixed
bounded cylinder.

\begin{definition}[Type I concentration sequence]
\label{p0:def:typeI-concentration}
Let \((u,p)\) be a suitable weak solution on
\(\mathbb R^3\times(T-\delta,T)\), and let \(r_n\downarrow0\) be a positive
local concentration sequence at \((x_0,T)\).  The sequence is called Type I
if the concentration point \((x_0,T)\) satisfies a local Type I
scale-invariant bound:
there are \(\rho>0\) and \(M<\infty\), with \(\rho^2<\delta\), such that
\[
  \operatorname*{ess\,sup}_{T-\rho^2<t<T}
  \sqrt{T-t}\,\|u(t)\|_{L^\infty(B_\rho(x_0))}\le M.
\]
This is the Type I alternative.
\end{definition}

\begin{definition}[Local Type II alternative]
\label{p0:def:typeII-alternative}
Let \((u,p)\) be a suitable weak solution on
\(\mathbb R^3\times(T-\delta,T)\), and let \(r_n\downarrow0\) be a positive
local concentration sequence at \((x_0,T)\).  The sequence belongs to the local
Type II alternative if the concentration point \((x_0,T)\) satisfies no
local Type I scale-invariant bound.  Equivalently, for every \(\rho>0\) with
\(\rho^2<\delta\),
\[
  \operatorname*{ess\,sup}_{T-\rho^2<t<T}
  \sqrt{T-t}\,\|u(t)\|_{L^\infty(B_\rho(x_0))}=\infty.
\]
\end{definition}

\begin{theorem}[Local blow-up rate dichotomy]\label{p0:thm:typeI-typeII-dichotomy}
Let \((u,p)\) be a suitable weak solution on
\(\mathbb R^3\times(T-\delta,T)\), and suppose \(\Sigma(T)\ne\emptyset\).
Then there exist \(x_0\in\Sigma(T)\),
\(\eta>0\), and \(r_n\downarrow0\) such that
\[
  C((x_0,T),r_n)+D((x_0,T),r_n)\ge\eta
  \qquad\text{for all }n.
\]
The sequence \(r_n\) is a positive local concentration sequence.  The
corresponding concentration regime satisfies either the local Type I bound or
the local Type II alternative.  Consequently, once the vanishing case has been
excluded by \(\varepsilon\)-regularity, every finite-time singular case belongs
to exactly one of the two blow-up rate alternatives.
\end{theorem}

\begin{proof}
If the local scale-invariant quantities vanished at every point of
\(\Sigma(T)\), then \cref{p0:thm:no-concentration} would give
\(\Sigma(T)=\emptyset\), contrary to the hypothesis.  Therefore
\cref{p0:thm:singular-positive} gives
\(x_0\in\Sigma(T)\), \(\eta>0\), and \(r_n\downarrow0\) with the stated lower
bound.  The rescaled sequence has a positive local scale-invariant lower bound
by \cref{p0:prop:critical-scaling}.  By definition, either the Type I bound in
\cref{p0:def:typeI-concentration} holds, or the local Type II alternative in
\cref{p0:def:typeII-alternative} holds.
\end{proof}

\begin{remark}
This theorem is local.  Its conclusion is the existence of a positive
mean-subtracted scale-invariant concentration sequence and the dichotomy
between a local Type I bound and failure of every such bound at the
concentration point.
\end{remark}

\section{Local Pointwise Type I Reduction to Seregin Extraction}
\label{p0:sec:local-typeI-terminal-bridge-paper0}

We use the following consequence of the weak-Serrin endpoint argument.
The local pointwise Type I alternative in
\cref{p0:def:typeI-concentration} is strong enough, after rescaling at the singular
point, to give compactness on every compact negative-time cylinder.  The
endpoint weak-Serrin implication reads:
for suitable weak solutions the local pointwise Type I envelope belongs to
\(L^{2,\infty}_tL^\infty_x\), hence the Albritton--Barker weak Serrin local
Type I estimate gives the scale-invariant \(A\)-bound on smaller cylinders and
therefore terminal velocity no-escape.  Thus positive local Type I
concentration sequences are admissible for the Seregin extraction theorem \Cref{p1:thm:extraction}.

Throughout this section \(z_*=(x_*,T)\), and
\[
  Q_1(0,0)=B_1(0)\times(-1,0).
\]

\begin{definition}[Admissible local Type I concentration sequence]
\label{p0:def:paper0-admissible-local-bridge-sequence}
For a local pointwise Type I point \(z_*=(x_*,T)\), a sequence
\(r_n\downarrow0\) is an admissible local Type I concentration sequence if the rescaled
velocities
\[
  u_n(x,t)=r_nu(x_*+r_nx,T+r_n^2t)
\]
satisfy, for some \(\eta_v>0\),
\[
   \int_{Q_1(0,0)} |u_n|^3\,dx\,dt\ge\eta_v
   \qquad\text{for every }n,
\]
and
\[
   \lim_{\sigma\downarrow0}
   \sup_n
   \int_{-\sigma}^{0}\int_{B_1}|u_n|^3\,dx\,dt=0 .
\]
\end{definition}

\begin{proposition}[Local pointwise Type I rescalings have terminal compactness]
\label{p0:prop:paper0-local-typeI-terminal-compactness}
Let \((u,p)\) be a suitable weak solution on
\(\mathbb R^3\times(0,T)\) satisfying the finite-energy bounds
\[
  u\in L^\infty(0,T;L^2(\mathbb R^3))
  \cap L^2(0,T;\dot H^1(\mathbb R^3)).
\]
Let \(z_*=(x_*,T)\) be a singular point, and assume that \(z_*\) satisfies the
local pointwise Type I bound: there exist \(\rho>0\) and \(M<\infty\) such that
\begin{equation}\label{p0:eq:paper0-local-typeI-bridge-bound}
  \operatorname*{ess\,sup}_{T-\rho^2<t<T}
  \sqrt{T-t}\,\|u(t)\|_{L^\infty(B_\rho(x_*))}\le M .
\end{equation}
For any sequence \(r_n\downarrow0\), define
\begin{equation}\label{p0:eq:paper0-local-typeI-bridge-rescaling}
  u_n(x,t)=r_nu(x_*+r_nx,T+r_n^2t),
  \qquad
  p_n(x,t)=r_n^2p(x_*+r_nx,T+r_n^2t).
\end{equation}
Then the following hold.
\begin{enumerate}[label=\textup{(\roman*)}]
\item If
\[
  K=B_R\times[-S,-\sigma]\Subset\mathbb R^3\times(-\infty,0),
  \qquad R,S<\infty,\quad \sigma>0,
\]
then, for all sufficiently large \(n\),
\begin{equation}\label{p0:eq:paper0-local-typeI-bridge-Linfty}
  \|u_n\|_{L^\infty(K)}\le M\sigma^{-1/2}.
\end{equation}
\item On each such compact cylinder \(K\), after subtracting time-dependent
local pressure gauges \(a_{n,K}(t)\), the sequence satisfies uniform local
Seregin bounds
\begin{equation}\label{p0:eq:paper0-local-typeI-bridge-Seregin-bounds}
  \|u_n\|_{L^\infty_tL^2_x(K)}
  +\|\nabla u_n\|_{L^2(K)}
  +\|p_n-a_{n,K}(t)\|_{L^{3/2}(K)}
  \le C_K .
\end{equation}
Here \(C_K\) may depend on \(K,M,\rho\), and on the finite Leray energy level
of the original solution, but it does not depend on \(n\) and does not use a
terminal no-escape assumption.
\item After passing to a subsequence, \((u_n,p_n-a_{n,K})\) converges on every
compact \(K\Subset\mathbb R^3\times(-\infty,0)\) to an ancient suitable weak
solution \((U,Q)\), with
\[
  u_n\to U\quad\text{strongly in }L^q(K),\qquad 1\le q<\infty,
\]
\[
  p_n-a_{n,K}(t)\rightharpoonup Q
  \quad\text{weakly in }L^{3/2}(K),
\]
and
\[
  |U(x,t)|\le M(-t)^{-1/2},\qquad t<0.
\]
\end{enumerate}
\end{proposition}

\begin{proof}
Fix \(K=B_R\times[-S,-\sigma]\Subset\mathbb R^3\times(-\infty,0)\).  For
large \(n\),
\[
  r_nR<\rho/2,\qquad r_n^2S<\rho^2/2.
\]
Thus \(x_*+r_nx\in B_\rho(x_*)\) and
\(T+r_n^2t\in(T-\rho^2,T)\) for a.e. \((x,t)\in K\).  The local pointwise Type I bound
\eqref{p0:eq:paper0-local-typeI-bridge-bound} gives
\[
  |u_n(x,t)|
  =r_n|u(x_*+r_nx,T+r_n^2t)|
  \le r_nM(T-(T+r_n^2t))^{-1/2}
  =M(-t)^{-1/2}
  \le M\sigma^{-1/2}.
\]
This proves \eqref{p0:eq:paper0-local-typeI-bridge-Linfty}.  It also gives
uniform local \(L^\infty_tL^2_x\) and \(L^3_{x,t}\) bounds on \(K\).

It remains to explain the pressure and dissipation bounds.  Use the
whole-space Leray pressure representative, which is legitimate in the
finite-energy class up to addition of functions of time.  Choose
\(\chi_R\in C_c^\infty(B_{8R})\), \(\chi_R\equiv1\) on \(B_{4R}\), and write
inside \(B_{4R}\)
\[
  p_n=q_n+h_n,\qquad
  q_n=\mathcal R_i\mathcal R_j(\chi_Ru_{n,i}u_{n,j}).
\]
Calderon--Zygmund boundedness and the local \(L^3\)-bound for \(u_n\) give a
uniform \(L^{3/2}\)-bound for \(q_n\).  For \(h_n\), take the gauge
\(a_{n,K}(t)=h_n(0,t)\).  The harmonic pressure kernel estimate gives
\[
  |h_n(x,t)-a_{n,K}(t)|
  \le
  C_R\sum_{m\ge2}2^{-4m}
  \int_{2^mR<|z|<2^{m+1}R}|u_n(z,t)|^2\,dz,
  \qquad x\in B_{2R}.
\]
For annuli whose physical image remains inside \(B_\rho(x_*)\), the local
Type I bound controls the \(L^2\)-mass by
\[
  C M^2\sigma^{-1}(2^mR)^3,
\]
and the resulting series is geometric.  For the outer annuli, the finite
Leray energy gives
\[
  \int_{\mathbb R^3}|u_n(z,t)|^2\,dz
  =r_n^{-1}\int_{\mathbb R^3}|u(y,T+r_n^2t)|^2\,dy,
\]
while the first outer annulus has radius comparable to \(\rho/r_n\).  Hence
the tail of the kernel series is bounded by
\[
  C_{\rho,R}\,r_n^3\|u\|_{L^\infty(0,T;L^2)}^2,
\]
which is uniform in \(n\).  This proves the pressure-gauge part of
\eqref{p0:eq:paper0-local-typeI-bridge-Seregin-bounds}.

Apply the local energy inequality for \((u_n,p_n)\) with a cutoff equal to one
on \(K\) and supported in a slightly larger compact cylinder.  The local
kinetic, cubic velocity, and pressure-flux terms are bounded by the preceding
estimates, so \(\|\nabla u_n\|_{L^2(K)}\le C_K\).  The Navier--Stokes
equation then bounds \(\partial_tu_n\) locally in a negative Sobolev space.
Aubin--Lions compactness and the pressure bound give the asserted Seregin
subsequence convergence.  Passing to the limit in the equations and local
energy inequality gives an ancient suitable weak solution, and the pointwise
Type I bound passes to the limit.
\end{proof}

\begin{corollary}[Local pointwise Type I concentration sequences are admissible]
\label{p0:cor:paper0-local-typeI-bridge-admissible}
Under the hypotheses of
\cref{p0:prop:paper0-local-typeI-terminal-compactness}, let \(r_n\downarrow0\) be
any sequence such that the rescaled velocities satisfy
\[
  \int_{Q_1(0,0)} |u_n|^3\,dx\,dt\ge\eta_v>0
  \qquad\text{for all }n .
\]
Then \(r_n\) is an admissible local Type I concentration sequence in the sense of
\cref{p0:def:paper0-admissible-local-bridge-sequence}.
\end{corollary}

\begin{proof}
The displayed positive concentration lower bound is the first requirement in
\cref{p0:def:paper0-admissible-local-bridge-sequence}.  The Albritton--Barker
local weak-Serrin Type I estimate implies
that the local pointwise Type I envelope implies terminal velocity no-escape
for every rescaled sequence
\cite{AlbrittonBarker2019}.  This is exactly the second requirement in
\cref{p0:def:paper0-admissible-local-bridge-sequence}.
\end{proof}

\begin{lemma}[Local Type I concentration enters the Seregin extraction]
\label{p0:lem:paper0-local-typeI-into-terminal-dichotomy}
Under the hypotheses of
\cref{p0:prop:paper0-local-typeI-terminal-compactness}, let \(r_n\downarrow0\) be
a velocity concentration sequence satisfying
\[
  \int_{Q_1(0,0)} |u_n|^3\,dx\,dt\ge\eta_v>0
  \qquad\text{for all }n .
\]
Then \(r_n\) is admissible by
\cref{p0:cor:paper0-local-typeI-bridge-admissible}.  There are
\(\sigma_0\in(0,1)\) and \(\eta_1>0\) such that
\begin{equation}\label{p0:eq:paper0-local-typeI-bridge-away-mass}
  \int_{-1}^{-\sigma_0}\int_{B_1}|u_n|^3\,dx\,dt\ge\eta_1
  \qquad\text{for all }n.
\end{equation}
Consequently, the rescaled sequence retains positive local \(L^3\) mass on a
compact negative-time cylinder.  In particular the associated Seregin
extraction has no loss of the local \(L^3\) mass needed for
nontriviality.
\end{lemma}

\begin{proof}
The no-escape condition in
\cref{p0:def:paper0-admissible-local-bridge-sequence} gives \(\sigma_0>0\) such
that the terminal layer \(B_1\times(-\sigma_0,0)\) carries at most half of the
fixed \(Q_1\) velocity mass.  This proves
\eqref{p0:eq:paper0-local-typeI-bridge-away-mass}.  The compactness part follows
from \cref{p0:prop:paper0-local-typeI-terminal-compactness} on the compact
negative-time cylinder \(B_1\times[-1,-\sigma_0]\).  Strong local \(L^3\)
convergence on this cylinder passes the positive mass to the ancient limit.
This is the only point at which terminal no-escape is used.

Thus the local Type I concentration sequence is admissible for the Seregin
extraction and retains positive compact \(L^3\)-mass away from the terminal
slice.
\end{proof}

\begin{theorem}[Conditional local Type I exclusion]
\label{p0:thm:paper0-local-typeI-entry-pipeline}
Assume the residual-class hypothesis \Cref{p1:hyp:no-remainder}.
Let \((u,p)\) be a finite-energy suitable weak solution on
\(\mathbb R^3\times(0,T)\).  Then no singular point \(z_*=(x_*,T)\) satisfying
the local pointwise Type I bound \eqref{p0:eq:paper0-local-typeI-bridge-bound}
exists.  More precisely, any such point would generate a bounded centered Type I
Seregin limit with positive local concentration, and that limit is excluded
by the resolved Liouville alternatives together with
\Cref{p1:hyp:no-remainder}.
\end{theorem}

\begin{proof}
Assume \Cref{p1:hyp:no-remainder}.
Suppose, toward a contradiction, that \(z_*=(x_*,T)\) is a singular point
satisfying \eqref{p0:eq:paper0-local-typeI-bridge-bound}.
Choose the positive \(Q_1\) velocity concentration sequence supplied by the
local concentration theorem.  By
\cref{p0:cor:paper0-local-typeI-bridge-admissible}, this sequence satisfies
terminal no-escape and is admissible.
By \cref{p0:lem:paper0-local-typeI-into-terminal-dichotomy}, the local Type I
sequence has local compactness and positive compact \(L^3\)-mass away from the
terminal slice.
The compactness statement gives an ancient Type I limit satisfying
\(|U(x,t)|\le M(-t)^{-1/2}\); after the centered change of variables this is a
bounded centered Type I Seregin limit.  By the class exhaustion in
\cref{p1:prop:exhaustion}, this limit lies either in one of the resolved
Liouville classes or in the remaining class \(\calR(\calS)\) for the
collection \(\calS\) generated by the extraction.  The resolved classes are
excluded by the small-amplitude, stationary, uniformly tight, and
structure/decay Liouville results proved or cited in this paper.  The remaining
class is empty by the residual-class hypothesis.  This contradicts the
local \(L^3\) nonvanishing, so the assumed local Type I singular point
cannot exist under \Cref{p1:hyp:no-remainder}.
\end{proof}

\begin{corollary}[Conditional exclusion of the local Type I alternative]
\label{p0:cor:paper0-local-typeII-after-bridge}
Assume the residual-class hypothesis \Cref{p1:hyp:no-remainder}.  Then the
local Type I alternative is excluded for every finite-energy singular point.
Hence every finite-energy singular point belongs to the local Type II
alternative in the dichotomy above.
\end{corollary}

\begin{proof}
This is \cref{p0:thm:paper0-local-typeI-entry-pipeline} restated in the terminology
of the local Type I/local Type II dichotomy.
\end{proof}



\clearpage
\section*{Part II. Type I Blow-up Limits and the Residual Hypothesis}

\subsection*{Overview}

We reduce an admissible finite-energy Type I singularity for the
three-dimensional Navier--Stokes equations to a normalized nonzero bounded
ancient solution of the centered self-similar Navier--Stokes equation.  The
construction keeps track of pressure gauges, local compactness, stability of
the local energy inequality, and the velocity concentration supplied by local
Caffarelli--Kohn--Nirenberg theory.  It yields a smooth bounded ancient solution
for the centered equation on \(\R^3\times\R\), together with a quantitative
local \(L^3\) nonvanishing condition.

We then organize the possible normalized Seregin ancient limits by standard
analytic hypotheses: small amplitude, stationary \(L^3\), uniform
\(L^3\)-tightness, and structure or decay assumptions covered by cited
Liouville theorems.  The small-amplitude case is excluded by a bounded ancient
perturbative Liouville theorem proved below, and the stationary \(L^3\) case by
the theorem of Ne\v{c}as--R\r{u}\v{z}i\v{c}ka--\v{S}ver\'ak.  The
structure/decay class contains only hypotheses for which the cited Liouville
theorems force \(V\equiv0\).  The uniformly \(L^3\)-tight class is closed in
\Cref{p2:thm:tight-liouville}.  The remaining class is recorded as the
residual-class hypothesis \Cref{p1:hyp:no-remainder}.  Consequently the
Type I conclusion obtained from this extraction is conditional on that residual
hypothesis within the stated admissible finite-energy class.

\section{Introduction and main results}\label{p1:sec:introduction}

\subsection{Type I blow-up limits}

Let \(u:\R^3\times(0,T)\to\R^3\) solve the incompressible Navier--Stokes
equations
\begin{equation}\label{p1:eq:NS}
   \partial_tu+(u\cdot\nabla)u+\nabla p=\Delta u,\qquad
   \divergence u=0.
\end{equation}
The local regularity theory of Caffarelli--Kohn--Nirenberg \cite{CKN1982}
detects singular points through scale-invariant concentration of velocity and
pressure.  A terminal point \(z_*=(x_*,T)\) is of Type I if
\begin{equation}\label{p1:eq:type-I-bound}
   \limsup_{t\uparrow T}\sqrt{T-t}\,
   \|u(t)\|_{L^\infty(\R^3)}<\infty .
\end{equation}
The blow-up extraction uses a selected velocity concentration sequence from
the local concentration theorem and assumes only that this velocity
concentration does not escape into the terminal slice after rescaling.  The
resulting admissibility condition is defined in
\cref{p1:def:admissible-type-I}.  It replaces a stronger all-scale compactness
assumption on scale-invariant energy, pressure, and dissipation quantities by
the non-escape condition needed for nontriviality of the ancient limit.

\subsection{Normalized Seregin limits}

For scales \(\lambda\downarrow0\), the parabolic rescaling about \(z_*\) is
\[
   u^{(\lambda)}(x,t)
   =
   \lambda u(x_*+\lambda x,T+\lambda^2t),\qquad
   p^{(\lambda)}(x,t)
   =
   \lambda^2p(x_*+\lambda x,T+\lambda^2t).
\]
The rescaled time intervals exhaust \((-\infty,0)\).  The Type I bound
becomes a uniform bound
\[
   \|u^{(\lambda)}(t)\|_{L^\infty(\R^3)}
   \le M(-t)^{-1/2}
\]
on compact subintervals of \((-\infty,0)\), and compactness produces an
unrenormalized ancient solution \((U,Q)\).  The associated centered variables
are
\[
   y=\frac{x}{\sqrt{-t}},\qquad
   \tau=-\log(-t),\qquad
   V(y,\tau)=\sqrt{-t}\,U(y\sqrt{-t},t).
\]
They transform the ancient Navier--Stokes solution into the centered equation
\begin{equation}\label{p1:eq:centered}
   \partial_\tau V+(V\cdot\nabla)V+\nabla\Pi
   =
   \Delta V-\frac12(V+y\cdot\nabla V),
   \qquad \divergence V=0.
\end{equation}
The selected velocity concentration at the singular point persists through
the extraction and gives a compact cylinder \(K\Subset\R^3\times\R\) and a
constant \(\eta_0>0\) such that
\begin{equation}\label{p1:eq:nonvanishing}
   \iint_K |V|^3\,dy\,d\tau\ge\eta_0 .
\end{equation}

\subsection{Liouville alternatives and the remaining case}

No unconditional Liouville theorem is known for all bounded ancient solutions
of \eqref{p1:eq:centered}.  We therefore work with a finite decomposition based
on explicit Liouville theorems.  A normalized Seregin limit is compared with the
following analytic classes:
\[
   \calC_{\mathrm{small}},\qquad
   \calC_{\mathrm{stat},L^3},\qquad
   \calC_{L^3\mathrm{-tight}},\qquad
   \calC_{\mathrm{known}},\qquad
   \calR(\calS).
\]
They denote, respectively, small amplitude, stationary \(L^3\), uniform
\(L^3\)-tightness, structure or decay covered by cited Liouville theorems, and
the complement of these four classes inside a fixed collection \(\calS\) of
normalized limits.  The last set represents no additional analytic mechanism;
it is the remaining case after the stated Liouville alternatives have been
removed.

\subsection{Main theorems}

\begin{theorem}[Normalized Type I blow-up limit]\label{p1:thm:extraction}
Let \((u,p,z_*)\) be an admissible finite-energy Type I singularity.  Then
there exist scales \(\lambda_k\downarrow0\), pressure gauges \(a_k(t)\), an
ancient suitable weak solution \((U,Q)\) on
\(\R^3\times(-\infty,0)\), and a centered solution \((V,\Pi)\) on
\(\R^3\times\R\), such that, after passing to a subsequence,
\[
   u_k\to U
   \quad\hbox{strongly in }L^q_{\loc}(\R^3\times(-\infty,0)),
   \qquad 1\le q<\infty,
\]
and
\[
   p_k-a_k(t)\rightharpoonup Q
   \quad\hbox{weakly in }L^{3/2}_{\loc}(\R^3\times(-\infty,0)).
\]
The centered velocity \(V\) is smooth, bounded, solves \eqref{p1:eq:centered}, and
satisfies the local \(L^3\) nonvanishing condition \eqref{p1:eq:nonvanishing}.
\end{theorem}

For each \(V\in\calS\), the class definitions give the exhaustive alternative
\[
   V\in
   \calC_{\mathrm{small}}
   \cup
   \calC_{\mathrm{stat},L^3}
   \cup
   \calC_{L^3\mathrm{-tight}}
   \cup
   \calC_{\mathrm{known}}
   \cup
   \calR(\calS).
\]
This is stated later as \cref{p1:prop:exhaustion}.  It is a direct consequence
of the definition of the remaining class, not an independent Liouville theorem.

\begin{theorem}[Conditional Type I exclusion]
\label{p1:thm:remaining-criterion}
Assume the residual-class hypothesis \Cref{p1:hyp:no-remainder}.  Then no
admissible finite-energy Type I singularity exists.
\end{theorem}

The proof uses one proved closure result and one explicit residual-class
hypothesis.  The uniformly \(L^3\)-tight class is excluded by
\Cref{p2:thm:tight-liouville}.
The remaining class is excluded only under the residual-class hypothesis
\Cref{p1:hyp:no-remainder}.  The structure/decay
class is restricted to hypotheses covered by cited Liouville theorems.

The logical implication is
\[
\begin{array}{c}
\text{admissible finite-energy Type I singularity}
\\
\Downarrow
\\
\text{nonzero normalized Seregin limit}
\\
\Downarrow
\\
\text{small/stationary/tight/structure--decay/remaining alternative}
\\
\Downarrow
\\
\text{contradiction by the resolved classes and the residual-class hypothesis.}
\end{array}
\]

\begin{itemize}[leftmargin=*,itemsep=2pt]
\item Small amplitude: excluded here in \Cref{p1:sec:liouville}.
\item Stationary \(L^3\): excluded here by \cite{NRS1996}.
\item Uniformly \(L^3\)-tight: proved in \Cref{p1:hyp:tight-liouville}, using
\Cref{p2:thm:tight-liouville}.
\item Known structure/decay: excluded by the cited theorems in
\Cref{p1:thm:known-structure-decay-liouville}.
\item Remaining class: excluded under the residual-class hypothesis
\Cref{p1:hyp:no-remainder}.
\end{itemize}

\subsection{Dependencies and conditional hypothesis}

The proof of the extraction theorem uses the local velocity concentration
result stated as \cref{p1:prop:local-velocity-concentration}; this is the only
external result needed for the blow-up construction.  The uniformly
\(L^3\)-tight Liouville theorem and the remaining-class assertion are not used
to construct the ancient limit.  They are invoked only after the class
exhaustion.  The statement \cref{p1:hyp:tight-liouville} follows from the
endpoint \(L^3\) closure in \Cref{p2:thm:tight-liouville}, and
\cref{p1:hyp:no-remainder} is the residual-class hypothesis.  The
structure/decay class is restricted by definition to assumptions covered by
the cited Liouville theorems.  Thus the final Type I conclusion is conditional
on the admissibility condition in \Cref{p0:cor:auto-paperI-admissibility}, the
extraction proved here, the resolved Liouville-class results cited above, and the
residual-class hypothesis.

\subsection{Organization of the paper}

\Cref{p1:sec:prelim} introduces suitable solutions, scale-invariant quantities,
the admissibility condition, scaling, pressure gauges, and local concentration.
\Cref{p1:sec:blowup-compactness} proves the compactness theorem for Type I
rescalings away from the terminal time and constructs the ancient suitable
weak limit.  \Cref{p1:sec:nontriviality} proves persistence of concentration and
nontriviality.  \Cref{p1:sec:seregin} derives the centered equation and defines
normalized Seregin limits and collections.  \Cref{p1:sec:classes} defines the
Liouville classes and the remaining class.  \Cref{p1:sec:liouville} proves the
class exclusions and the uniformly tight theorem.  \Cref{p1:sec:criterion} proves
the conditional remaining-case criterion.  Appendix \ref{p1:app:centered-symmetries} gives
the elementary symmetries of the centered equation.

\section{Suitable solutions, scale-invariant quantities, and local concentration}
\label{p1:sec:prelim}

For \(z_0=(x_0,t_0)\), write
\[
   Q_r(z_0)=B_r(x_0)\times(t_0-r^2,t_0).
\]
When \(z_0=(0,0)\), write \(Q_r=Q_r(0,0)\).  The spatial average of \(f\) over
\(B_r(x_0)\) is
\[
   (f)_{B_r(x_0)}(t)=\frac1{|B_r|}\int_{B_r(x_0)}f(x,t)\,dx.
\]
Throughout the paper, \(\mathbb P\) denotes the Helmholtz--Leray projection
onto divergence-free vector fields on \(\R^3\).

\begin{definition}[Terminal singular point]\label{p1:def:terminal-singular}
Let \((u,p)\) be a suitable weak solution on \(\R^3\times(0,T)\).  A terminal
point \(z_*=(x_*,T)\) is regular if there are \(r>0\) and \(K<\infty\) such
that
\[
   \esssup_{Q_r(z_*)}|u|\le K.
\]
It is singular otherwise.
\end{definition}

\begin{definition}[Suitable weak solution]\label{p1:def:suitable}
Let \(\Omega\subset\R^3\times\R\) be open.  A pair \((u,p)\) is a suitable
weak solution of \eqref{p1:eq:NS} in \(\Omega\) if
\[
   u\in L^\infty_tL^2_x(\Omega')\cap L^2_t\dot H^1_x(\Omega'),
   \qquad
   p\in L^{3/2}(\Omega')
\]
for every compact cylinder \(\Omega'\Subset\Omega\), if \((u,p)\) solves
\eqref{p1:eq:NS} distributionally, \(\divergence u=0\), and if the local energy
inequality
\begin{align}\label{p1:eq:lei}
   &\int_{\R^3}|u(x,t_2)|^2\phi(x,t_2)\,dx
   +2\int_{t_1}^{t_2}\int_{\R^3}|\nabla u|^2\phi\,dx\,dt  \notag\\
   &\le
   \int_{\R^3}|u(x,t_1)|^2\phi(x,t_1)\,dx
   +\int_{t_1}^{t_2}\int_{\R^3}|u|^2(\partial_t\phi+\Delta\phi)\,dx\,dt
   \notag\\
   &\quad
   +\int_{t_1}^{t_2}\int_{\R^3}(|u|^2+2p)u\cdot\nabla\phi\,dx\,dt
\end{align}
holds for every non-negative \(\phi\in C_c^\infty(\Omega)\) and almost every
\(t_1<t_2\) for which the displayed terms are defined.
\end{definition}

This is the Caffarelli--Kohn--Nirenberg notion of suitability
\cite{CKN1982}; the global finite-energy background is Leray's
Leray--Hopf class \cite{Leray1934}.

\begin{definition}[Suitable Leray--Hopf solution on \(\R^3\times(0,T)\)]
\label{p1:def:suitable-leray-hopf}
A suitable Leray--Hopf solution on \(\R^3\times(0,T)\) is a suitable weak
solution on \(\R^3\times(0,T)\) satisfying
\[
   u\in L^\infty(0,T;L^2(\R^3))
   \cap L^2(0,T;\dot H^1(\R^3)),
   \qquad
   u\in C_w([0,T);L^2(\R^3)),
\]
and the global Leray energy inequality.  The pressure is taken in
\(L^{3/2}_{\loc}(\R^3\times(0,T))\), modulo functions of time.
\end{definition}

Define
\begin{equation}\label{p1:eq:C-D-def}
   A(u;z_0,r)=r^{-1}\esssup_{t_0-r^2<t<t_0}
   \int_{B_r(x_0)}|u(x,t)|^2\,dx,
\end{equation}
\begin{equation}\label{p1:eq:C-D-def-2}
   C(u;z_0,r)=r^{-2}\iint_{Q_r(z_0)}|u|^3\,dx\,dt,
   \qquad
   D(p;z_0,r)=r^{-2}\iint_{Q_r(z_0)}
   |p-(p)_{B_r(x_0)}(t)|^{3/2}\,dx\,dt,
\end{equation}
\begin{equation}\label{p1:eq:E-I-def}
   E(u;z_0,r)=r^{-1}\iint_{Q_r(z_0)}|\nabla u|^2\,dx\,dt,
\end{equation}
\begin{definition}[Admissible finite-energy Type I singularity]
\label{p1:def:admissible-type-I}
A triple \((u,p,z_*)\), with \(z_*=(x_*,T)\), is an admissible finite-energy
Type I singularity if:
\begin{enumerate}[label=(\roman*)]
\item \((u,p)\) is a suitable Leray--Hopf solution on
\(\R^3\times(0,T)\);
\item \(z_*=(x_*,T)\) is a singular point;
\item the Type I bound \eqref{p1:eq:type-I-bound} holds;
\item there are scales \(\lambda_k\downarrow0\) and a constant
\(\eta_v>0\) such that
\[
   C(u;z_*,\lambda_k)\ge\eta_v
   \qquad\hbox{for every }k;
\]
\item for the corresponding rescaled velocities
\[
   u_k(x,t)=\lambda_k u(x_*+\lambda_kx,T+\lambda_k^2t),
\]
the velocity no-escape condition holds:
\begin{equation}\label{p1:eq:velocity-no-escape}
   \lim_{\sigma\downarrow0}
   \sup_k
   \int_{-\sigma}^{0}\int_{B_1}|u_k(x,t)|^3\,dx\,dt=0 .
\end{equation}
\end{enumerate}
The sequence \(\{\lambda_k\}\) is called an admissible concentration sequence
at \(z_*\).
\end{definition}

\begin{remark}
Item (iv) is the velocity-only concentration supplied by the local
concentration theorem; see \cref{p1:prop:local-velocity-concentration}.  Thus,
for a singular point, the additional admissibility requirement is the
no-escape condition \eqref{p1:eq:velocity-no-escape}.  This is the only
terminal-slice compactness assumption used to prove that the ancient limit is
nonzero.  No uniform bound on \(A+C+D+E\) over all terminal scales is assumed.
\end{remark}

\begin{lemma}[Uniform kinetic energy implies no escape]
\label{p1:lem:A-seq-no-escape}
Let \((u,p)\) satisfy the Type I bound \eqref{p1:eq:type-I-bound}, and let
\(\lambda_k\downarrow0\).  Suppose that
\[
   \sup_k A(u;z_*,\lambda_k)\le B<\infty .
\]
Then the corresponding rescaled velocities satisfy the no-escape condition
\eqref{p1:eq:velocity-no-escape}.
\end{lemma}

\begin{proof}
Choose \(M<\infty\) and \(t_0<T\) such that
\[
   \|u(t)\|_{L^\infty(\R^3)}
   \le M(T-t)^{-1/2},
   \qquad t_0<t<T.
\]
For \(k\) sufficiently large, \(T+\lambda_k^2t>t_0\) whenever
\(-1<t<0\).  Hence
\[
   \|u_k(t)\|_{L^\infty(B_1)}
   \le M(-t)^{-1/2},\qquad -1<t<0.
\]
By scale invariance of \(A\),
\[
   \esssup_{-1<t<0}\int_{B_1}|u_k(x,t)|^2\,dx
   =
   A(u;z_*,\lambda_k)\le B.
\]
Therefore, for \(0<\sigma<1\),
\[
\begin{aligned}
   \int_{-\sigma}^{0}\int_{B_1}|u_k|^3\,dx\,dt
   &\le
   \int_{-\sigma}^{0}
   \|u_k(t)\|_{L^\infty(B_1)}
   \int_{B_1}|u_k(x,t)|^2\,dx\,dt  \\
   &\le MB\int_{-\sigma}^{0}(-t)^{-1/2}\,dt
    =2MB\sigma^{1/2}.
\end{aligned}
\]
Taking the supremum over the tail \(k\ge k_0\) and then letting
\(\sigma\downarrow0\) gives \eqref{p1:eq:velocity-no-escape} for the tail of the
sequence.  For each of the finitely many remaining indices
\(1\le k<k_0\), the rescaled velocity is defined on
\(B_1\times(-\delta_k,0)\) for some \(\delta_k>0\), and belongs to
\(L^3\) there.  Therefore absolute continuity of the \(L^3\)-integral gives
\[
   \lim_{\sigma\downarrow0}
   \max_{1\le k<k_0}
   \int_{-\sigma}^{0}\int_{B_1}|u_k|^3\,dx\,dt=0.
\]
Combining the finite initial segment with the tail proves the full supremum
in \eqref{p1:eq:velocity-no-escape}.
\end{proof}

\begin{proposition}[Local velocity concentration]
\label{p1:prop:local-velocity-concentration}
Let \((u,p)\) be suitable in a parabolic neighborhood of
\(z_*=(x_*,T)\).  If \(z_*\) is singular, then there exist
\(\lambda_k\downarrow0\) and \(\eta_v>0\) such that
\[
   C(u;z_*,\lambda_k)\ge\eta_v
   \qquad\hbox{for every }k.
\]
\end{proposition}

\begin{proof}
This is the velocity-only concentration corollary of the local concentration
result in \Cref{p0:thm:typeI-typeII-dichotomy}.  That result follows from the
velocity \(\varepsilon\)-regularity criterion: if
\[
   \limsup_{r\downarrow0}C(u;z_*,r)=0,
\]
then \(u\) is bounded in a backward parabolic cylinder ending at \(z_*\),
contradicting singularity of \(z_*\).
\end{proof}

\begin{lemma}[Scaling]\label{p1:lem:scaling}
Let \((u,p)\) be suitable in a neighborhood of \((x_*,T)\).  For
\(\lambda>0\), set
\[
   u^\lambda(x,t)=\lambda u(x_*+\lambda x,T+\lambda^2t),
   \qquad
   p^\lambda(x,t)=\lambda^2p(x_*+\lambda x,T+\lambda^2t).
\]
Then \((u^\lambda,p^\lambda)\) is suitable on the rescaled domain.  Moreover
\(A,C,D,E\) are invariant under this scaling.
\end{lemma}

\begin{proof}
Suitability of the rescaled pair and preservation of pressure gauges follow
by applying \eqref{p1:eq:lei} to the test function
\[
   \phi^\lambda(X,S)
   =
   \phi\left(\frac{X-x_*}{\lambda},\frac{S-T}{\lambda^2}\right)
\]
and then changing variables
\((X,S)=(x_*+\lambda x,T+\lambda^2t)\).  The differential identities
\[
   \partial_tu^\lambda=\lambda^3(\partial_tu)(X,S),\quad
   (u^\lambda\cdot\nabla)u^\lambda
      =\lambda^3((u\cdot\nabla)u)(X,S),
\]
\[
   \nabla p^\lambda=\lambda^3(\nabla p)(X,S),\qquad
   \Delta u^\lambda=\lambda^3(\Delta u)(X,S)
\]
give the rescaled equation and the divergence condition.

For the pressure average,
\[
   (p^\lambda)_{B_r}(t)
   =
   \lambda^2(p)_{B_{\lambda r}(x_*)}(S).
\]
Consequently
\[
   \iint_{Q_r}|u^\lambda|^3\,dx\,dt
   =
   \lambda^{-2}\iint_{Q_{\lambda r}(z_*)}|u|^3\,dX\,dS
\]
and
\[
   \iint_{Q_r}|p^\lambda-(p^\lambda)_{B_r}(t)|^{3/2}\,dx\,dt
   =
   \lambda^{-2}
   \iint_{Q_{\lambda r}(z_*)}
   |p-(p)_{B_{\lambda r}(x_*)}(S)|^{3/2}\,dX\,dS .
\]
Multiplication by \(r^{-2}\) therefore gives invariance of \(C\) and \(D\).
The remaining two quantities have the same critical dimensional balance:
\begin{align*}
   \int_{B_r}|u^\lambda(x,t)|^2\,dx
   &=
   \lambda^{-1}\int_{B_{\lambda r}(x_*)}|u(X,S)|^2\,dX,\\
   \iint_{Q_r}|\nabla u^\lambda|^2\,dx\,dt
   &=
   \lambda^{-1}\iint_{Q_{\lambda r}(z_*)}|\nabla u|^2\,dX\,dS .
\end{align*}
Multiplication by \(r^{-1}\) on the rescaled cylinder is the same as
multiplication by \((\lambda r)^{-1}\) on the original cylinder.  Hence
\(A\) and \(E\) are invariant as well.
\end{proof}

\begin{lemma}[Pressure gauges do not change the equations]
\label{p1:lem:gauges}
Let \((v,q)\) be suitable in a cylinder \(Q\), and let
\(\alpha(t)\in L^{3/2}_{\loc}\) depend only on time.  Then
\((v,q-\alpha(t))\) satisfies the same distributional Navier--Stokes equation,
the same divergence condition, and the same local energy inequality.
\end{lemma}

\begin{proof}
In the distributional equation, replacing \(q\) by \(q-\alpha(t)\) adds
\(-\nabla\alpha(t)=0\).  In the local energy inequality the only new term is
\[
   -2\int \alpha(t)v(x,t)\cdot\nabla\phi(x,t)\,dx\,dt .
\]
The product is integrable on the support of \(\phi\): by suitability,
\(v\in L^\infty_tL^2_x\) locally, and
\[
   \left|\int v(\cdot,t)\cdot\nabla\phi(\cdot,t)\,dx\right|
   \le
   \|v(t)\|_{L^2(\supp\phi(t))}
   \|\nabla\phi(t)\|_{L^2}.
\]
For almost every \(t\), \(\divergence v(\cdot,t)=0\) in distributions and
\(\phi(\cdot,t)\) is compactly supported, so
\[
   \int v(x,t)\cdot\nabla\phi(x,t)\,dx=0 .
\]
Thus the additional term vanishes.
\end{proof}

\section{The Type I blow-up sequence and compactness}
\label{p1:sec:blowup-compactness}

\subsection{The rescaled sequence}

Choose \(M<\infty\) so that, after decreasing \(T_0<T\) if necessary,
\begin{equation}\label{p1:eq:type-I-pointwise}
   \|u(t)\|_{L^\infty(\R^3)}
   \le \frac{M}{\sqrt{T-t}},
   \qquad T_0<t<T.
\end{equation}
For \(\lambda_k\downarrow0\), define
\begin{equation}\label{p1:eq:rescaled-sequence}
   u_k(x,t)=\lambda_k u(x_*+\lambda_k x,T+\lambda_k^2t),
   \qquad
   p_k(x,t)=\lambda_k^2p(x_*+\lambda_k x,T+\lambda_k^2t).
\end{equation}
Then \(u_k\) is defined on \(\R^3\times(-T/\lambda_k^2,0)\), and these
domains exhaust \(\R^3\times(-\infty,0)\).

\subsection{Inherited Type I bounds}

\begin{lemma}[Inherited Type I bound]\label{p1:lem:inherited-bound}
For every compact interval \(I\Subset(-\infty,0)\),
\[
   \sup_k\|u_k\|_{L^\infty(\R^3\times I)}<\infty.
\]
More precisely, for every fixed \(t<0\) and all sufficiently large \(k\),
\[
   \|u_k(t)\|_{L^\infty(\R^3)}
   \le \frac{M}{\sqrt{-t}}.
\]
\end{lemma}

\begin{proof}
Since \(T-(T+\lambda_k^2t)=\lambda_k^2(-t)\), \eqref{p1:eq:type-I-pointwise}
gives, for all large \(k\),
\[
   \|u_k(t)\|_\infty
   =
   \lambda_k\|u(T+\lambda_k^2t)\|_\infty
   \le
   \lambda_k\frac{M}{\sqrt{\lambda_k^2(-t)}}
   =
   \frac{M}{\sqrt{-t}}.
\]
Taking the supremum over \(t\in I\) gives the interval statement.
\end{proof}

\subsection{Local energy, pressure, and time-derivative estimates}

\begin{lemma}[Pressure gauges and local pressure bounds]\label{p1:lem:pressure}
After fixing the whole-space Leray pressure representative below, set
\[
   a_k(t)=(p_k)_{B_1}(t)
\]
whenever the rescaled solution is defined.  Then for every
\[
   Q=B_R\times[-S,-\sigma]\Subset\R^3\times(-\infty,0)
\]
one has
\[
   \sup_k\|p_k-a_k(t)\|_{L^{3/2}(Q)}<\infty .
\]
\end{lemma}

\begin{proof}
For large \(k\), the original times \(T+\lambda_k^2t\), \(t\in[-S,-\sigma]\),
lie in the interval where \eqref{p1:eq:type-I-pointwise} holds.  On those slices
\(u(t)\in L^2(\R^3)\cap L^\infty(\R^3)\), hence
\[
   u_i(t)u_j(t)\in L^1(\R^3)\cap L^{3/2}(\R^3).
\]
Use the whole-space Leray pressure representative
\[
   p^L=\mathcal R_i\mathcal R_j(u_i u_j),
\]
modulo functions of time.  We use the standard convention
\[
   \widehat{\mathcal R_i f}(\xi)=i\xi_i|\xi|^{-1}\widehat f(\xi),
\]
so that \(\mathcal R_i\mathcal R_j\) has multiplier
\(-\xi_i\xi_j|\xi|^{-2}\) and
\(-\Delta(\mathcal R_i\mathcal R_j f)=\partial_i\partial_jf\).  Thus this is
the pressure reconstruction from
\(-\Delta p=\partial_i\partial_j(u_i u_j)\) and the Calderon--Zygmund theory
of Riesz transforms \cite{Stein1970,Sohr2001}.  Since
\(u(t)\in L^2(\R^3)\cap L^\infty(\R^3)\), interpolation gives
\(u(t)\in L^3(\R^3)\), hence \(u_i(t)u_j(t)\in L^{3/2}(\R^3)\) and
\(p^L(t)\in L^{3/2}(\R^3)\) for a.e. \(t\).

The given pressure and \(p^L\) differ by a function of time.  Indeed, applying
the Helmholtz projection to the whole-space weak momentum equation gives
\[
   \partial_tu-\Delta u+\mathbb P\divergence(u\otimes u)=0
\]
in solenoidal distributions.  The complementary gradient part of
\(\divergence(u\otimes u)\) is precisely
\(-\nabla p^L\) with the convention above.  Therefore the same velocity field
solves the distributional equation with \(p^L\) in place of \(p\).  Subtracting
the two formulations gives
\[
   \nabla(p-p^L)=0
   \quad\hbox{in }\mathcal D'(\R^3)
\]
for a.e. time.  By the de Rham theorem on \(\R^3\), \(p-p^L\) is spatially
constant for a.e. time.  Replacing the pressure by \(p^L\) is therefore a
time-dependent gauge change and does not affect suitability by
\cref{p1:lem:gauges}.

Fix \(\chi_R\in C_c^\infty(B_{8R})\) with \(\chi_R=1\) on \(B_{4R}\).  For
\(t\in[-S,-\sigma]\), choose the representative so that, in \(B_{2R}\),
\[
   p_k=q_k+h_k,\qquad
   q_k=\mathcal R_i\mathcal R_j(\chi_Ru_{k,i}u_{k,j}),
\]
and \(h_k(\cdot,t)\) is harmonic in \(B_{4R}\) for a.e. \(t\).  This choice
is made before \(a_k(t)=(p_k)_{B_1}(t)\) is evaluated.  The harmonicity holds
because \((1-\chi_R)u_{k,i}u_{k,j}\) is supported outside \(B_{4R}\), so the
localized Riesz potential generated by this exterior part has zero Laplacian
in \(B_{4R}\).  Calderon--Zygmund boundedness gives, for a.e. \(t\),
\[
   \|q_k(t)\|_{L^{3/2}(B_{2R})}
   \le C\|u_k(t)\|_{L^3(B_{8R})}^2
   \le C(R,\sigma,M).
\]
Integrating over \([-S,-\sigma]\) gives a uniform
\(L^{3/2}_{x,t}\)-bound for \(q_k\).
Let \(c_{k,R}(t)=h_k(0,t)\), where \(h_k(\cdot,t)\) is represented by its
smooth harmonic representative.  For \(x\in B_R\), the kernel estimate
\[
   |K_{ij}(x-z)-K_{ij}(-z)|\le C R|z|^{-4},\qquad |z|>4R,
\]
implies
\[
   |h_k(x,t)-c_{k,R}(t)|
   \le C R\sum_{m=2}^\infty(2^mR)^{-4}
   \int_{2^mR<|z|<2^{m+1}R}|u_k(z,t)|^2\,dz .
\]
Using \cref{p1:lem:inherited-bound},
\[
   \int_{2^mR<|z|<2^{m+1}R}|u_k(z,t)|^2\,dz
   \le C M^2\sigma^{-1}(2^mR)^3,
\]
and therefore
\[
   |h_k(x,t)-c_{k,R}(t)|
   \le C M^2\sigma^{-1}\sum_{m=2}^{\infty}2^{-m},
   \qquad x\in B_R,
\]
which implies
\[
   \|h_k(t)-c_{k,R}(t)\|_{L^\infty(B_R)}\le C(M,\sigma).
\]
Thus
\[
   \sup_k\|p_k-c_{k,R}(t)\|_{L^{3/2}(B_R\times[-S,-\sigma])}<\infty .
\]
If \(R<1\), enlarge \(R\) to \(1\).  Otherwise \(B_1\subset B_R\).  Jensen's
inequality gives, for a.e. \(t\),
\[
   |a_k(t)-c_{k,R}(t)|^{3/2}
   \le
   |B_1|^{-1}\int_{B_1}|p_k-c_{k,R}(t)|^{3/2}\,dx .
\]
Integrating in time and adding constants on \(B_R\) gives the desired bound
for \(p_k-a_k(t)\).
\end{proof}

\begin{lemma}[Compatibility of pressure gauges]\label{p1:lem:gauge-compatibility}
Let \(B_1\subset B_R\), \(I\Subset(-\infty,0)\), and suppose
\[
   p_k-c_k(t)
   \quad\hbox{is bounded in }L^{3/2}(B_R\times I)
\]
for some time-dependent gauges \(c_k(t)\).  If
\(a_k(t)=(p_k)_{B_1}(t)\), then \(p_k-a_k(t)\) is bounded in
\(L^{3/2}(B_{R'}\times I)\) for every \(R'\le R\).  Changing between such
gauges changes the pressure only by a function of time and therefore does not
affect the equation, the local energy inequality, or the centered equation.
\end{lemma}

\begin{proof}
For a.e. \(t\), Jensen's inequality gives
\[
   |a_k(t)-c_k(t)|^{3/2}
   \le |B_1|^{-1}\int_{B_1}|p_k-c_k(t)|^{3/2}\,dx .
\]
Integrating in \(t\) and using \(B_{R'}\subset B_R\) yields
\[
   \|a_k-c_k\|_{L^{3/2}(I)}^{3/2}
   \le
   |B_1|^{-1}\|p_k-c_k(t)\|_{L^{3/2}(B_R\times I)}^{3/2}.
\]
Thus
\[
   \|p_k-a_k(t)\|_{L^{3/2}(B_{R'}\times I)}
   \le
   \|p_k-c_k(t)\|_{L^{3/2}(B_{R'}\times I)}
   +|B_{R'}|^{2/3}\|a_k-c_k\|_{L^{3/2}(I)}
\]
and the right-hand side is bounded by
\[
   C(R')\|p_k-c_k(t)\|_{L^{3/2}(B_R\times I)} .
\]
This proves the asserted bound.
The last assertion is \cref{p1:lem:gauges}, and centering transforms the
time-dependent gauge into another time-dependent gauge.
\end{proof}

\begin{lemma}[Local energy and dissipation]\label{p1:lem:energy-diss}
Let
\[
   Q=B_R\times[-S,-\sigma]\Subset\R^3\times(-\infty,0).
\]
Then
\[
   \sup_k\|u_k\|_{L^\infty_tL^2_x(Q)}
   +\sup_k\|\nabla u_k\|_{L^2(Q)}<\infty .
\]
\end{lemma}

\begin{proof}
The \(L^\infty_tL^2_x\) estimate follows from
\[
   \|u_k(t)\|_{L^2(B_R)}
   \le |B_R|^{1/2}M\sigma^{-1/2},
   \qquad t\in[-S,-\sigma].
\]
Choose \(\phi\in C_c^\infty(B_{2R}\times[-2S,-\sigma/2])\), \(0\le\phi\le1\),
with \(\phi=1\) on \(Q\).  Apply the local energy inequality to
\((u_k,p_k-a_k(t))\), using \cref{p1:lem:pressure} on the larger cylinder.  For
a.e. admissible times and then by monotone approximation of temporal cutoffs,
\[
   2\iint |\nabla u_k|^2\phi
   \le
   \iint |u_k|^2|\partial_t\phi+\Delta\phi|
   +\iint (|u_k|^2+2|p_k-a_k(t)|)|u_k||\nabla\phi|.
\]
The first term is bounded because \(\partial_t\phi+\Delta\phi\) is bounded
and \(u_k\) is uniformly bounded in \(L^\infty\) on the finite-measure support
of \(\phi\).  The cubic velocity term satisfies
\[
   \iint |u_k|^3|\nabla\phi|
   \le
   \|\nabla\phi\|_\infty
   |\supp\nabla\phi|\,\|u_k\|_{L^\infty(\supp\nabla\phi)}^3
   \le C(Q,\phi,M,\sigma).
\]
For the pressure term, H\"older's inequality on \(\supp\nabla\phi\) gives
\[
   \iint |p_k-a_k(t)|\,|u_k|\,|\nabla\phi|
   \le
   \|p_k-a_k(t)\|_{L^{3/2}}
   \|u_k\nabla\phi\|_{L^3},
\]
and both factors are uniformly bounded by \cref{p1:lem:pressure} and
\cref{p1:lem:inherited-bound}.  Since \(\phi=1\) on \(Q\), the dissipation bound
follows.
\end{proof}

\begin{lemma}[Time derivative bound]\label{p1:lem:time-derivative}
For every compact cylinder \(Q\Subset\R^3\times(-\infty,0)\),
\[
   \partial_tu_k
   \quad\hbox{is bounded in }L^{3/2}_tW^{-1,3/2}_x(Q).
\]
\end{lemma}

\begin{proof}
The equation is
\[
   \partial_tu_k
   =
   \Delta u_k-\divergence(u_k\otimes u_k)-\nabla(p_k-a_k(t)).
\]
Let \(Q=B_R\times[-S,-\sigma]\).  The term \(\Delta u_k\) is bounded in
\(L^2_tH^{-1}_x(Q)\) by \cref{p1:lem:energy-diss}; since
\(W^{1,3}_0(B_R)\subset H^1_0(B_R)\) and the time interval is finite, this is
also an \(L^{3/2}_tW^{-1,3/2}_x\)-bound.  Explicitly, for
\(\zeta\in W^{1,3}_0(B_R)\),
\[
   |\langle\Delta u_k,\zeta\rangle|
   =
   \left|\int_{B_R}\nabla u_k:\nabla\zeta\,dx\right|
   \le
   C(R)\|\nabla u_k\|_{L^2(B_R)}\|\zeta\|_{W^{1,3}(B_R)}.
\]
The \(L^2_t\)-bound on \(\nabla u_k\), together with the finite length of
\([-S,-\sigma]\), gives the claimed \(L^{3/2}_t\)-bound.

For the nonlinear term, for
\(\zeta\in W^{1,3}_0(B_R)\),
\[
   \left|\int_{B_R}(u_k\otimes u_k):\nabla\zeta\,dx\right|
   \le
   \|u_k\otimes u_k\|_{L^{3/2}(B_R)}
   \|\nabla\zeta\|_{L^3(B_R)}.
\]
Moreover
\[
   \|u_k\otimes u_k\|_{L^{3/2}(Q)}
   \le
   \|u_k\|_{L^\infty(Q)}\|u_k\|_{L^{3/2}(Q)}
   \le C(Q).
\]
Thus \(\divergence(u_k\otimes u_k)\) is bounded in
\(L^{3/2}_tW^{-1,3/2}_x(Q)\).
For the pressure term,
\[
   \left|\int_{B_R}(p_k-a_k(t))\,\divergence\zeta\,dx\right|
   \le
   \|p_k-a_k(t)\|_{L^{3/2}(B_R)}
   \|\nabla\zeta\|_{L^3(B_R)}.
\]
\Cref{p1:lem:pressure} gives
\[
   p_k-a_k(t)\quad\hbox{bounded in }L^{3/2}(Q).
\]
Hence \(\nabla(p_k-a_k(t))\) is bounded in
\(L^{3/2}_tW^{-1,3/2}_x(Q)\).  Summing the three terms gives the result.
\end{proof}

\subsection{Compactness away from the terminal time}

\begin{proposition}[Compactness]\label{p1:prop:compactness}
After passing to a subsequence,
\[
   u_k\to U
\]
strongly in \(L^q_{\loc}(\R^3\times(-\infty,0))\) for every finite \(q\), and
\[
   p_k-a_k(t)\rightharpoonup Q
   \quad\hbox{weakly in }L^{3/2}_{\loc}(\R^3\times(-\infty,0)).
\]
\end{proposition}

\begin{proof}
Fix \(Q_0\Subset\R^3\times(-\infty,0)\).  By
\cref{p1:lem:energy-diss}, \(u_k\) is bounded in \(L^2_tH^1_x(Q_0)\), and by
\cref{p1:lem:time-derivative}, \(\partial_tu_k\) is bounded in
\(L^{3/2}_tW^{-1,3/2}_x(Q_0)\).  On bounded balls,
\[
   H^1\Subset L^2\hookrightarrow W^{-1,3/2}.
\]
The Aubin--Lions--Simon theorem \cite{Aubin1963,LionsMagenes,Simon1987} gives
relative compactness in \(L^2(Q_0)\).  Hence, after passing to a subsequence,
\(u_k\to U\) strongly in \(L^2(Q_0)\) and a.e. on \(Q_0\).

The inherited Type I bound gives a common \(L^\infty(Q_0)\) bound, and the
limit belongs to the same \(L^\infty\) ball by a.e. convergence.  For
\(q\ge2\),
\[
   \|u_k-U\|_{L^q(Q_0)}^q
   \le
   (2\sup_j\|u_j\|_{L^\infty(Q_0)}+2\|U\|_{L^\infty(Q_0)})^{q-2}
   \|u_k-U\|_{L^2(Q_0)}^2.
\]
For \(1\le q<2\), H\"older's inequality on the finite-measure set \(Q_0\)
gives
\[
   \|u_k-U\|_{L^q(Q_0)}
   \le |Q_0|^{1/q-1/2}\|u_k-U\|_{L^2(Q_0)}\to0.
\]
Choose an exhaustion \(Q^{(m)}\Subset\R^3\times(-\infty,0)\) by compact
cylinders.  The preceding argument gives a subsequence converging on
\(Q^{(1)}\); extracting successively on \(Q^{(m)}\) and taking the diagonal
subsequence gives strong local convergence on every compact cylinder.

For pressure, \cref{p1:lem:pressure} gives boundedness of
\(p_k-a_k(t)\) in \(L^{3/2}(Q^{(m)})\) for every \(m\).  Since
\(L^{3/2}\) is reflexive, each restriction has a weakly convergent
subsequence.  The same diagonal extraction gives a single limit \(Q\) in
\(L^{3/2}_{\loc}\), because the weak limits agree on overlaps by uniqueness
of distributional limits.
\end{proof}

\subsection{The ancient suitable weak limit}

\begin{lemma}[Limit equation]\label{p1:lem:limit-equation}
The limit \((U,Q)\) solves
\[
   \partial_tU+(U\cdot\nabla)U+\nabla Q=\Delta U,\qquad
   \divergence U=0
\]
in distributions on \(\R^3\times(-\infty,0)\).
\end{lemma}

\begin{proof}
For \(\varphi\in C_c^\infty(\R^3\times(-\infty,0);\R^3)\), the weak
formulation for \((u_k,p_k-a_k(t))\) is
\[
   \iint\left[
      -u_k\cdot\partial_t\varphi
      -(u_k\otimes u_k):\nabla\varphi
      +(p_k-a_k(t))\divergence\varphi
      +\nabla u_k:\nabla\varphi
   \right]\,dx\,dt=0.
\]
The linear time term passes by strong \(L^2_{\loc}\)-convergence of \(u_k\).
The pressure term passes by weak \(L^{3/2}_{\loc}\)-convergence of
\(p_k-a_k(t)\), since \(\divergence\varphi\in L^3\).  The viscous term passes
by weak \(L^2_{\loc}\)-convergence of \(\nabla u_k\), which follows from the
uniform local dissipation bound.  The nonlinear term
passes because
\[
   \|u_k\otimes u_k-U\otimes U\|_{L^1(\supp\varphi)}
   \le
   \|u_k-U\|_{L^2}
   \bigl(\|u_k\|_{L^2}+\|U\|_{L^2}\bigr)\to0.
\]
Testing with gradients of scalar functions gives \(\divergence U=0\).
\end{proof}

\begin{lemma}[Stability of suitability]\label{p1:lem:suitability}
The pair \((U,Q)\) is a suitable weak solution on
\(\R^3\times(-\infty,0)\).
\end{lemma}

\begin{proof}
Use the distributional form of the local energy inequality.  For every
non-negative \(\phi\in C_c^\infty(\R^3\times(-\infty,0))\),
\[
   2\iint |\nabla u_k|^2\phi
   \le
   \iint |u_k|^2(\partial_t\phi+\Delta\phi)
   +\iint (|u_k|^2+2P_k)u_k\cdot\nabla\phi ,
   \qquad P_k=p_k-a_k(t).
\]
This follows from \eqref{p1:eq:lei} by applying \eqref{p1:eq:lei} with temporal
cutoffs which increase to \(1\) on the time projection of \(\supp\phi\).  The
monotone convergence theorem passes the nonnegative dissipation term to the
displayed distributional form.

The terms with \(|u_k|^2\) converge strongly in \(L^1\) because
\[
   \||u_k|^2-|U|^2\|_{L^1(\supp\phi)}
   \le
   \|u_k-U\|_{L^2}
   \bigl(\|u_k\|_{L^2}+\|U\|_{L^2}\bigr)\to0.
\]
The cubic flux converges strongly since \(u_k\to U\) strongly in
\(L^3_{\loc}\), hence
\[
   |u_k|^2u_k\to |U|^2U
   \quad\hbox{in }L^1_{\loc}.
\]
For the pressure flux, \(u_k\cdot\nabla\phi\to U\cdot\nabla\phi\) strongly in
\(L^3\), while \(P_k\rightharpoonup Q\) weakly in \(L^{3/2}\).  Therefore
\[
   \iint P_k u_k\cdot\nabla\phi\to
   \iint Q\,U\cdot\nabla\phi .
\]
Finally,
\(\nabla u_k\rightharpoonup\nabla U\) in \(L^2_{\loc}\), and
\[
   w\mapsto\iint |w|^2\phi
\]
is weakly lower semicontinuous because \(\phi\ge0\) and the expression equals
\(\|\sqrt{\phi}\,w\|_{L^2}^2\).  Passing to the limit gives the local energy
inequality for \((U,Q)\) in distributional form.  The endpoint form in
\cref{p1:def:suitable} is recovered by applying the distributional inequality to
standard time cutoffs approximating characteristic functions of
\((t_1,t_2)\); weak \(L^2_{\loc}\)-continuity of local energy-class solutions
gives the endpoint traces for a.e. \(t_1<t_2\).
\end{proof}

\begin{lemma}[Inherited ancient Type I bound]\label{p1:lem:limit-bound}
For every compact interval \(I=[-S,-\sigma]\Subset(-\infty,0)\),
\[
   \|U\|_{L^\infty(\R^3\times I)}\le M\sigma^{-1/2}.
\]
\end{lemma}

\begin{proof}
\Cref{p1:lem:inherited-bound} gives
\[
   \|u_k\|_{L^\infty(\R^3\times I)}\le M\sigma^{-1/2}.
\]
Fix \(R<\infty\).  Since \(u_k\to U\) a.e. on \(B_R\times I\), the essential
supremum of \(U\) on \(B_R\times I\) cannot exceed \(M\sigma^{-1/2}\).  If it
did, \(|u_k|\) would exceed this bound by a positive amount on a positive
measure subset for large \(k\).  Letting \(R\to\infty\) through integers gives
the result on \(\R^3\times I\).
\end{proof}

\begin{proposition}[Ancient suitable Type I limit]\label{p1:prop:ancient-limit}
The pair \((U,Q)\) is an ancient suitable weak solution on
\(\R^3\times(-\infty,0)\) satisfying the interval Type I bound in
\cref{p1:lem:limit-bound}.
\end{proposition}

\begin{proof}
Combine \cref{p1:lem:limit-equation,p1:lem:suitability,p1:lem:limit-bound}.
\end{proof}

\section{Persistence of concentration and nontriviality}
\label{p1:sec:nontriviality}

\subsection{Persistence of concentration away from the terminal time}

The persistence of concentration is tracked through the velocity concentration sequence in
\cref{p1:def:admissible-type-I}, supplied by
\cref{p1:prop:local-velocity-concentration}.  The no-escape condition says that
this velocity mass does not concentrate entirely in the terminal layer \(t=0\)
after rescaling.

\subsection{Scale selection}

\begin{lemma}[Scale selection from local concentration]\label{p1:lem:scale-selection}
Let \(\lambda_k\downarrow0\) be an admissible concentration sequence in the
sense of \cref{p1:def:admissible-type-I}.  Then there are
\(\sigma\in(0,1)\), a compact cylinder
\[
   K_0=B_1\times(-1,-\sigma)\Subset\R^3\times(-\infty,0),
\]
and \(\eta_1>0\) such that
\begin{equation}\label{p1:eq:selected-lower-bound}
   \iint_{K_0}|u_k|^3\,dx\,dt\ge\eta_1
\end{equation}
for all \(k\).
\end{lemma}

\begin{proof}
By scale invariance of \(C\),
\[
   \int_{-1}^{0}\int_{B_1}|u_k|^3\,dx\,dt
   =
   C(u;z_*,\lambda_k)\ge\eta_v .
\]
The no-escape condition \eqref{p1:eq:velocity-no-escape} gives
\(\sigma\in(0,1)\) such that
\[
   \sup_k\int_{-\sigma}^{0}\int_{B_1}|u_k|^3\,dx\,dt
   \le \frac{\eta_v}{2}.
\]
Therefore
\[
   \int_{-1}^{-\sigma}\int_{B_1}|u_k|^3\,dx\,dt
   \ge \frac{\eta_v}{2}
\]
for every \(k\).  The assertion follows with
\(\eta_1=\eta_v/2\).
\end{proof}

\subsection{Nonzero ancient limits}

\begin{proposition}[Nonzero ancient limit]\label{p1:prop:nonzero}
For the scales chosen in \cref{p1:lem:scale-selection}, the limit \(U\) in
\cref{p1:prop:ancient-limit} is nonzero.  More precisely, there are a compact
cylinder \(K_0\Subset\R^3\times(-\infty,0)\) and \(\eta_1>0\) such that
\[
   \iint_{K_0}|U|^3\,dx\,dt\ge\eta_1 .
\]
\end{proposition}

\begin{proof}
Let \(K_0=B_1\times(-1,-\sigma)\) and \(\eta_1>0\) be given by
\cref{p1:lem:scale-selection}.  Since
\(K_0\Subset\R^3\times(-\infty,0)\), the compactness theorem gives, after
passing to the subsequence used to define \(U\),
\[
   u_k\to U\qquad\hbox{strongly in }L^3(K_0).
\]
Hence
\[
   \int_{K_0}|U|^3\,dx\,dt
   =
   \lim_{k\to\infty}\int_{K_0}|u_k|^3\,dx\,dt
   \ge\eta_1 .
\]
\end{proof}

This is the step where concentration generated by the singular point is shown to
propagate through the compactness limit.  Without the velocity concentration and
no-escape condition, the compactness theorem alone could produce the zero
ancient solution.

\section{Centered variables and normalized Seregin limits}\label{p1:sec:seregin}

\subsection{The centered self-similar equation}

Set
\[
   y=\frac{x}{\sqrt{-t}},\qquad
   \tau=-\log(-t),\qquad
   V(y,\tau)=\sqrt{-t}\,U(y\sqrt{-t},t),
\]
and
\[
   \Pi(y,\tau)=(-t)Q(y\sqrt{-t},t),
\]
up to addition of a function of \(\tau\).

\begin{lemma}[Renormalized equation]\label{p1:lem:renormalized-equation}
The pair \((V,\Pi)\) solves \eqref{p1:eq:centered}.
\end{lemma}

\begin{proof}
Write \(s=-t=e^{-\tau}\), \(x=s^{1/2}y\), and
\[
   U(x,t)=s^{-1/2}V(y,\tau).
\]
The identities are first obtained in distributions by testing the equation
for \(U\) against the inverse change of variables.  After
\cref{p1:lem:smoothness}, the same identities hold classically for the smooth
representative.  Because \(s_t=-1\),
\(\tau_t=s^{-1}\), and \(y_t=(2s)^{-1}y\),
\[
   \partial_tU
   =
   s^{-3/2}\left(\partial_\tau V+\frac12V+\frac12y\cdot\nabla V\right),
\]
\[
   (U\cdot\nabla_x)U=s^{-3/2}(V\cdot\nabla)V,\qquad
   \Delta_xU=s^{-3/2}\Delta V,\qquad
   \nabla_xQ=s^{-3/2}\nabla\Pi.
\]
Substitution into the equation for \(U\) gives \eqref{p1:eq:centered}, and the
divergence condition is preserved by the spatial scaling.
\end{proof}

\subsection{Boundedness, smoothness, and local energy}

\begin{lemma}[Smoothness]\label{p1:lem:smoothness}
The ancient solution \(U\) is smooth on \(\R^3\times(-\infty,0)\), and
\((V,\Pi)\) is smooth on \(\R^3\times\R\) after fixing a pressure gauge.
\end{lemma}

\begin{proof}
On each compact cylinder \(Q_0\Subset\R^3\times(-\infty,0)\),
\cref{p1:lem:limit-bound} gives \(U\in L^\infty(Q_0)\).  Hence
\(U\in L^s_tL^r_x(Q_0)\) for every finite \(r,s\).  Choose \(r=s=10\), so
\[
   \frac2s+\frac3r=\frac12<1.
\]
Serrin's interior regularity criterion \cite{Serrin1962,Sohr2001} applies and
gives local H\"older regularity of \(U\) on a slightly smaller cylinder.
Since \(Q_0\) was arbitrary, this removes every possible interior singular
point in \(\R^3\times(-\infty,0)\).  The bootstrap starts from
\[
   \partial_tU-\Delta U+\nabla Q=-\divergence(U\otimes U),\qquad
   \divergence U=0,
\]
with \(U\otimes U\) locally bounded after Serrin regularity.  Interior Stokes
estimates on nested cylinders give \(W^{2,1}_q\)-regularity for every finite
\(q\).  Parabolic Sobolev embedding then improves the H\"older regularity of
\(\nabla U\), and differentiating the equation and iterating gives
\[
   U\in C^\infty(\R^3\times(-\infty,0)).
\]
The pressure is recovered locally from
\(-\Delta Q=\partial_i\partial_j(U_iU_j)\) plus a harmonic function and is
smooth after fixing a time-dependent gauge.  The centered change of variables
is a smooth diffeomorphism for \(t<0\), so \(V\) and \(\Pi\) are smooth.
\end{proof}

\begin{corollary}[Classical Type I bound]\label{p1:cor:classical-type-I-bound}
For the smooth representative from \cref{p1:lem:smoothness},
\[
   \|U(t)\|_{L^\infty(\R^3)}\le \frac{M}{\sqrt{-t}},
   \qquad t<0.
\]
\end{corollary}

\begin{proof}
Fix \(t_0<0\).  Choose \(\delta_n\downarrow0\) with
\[
   I_n=[t_0-\delta_n,t_0+\delta_n]\Subset(-\infty,0).
\]
\Cref{p1:lem:limit-bound} gives
\[
   \|U\|_{L^\infty(\R^3\times I_n)}
   \le M(-t_0-\delta_n)^{-1/2}.
\]
Since the representative is continuous, the essential spacetime bound holds
pointwise.  Letting \(n\to\infty\) gives the estimate at \(t_0\).
\end{proof}

\begin{lemma}[Renormalized boundedness and nontriviality]
\label{p1:lem:renorm-bounds}
The centered field satisfies
\[
   \|V\|_{L^\infty(\R^3\times\R)}\le M.
\]
Moreover \eqref{p1:eq:nonvanishing} holds for some compact cylinder
\(K\Subset\R^3\times\R\) and \(\eta_0>0\).
\end{lemma}

\begin{proof}
The \(L^\infty\) bound follows from \cref{p1:cor:classical-type-I-bound}:
\[
   |V(y,\tau)|=\sqrt{-t}|U(y\sqrt{-t},t)|\le M.
\]
For nontriviality, \cref{p1:prop:nonzero} gives a positive local \(L^3\) lower bound
on a compact cylinder in \((x,t)\).  Under
\((x,t)=(e^{-\tau/2}y,-e^{-\tau})\),
\[
   dx\,dt=e^{-5\tau/2}\,dy\,d\tau,\qquad
   U=e^{\tau/2}V,
\]
and hence
\[
   |U|^3\,dx\,dt=e^{-\tau}|V|^3\,dy\,d\tau .
\]
On the compact image cylinder the factor \(e^{-\tau}\) is bounded above and
below by positive constants, so the positive lower bound becomes
\eqref{p1:eq:nonvanishing}, with a possibly different \(\eta_0\).
\end{proof}

\begin{lemma}[Local energy bounds in centered variables]
\label{p1:lem:centered-local-energy}
Let \(V\) be obtained from an admissible finite-energy Type I singularity.  For every
\(R<\infty\) and
compact interval \(I\subset\R\),
\[
   \int_I\int_{B_R}|V|^2\,dy\,d\tau
   +
   \int_I\int_{B_R}|\nabla V|^2\,dy\,d\tau<\infty .
\]
After subtracting a function of \(\tau\), the associated pressure satisfies
\[
   \Pi-b(\tau)\in L^{3/2}(B_R\times I).
\]
\end{lemma}

\begin{proof}
The unrenormalized limit \((U,Q)\) is suitable and has local energy,
dissipation, and pressure bounds on compact negative-time cylinders by
\cref{p1:lem:energy-diss,p1:lem:pressure,p1:prop:ancient-limit}; the bounds for the
limit follow from the weak lower semicontinuity and weak compactness already
used in \cref{p1:lem:suitability}.  On each compact
\((y,\tau)\)-cylinder, the map
\[
   (y,\tau)\mapsto (x,t)=(e^{-\tau/2}y,-e^{-\tau})
\]
has Jacobian \(dx\,dt=e^{-5\tau/2}\,dy\,d\tau\), with
\[
   V=e^{-\tau/2}U,\qquad
   \nabla_yV=e^{-\tau}\nabla_xU,\qquad
   \Pi=e^{-\tau}Q .
\]
If \(I=[\tau_1,\tau_2]\), then \(e^{-\tau}\) is bounded above and below by
positive constants on \(I\).  Therefore
\[
   \int_I\int_{B_R}|V|^2\,dy\,d\tau
   \le C(I)\int_{-e^{-\tau_1}}^{-e^{-\tau_2}}
       \int_{B_{C(I)R}}|U|^2\,dx\,dt,
\]
and similarly
\[
   \int_I\int_{B_R}|\nabla_yV|^2\,dy\,d\tau
   \le C(I)\int_{-e^{-\tau_1}}^{-e^{-\tau_2}}
       \int_{B_{C(I)R}}|\nabla_xU|^2\,dx\,dt.
\]
For the pressure, subtract a spatial-average gauge \(a(t)\) for \(Q\) on the
corresponding \(x\)-ball and set \(b(\tau)=e^{-\tau}a(-e^{-\tau})\).  Then
\(\Pi-b=e^{-\tau}(Q-a)\), and the same bounded-weight change of variables
gives \(\Pi-b\in L^{3/2}(B_R\times I)\).
\end{proof}

\begin{remark}\label{p1:rem:no-global-tightness}
\Cref{p1:lem:centered-local-energy} is local.  It does not give uniform
\(L^3\)-tightness in \(y\), nor a global \(L^\infty_\tau L^3_y\) bound.  This
is the reason for retaining the remaining case.
\end{remark}

\subsection{Normalized Seregin ancient limits}

\begin{definition}[Normalized Seregin ancient limit]\label{p1:def:seregin-limit}
A smooth bounded solution \((V,\Pi)\) of \eqref{p1:eq:centered} is a normalized
Seregin ancient limit generated by an admissible finite-energy Type I
singularity \((u,p,z_*)\) if it arises from the blow-up construction of
\cref{p1:thm:extraction} along an admissible concentration sequence and satisfies
the local \(L^3\) nonvanishing condition \eqref{p1:eq:nonvanishing}.
\end{definition}

\begin{definition}[Collection of normalized Seregin limits]
\label{p1:def:seregin-collection}
A collection \(\calS\) is any nonempty set of normalized Seregin ancient
limits generated by admissible finite-energy Type I singularities.  When
needed, \(\calS(u,p,z_*)\) denotes a collection associated with a fixed
singular point \(z_*\).
\end{definition}

\begin{remark}
The local \(L^3\) nonvanishing condition is part of the normalization, not an
additional hypothesis.  It is inherited from the selected velocity
concentration and the no-escape condition through
\cref{p1:lem:scale-selection,p1:prop:nonzero}.  Thus every normalized Seregin
ancient limit is nonzero in the sense of \eqref{p1:eq:nonvanishing}.
\end{remark}

\begin{proof}[Proof of \cref{p1:thm:extraction}]
Choose an admissible concentration sequence
\(\lambda_k\downarrow0\) from \cref{p1:def:admissible-type-I}, and form the
rescaled sequence \eqref{p1:eq:rescaled-sequence}.  The compact ancient limit is
\cref{p1:prop:ancient-limit}; the strong and weak convergence are
\cref{p1:prop:compactness}.  The no-escape argument in
\cref{p1:lem:scale-selection}, followed by strong \(L^3\)-compactness away from
\(t=0\), gives nontriviality through \cref{p1:prop:nonzero}.  The pointwise
Type I bound is \cref{p1:cor:classical-type-I-bound}, and the centered equation,
smoothness, boundedness, and local \(L^3\) nonvanishing are
\cref{p1:lem:renormalized-equation,p1:lem:smoothness,p1:lem:renorm-bounds}.
\end{proof}

\section{Liouville classes and the remaining class}\label{p1:sec:classes}

Fix a collection \(\calS\) of normalized Seregin ancient limits.  Complements
below are taken relative to \(\calS\).  Structural conditions involving the
unrenormalized ancient field refer to
\[
   U(x,t)=(-t)^{-1/2}V(x/\sqrt{-t},-\log(-t)).
\]

\begin{definition}[Small-amplitude class]\label{p1:def:small-class}
Let \(\varepsilon_0>0\) be the small-data Liouville constant for bounded
ancient solutions of \eqref{p1:eq:centered}.  We say
\(V\in\calC_{\mathrm{small}}\) if
\[
   \|V\|_{L^\infty(\R^3\times\R)}\le \varepsilon_0.
\]
\end{definition}

\begin{definition}[Stationary \(L^3\) class]\label{p1:def:stationary-class}
We say \(V\in\calC_{\mathrm{stat},L^3}\) if
\[
   V(y,\tau)=W(y)
   \quad\hbox{and}\quad
   W\in L^3(\R^3).
\]
\end{definition}

\begin{definition}[Uniformly \(L^3\)-tight class]\label{p1:def:tight-class}
We say \(V\in\calC_{L^3\mathrm{-tight}}\) if
\begin{equation}\label{p1:eq:L3-tightness}
   \lim_{R\to\infty}
   \sup_{\tau\in\R}\int_{|y|>R}|V(y,\tau)|^3\,dy=0 .
\end{equation}
\end{definition}

Every \(V\in\calS\) is bounded.  \Cref{p2:thm:tight-liouville} shows that, for such limits, the
tail condition \eqref{p1:eq:L3-tightness} gives the endpoint sequence-\(L^3\)
control needed after physical pullback; no separate
\(L^\infty_\tau L^3_y\) hypothesis is required for the tight-class exclusion.

\begin{definition}[Liouville-covered structure and decay class]
\label{p1:def:known-structure-decay-class}
A normalized Seregin limit \(V\) belongs to \(\calC_{\mathrm{known}}\) if its
physical ancient representative \(U\) satisfies at least one of the following
analytic hypotheses:
\begin{enumerate}[label=(\roman*)]
\item \(U\) is axisymmetric without swirl;
\item \(U\) is axisymmetric and satisfies the pointwise scale-invariant bound
\[
   |U(x,t)|\le \frac{C}{r}
\]
for some \(C<\infty\), where \(r\) is the cylindrical radius;
\item \(U\) is axisymmetric and
\[
   rU_\theta\in L_t^\infty L_x^p(\R^3\times(-\infty,0))
\]
for some \(1\le p<\infty\);
\item \(U\) is axisymmetric, periodic in the axial variable, and
\[
   rU_\theta\in L^\infty(\R^3\times(-\infty,0));
\]
\item the associated rescaled vorticity satisfies the perturbative
Gallay--Wayne weighted-vorticity hypotheses used in
\Cref{p3:thm:typeI-zero-hypotheses};
\end{enumerate}
\end{definition}

\begin{remark}
Only hypotheses for which a cited Liouville theorem gives \(U\equiv0\) are
included in \(\calC_{\mathrm{known}}\).  Broader qualitative descriptions are
not included unless they imply one of the listed hypotheses.
\end{remark}

\begin{lemma}[Structural quantities under centering]
\label{p1:lem:structure-centering}
The centered change of variables preserves axisymmetry and absence of swirl.
If \((r,z,t)\) and \((\rho,Z,\tau)\) are related by
\[
   \rho=\frac r{\sqrt{-t}},\qquad Z=\frac z{\sqrt{-t}},
   \qquad \tau=-\log(-t),
\]
then the swirl variables satisfy
\[
   \rho V_\theta(\rho,Z,\tau)=rU_\theta(r,z,t).
\]
\end{lemma}

\begin{proof}
The map \(x\mapsto y=x/\sqrt{-t}\) is radial in the cylindrical variables and
does not mix the \(e_r,e_\theta,e_z\) frame.  Since
\(V(y,\tau)=\sqrt{-t}\,U(x,t)\), the \(\theta\)-components obey
\(V_\theta=\sqrt{-t}\,U_\theta\), and the displayed identities follow.
\end{proof}

\begin{definition}[Remaining class]\label{p1:def:remaining-class}
The remaining class is
\[
   \calR(\calS)
   =
   \calS
   \setminus
   \left(
      \calC_{\mathrm{small}}
      \cup
      \calC_{\mathrm{stat},L^3}
      \cup
      \calC_{L^3\mathrm{-tight}}
      \cup
      \calC_{\mathrm{known}}
   \right).
\]
\end{definition}

\begin{remark}
The first four classes are not intended to be disjoint.  The remaining class is
the complement of their union, so overlaps among the Liouville classes cause no
ambiguity.
\end{remark}

\begin{proposition}[Exhaustion by defined classes]\label{p1:prop:exhaustion}
Every \(V\in\calS\) belongs either to one of the four Liouville classes or to
\(\calR(\calS)\).
\end{proposition}

\begin{proof}
Let
\[
   \calU=
      \calC_{\mathrm{small}}
      \cup
      \calC_{\mathrm{stat},L^3}
      \cup
      \calC_{L^3\mathrm{-tight}}
      \cup
      \calC_{\mathrm{known}}.
\]
For \(V\in\calS\), either \(V\in\calU\), in which case it lies in at least one
Liouville class, or \(V\in\calS\setminus\calU=\calR(\calS)\).
\end{proof}

\section{Liouville hypotheses and class exclusions}\label{p1:sec:liouville}

\begin{lemma}[Centered Stokes kernel estimates]
\label{p1:lem:centered-stokes-estimates}
Let
\[
   L=\Delta-\frac12y\cdot\nabla-\frac12,
   \qquad S(s)=e^{sL}.
\]
There is a universal constant \(C_0\) such that, for all \(s>0\),
\[
   \|S(s)f\|_{L^\infty}\le e^{-s/2}\|f\|_{L^\infty}
\]
and
\[
   \|S(s)\mathbb P\nabla\cdot F\|_{L^\infty}
   \le C_0 e^{-s/2}(1-e^{-s})^{-1/2}\|F\|_{L^\infty}.
\]
\end{lemma}

\begin{proof}
The first estimate follows from the explicit Ornstein--Uhlenbeck kernel:
\[
   S(s)f(y)=e^{-s/2}
   \int_{\R^3}G_{1-e^{-s}}(e^{-s/2}y-z)f(z)\,dz,
\]
where \(G_a\) is the heat kernel with variance \(a\).  The kernel has total
mass one.  For the projected-divergence estimate, write the corresponding
kernel as the Oseen projected-gradient kernel after the same change of
variables.  More explicitly, the semigroup for
\(\Delta-\frac12y\cdot\nabla\) is
\[
   (e^{s(\Delta-\frac12y\cdot\nabla)}g)(y)
   =
   (e^{(1-e^{-s})\Delta}g)(e^{-s/2}y).
\]
Applying this to the distribution \(g=\mathbb P\nabla\cdot F\) and using the
projected-gradient Oseen kernel estimate
\[
   \|e^{\theta\Delta}\mathbb P\nabla\cdot F\|_\infty
   \le C\theta^{-1/2}\|F\|_\infty,\qquad \theta>0,
\]
gives the displayed estimate with \(\theta=1-e^{-s}\), together with the
extra factor \(e^{-s/2}\) from the zeroth-order term in \(L\).  The
projected-gradient Oseen estimate follows from the \(L^1\)-bound
\[
   \|\nabla O(\cdot,\theta)\|_{L^1}\le C\theta^{-1/2}
\]
for the gradient of the Oseen kernel \(O(\cdot,\theta)\).  This is the
bounded-data Stokes estimate used by Kato \cite{Kato1984}; it does not require
boundedness of the Leray projector itself on \(L^\infty\).
\end{proof}

\begin{lemma}[Mild representation for normalized limits]
\label{p1:lem:mild-representation}
Every normalized Seregin ancient limit satisfies, for all \(\Theta>0\),
\[
   V(\tau)
   =
   S(\Theta)V(\tau-\Theta)
   -
   \int_0^\Theta
   S(\Theta-\sigma)\mathbb P\nabla\cdot
   (V\otimes V)(\tau-\Theta+\sigma)\,d\sigma .
\]
\end{lemma}

\begin{proof}
Normalized Seregin limits are smooth and bounded by
\cref{p1:lem:smoothness,p1:lem:renorm-bounds}.  On a finite interval
\([\tau-\Theta,\tau]\), the tensor \(V\otimes V\) is bounded, so
\cref{p1:lem:centered-stokes-estimates} makes the Duhamel integral finite in
\(L^\infty\).  In the sense of solenoidal distributions, applying the Leray
projector to \eqref{p1:eq:centered} removes the pressure and gives
\[
   \partial_\tau V=LV-\mathbb P\nabla\cdot(V\otimes V)
\]
on every finite interval.  Since \(V\) is smooth, the map
\[
   s\mapsto S(\tau-s)V(s)
\]
is differentiable as a tempered distribution for \(\tau-\Theta<s<\tau\), and
\[
   \frac{d}{ds}\bigl(S(\tau-s)V(s)\bigr)
   =
   -S(\tau-s)\mathbb P\nabla\cdot(V\otimes V)(s).
\]
Integrating this identity from \(\tau-\Theta\) to \(\tau\) gives the stated
Duhamel formula.
\end{proof}

\begin{lemma}[Small bounded ancient Liouville theorem]\label{p1:lem:small-liouville}
There is \(\varepsilon_0>0\) such that every bounded mild ancient solution of
\eqref{p1:eq:centered} with
\[
   \|V\|_{L^\infty(\R^3\times\R)}\le\varepsilon_0
\]
is identically zero.
\end{lemma}

\begin{proof}
Write the projected centered equation as
\[
   \partial_\tau V=LV-\mathbb P\nabla\cdot(V\otimes V),
   \qquad
   L=\Delta-\frac12 y\cdot\nabla-\frac12.
\]
Let \(S(s)=e^{sL}\).  By \cref{p1:lem:centered-stokes-estimates}, the nonlinear
kernel is integrable on \((0,\infty)\); set
\[
   A=C_0\int_0^\infty e^{-s/2}(1-e^{-s})^{-1/2}\,ds<\infty .
\]
For fixed \(\tau\) and \(T>0\), the mild identity on \([\tau-T,\tau]\) gives
\[
   V(\tau)
   =
   S(T)V(\tau-T)
   -
   \int_0^T
   S(T-\sigma)\mathbb P\nabla\cdot(V\otimes V)(\tau-T+\sigma)\,d\sigma .
\]
If \(M=\|V\|_{L^\infty(\R^3\times\R)}\), then
\[
   \|V(\tau)\|_\infty\le e^{-T/2}M+AM^2.
\]
Letting \(T\to\infty\) and taking the supremum in \(\tau\) gives
\[
   M\le AM^2.
\]
If \(0<M<1/A\), this is impossible.  Choose \(\varepsilon_0<1/A\); then
\(M\le\varepsilon_0\) implies \(M=0\).
\end{proof}

\begin{lemma}[Stationary centered profiles]\label{p1:lem:stationary-profile}
If \(V(y,\tau)=W(y)\) is a smooth solution of \eqref{p1:eq:centered}, then after
subtracting a function of \(\tau\) from the pressure, \(W\) solves
\[
   (W\cdot\nabla)W+\nabla P
   =
   \Delta W-\frac12(W+y\cdot\nabla W),
   \qquad \divergence W=0.
\]
\end{lemma}

\begin{proof}
Since \(V\) is independent of \(\tau\), \(\partial_\tau V=0\).  Set
\[
   F(y)=
   \Delta W-\frac12(W+y\cdot\nabla W)-(W\cdot\nabla)W .
\]
The centered equation gives \(\nabla_y\Pi(\cdot,\tau)=F\) for a.e. \(\tau\).
Fix one such time \(\tau_0\) and put \(P(y)=\Pi(y,\tau_0)\).  Then
\(\nabla_y(\Pi(\cdot,\tau)-P)=0\) for a.e. \(\tau\), so
\(\Pi(y,\tau)-P(y)=b(\tau)\) is spatially constant.  Subtracting this gauge
from \(\Pi\) leaves the equation unchanged and yields the displayed
stationary self-similar equation for \(W\).
\end{proof}

\begin{theorem}[Small and stationary class exclusions]
\label{p1:thm:classical-classes}
Let \(V\) be a normalized Seregin ancient limit.  Then
\[
   V\notin \calC_{\mathrm{small}},
   \qquad
   V\notin \calC_{\mathrm{stat},L^3}.
\]
\end{theorem}

\begin{proof}
If \(V\in\calC_{\mathrm{small}}\), then
\cref{p1:lem:mild-representation,p1:lem:small-liouville} give \(V\equiv0\),
contradicting the local \(L^3\) nonvanishing condition
\eqref{p1:eq:nonvanishing}.

If \(V\in\calC_{\mathrm{stat},L^3}\), then
\(V(y,\tau)=W(y)\) with \(W\in L^3(\R^3)\).  By
\cref{p1:lem:stationary-profile}, \(W\) satisfies the stationary self-similar
profile equation after a pressure gauge change.
The theorem of Ne\v{c}as--R\r{u}\v{z}i\v{c}ka--\v{S}ver\'ak
\cite{NRS1996} gives \(W\equiv0\).  Again
\eqref{p1:eq:nonvanishing} gives the contradiction.
\end{proof}

\begin{theorem}[Uniformly tight \(L^3\) closure]
\label{p1:hyp:tight-liouville}
No nonzero normalized Seregin ancient limit belongs to
\(\calC_{L^3\mathrm{-tight}}\).
\end{theorem}

\begin{proof}
This is the tight-class Liouville theorem proved in
\Cref{p2:thm:tight-liouville,p2:cor:tight-input-for-paper-I}.
\Cref{p2:thm:tight-liouville} proves that, for bounded centered ancient solutions,
the tail-tightness condition gives the endpoint sequence-\(L^3\) control after
physical pullback.  It then applies the Albritton--Barker endpoint ancient
Liouville theorem and contradicts the compact velocity normalization
\Cref{p1:prop:nonzero}.
\end{proof}

\begin{theorem}[Liouville-covered structure and decay exclusions]
\label{p1:thm:known-structure-decay-liouville}
No nonzero normalized Seregin ancient limit belongs to
\(\calC_{\mathrm{known}}\).
\end{theorem}

\begin{proof}
By \cref{p1:def:known-structure-decay-class}, membership in
\(\calC_{\mathrm{known}}\) means that the physical ancient representative
\(U\) satisfies one of the listed analytic hypotheses.  The axisymmetric
no-swirl and pointwise \(1/r\) cases are excluded by
Koch--Nadirashvili--Seregin--\v{S}ver\'ak \cite{KNSS2009}.  The finite
\(L^p\) circulation case is excluded by Lei--Zhang--Zhao \cite{LZZ2017}, and
the axial-periodic bounded-circulation case is excluded by Lei--Ren--Zhang
\cite{LRZ2022}.  The perturbative Gallay--Wayne weighted-vorticity hypotheses
are excluded by the corresponding Liouville theorem in
\Cref{p3:thm:typeI-zero-hypotheses}.  Each
alternative gives \(V\equiv0\), contradicting the local \(L^3\) nonvanishing
condition \eqref{p1:eq:nonvanishing}.  Hence no nonzero normalized Seregin
ancient limit lies in \(\calC_{\mathrm{known}}\).
\end{proof}

\begin{proposition}[Exclusion of the Liouville classes]
\label{p1:prop:known-class-exclusion}
No normalized Seregin ancient limit can belong to
\[
   \calC_{\mathrm{small}}
   \cup
   \calC_{\mathrm{stat},L^3}
   \cup
   \calC_{L^3\mathrm{-tight}}
   \cup
   \calC_{\mathrm{known}}.
\]
\end{proposition}

\begin{proof}
The small and stationary \(L^3\) classes are excluded by
\cref{p1:thm:classical-classes}.  The uniformly tight class is excluded by
\cref{p1:hyp:tight-liouville}.  The structure/decay class is excluded by
\cref{p1:thm:known-structure-decay-liouville}.
\end{proof}

\section{The remaining-case criterion}\label{p1:sec:criterion}

\begin{hypothesis}[Residual-class hypothesis]\label{p1:hyp:no-remainder}
For every collection \(\calS\) of normalized Seregin ancient limits generated
by admissible finite-energy Type I singularities,
\[
   \calR(\calS)=\varnothing .
\]
\end{hypothesis}

\begin{corollary}[Rigidity in the remaining case]\label{p1:hyp:remainder-rigidity}
Assume \Cref{p1:hyp:no-remainder}.
For every collection \(\calS\) of normalized Seregin ancient limits generated
by admissible finite-energy Type I singularities, every element of
\(\calR(\calS)\) is identically zero.
\end{corollary}

\begin{proof}
This is immediate from \cref{p1:hyp:no-remainder}.
\end{proof}

\begin{proposition}[Equivalence on normalized Seregin limits]
\label{p1:prop:remainder-equivalence}
For collections of normalized Seregin limits, the residual-class hypothesis
\cref{p1:hyp:no-remainder} implies \cref{p1:hyp:remainder-rigidity}.  Conversely,
the corresponding rigidity assertion implies emptiness of the remaining class.
\end{proposition}

\begin{proof}
If \cref{p1:hyp:no-remainder} holds, the rigidity statement is vacuous.
Conversely, assume \cref{p1:hyp:remainder-rigidity}.  If some normalized
collection had \(V\in\calR(\calS)\), the rigidity hypothesis would give
\(V\equiv0\), directly contradicting the local \(L^3\) nonvanishing condition
\eqref{p1:eq:nonvanishing}.  Hence the remaining class is empty.
\end{proof}

\begin{proof}[Proof of \cref{p1:thm:remaining-criterion}]
Assume \Cref{p1:hyp:no-remainder}.
Suppose that \((u,p,z_*)\) is an admissible finite-energy Type I singularity.
By \cref{p1:thm:extraction}, there exists a normalized nonzero Seregin ancient
limit \(V\).  Let \(\calS\) be any collection of normalized limits containing
\(V\).

By \cref{p1:prop:exhaustion}, either \(V\) belongs to one of the Liouville classes
or \(V\in\calR(\calS)\).  The first alternative is impossible by
\cref{p1:prop:known-class-exclusion}.  The second alternative is impossible by
\cref{p1:hyp:no-remainder}.  Hence the assumed admissible Type I singularity
cannot exist.
\end{proof}

The preceding theorem is conditional on the residual-class hypothesis for the
stated admissible finite-energy class.  The compactness construction, pressure
normalization, and class exhaustion are contained in this paper, with the local
velocity concentration result stated explicitly in
\cref{p1:prop:local-velocity-concentration}.  The structure/decay class is
restricted to the assumptions in \cref{p1:def:known-structure-decay-class}; any
other structured behavior lies in the remaining class and is addressed only by
\cref{p1:hyp:no-remainder}.  No whole-space critical-norm estimate on the
original solution is used in the conditional closure step.

\section{Elementary symmetries of the centered equation}
\label{p1:app:centered-symmetries}

\begin{lemma}[Logarithmic-time translation]\label{p1:lem:tau-translation}
If \(V\) solves \eqref{p1:eq:centered}, then
\[
   V_{\tau_0}(y,\tau)=V(y,\tau+\tau_0)
\]
also solves \eqref{p1:eq:centered}.  If
\[
   U^\lambda(x,t)=\lambda U(\lambda x,\lambda^2t),
\]
then the renormalized field associated with \(U^\lambda\) is
\[
   V^\lambda(y,\tau)=V(y,\tau-2\log\lambda).
\]
\end{lemma}

\begin{proof}
The first statement follows because \eqref{p1:eq:centered} is autonomous in
\(\tau\).  For the second, put \(s=-t=e^{-\tau}\).  Then
\[
   \sqrt{s}\,U^\lambda(y\sqrt{s},-s)
   =
   \lambda\sqrt{s}\,U(\lambda y\sqrt{s},-\lambda^2s)
   =
   V(y,-\log(\lambda^2s))
   =
   V(y,\tau-2\log\lambda).
\]
\end{proof}

\begin{proposition}[Dilation and translation are not symmetries]
\label{p1:prop:no-cascade}
The centered equation \eqref{p1:eq:centered} is autonomous in \(\tau\), but it is
not invariant under
\[
   V(y,\tau)\mapsto \alpha V(\alpha y,\alpha^2\tau)
\]
unless \(\alpha=1\).  It is also not invariant under constant translations
\[
   V(y,\tau)\mapsto V(y-a,\tau)
\]
unless \(a=0\).
\end{proposition}

\begin{proof}
For scaling, set \(Y=\alpha y\), \(T=\alpha^2\tau\), and
\(\widetilde V(y,\tau)=\alpha V(Y,T)\),
\(\widetilde\Pi(y,\tau)=\alpha^2\Pi(Y,T)\).  The non-drift part scales with
factor \(\alpha^3\), while
\[
   -\frac12(\widetilde V+y\cdot\nabla\widetilde V)
   =
   -\frac{\alpha}{2}(V+Y\cdot\nabla_YV)(Y,T).
\]
Using the equation for \(V\), the remaining term in the equation for \(\widetilde V\) is
\[
   \frac{\alpha-\alpha^3}{2}(V+Y\cdot\nabla_YV)(Y,T).
\]
If \(\alpha\ne1\), this forces \(V+Y\cdot\nabla V=0\).  Along each ray,
\[
   \frac{d}{dr}\bigl(rV(r\omega,T)\bigr)=0.
\]
Boundedness and smoothness at the origin force the constant to be zero, hence
\(V\equiv0\).  Thus a nonzero centered solution has no such scaling symmetry.

For translations, all translation-invariant terms transform by composition,
but the drift satisfies
\[
   y\cdot\nabla V(y-a,\tau)
   =
   (y-a)\cdot\nabla V(y-a,\tau)+a\cdot\nabla V(y-a,\tau).
\]
The extra term is absent from the original equation.  Thus spatial
translation is not a symmetry of the fixed centered equation unless \(a=0\),
apart from exceptional fixed solutions satisfying \(a\cdot\nabla V\equiv0\).
\end{proof}



\clearpage
\section*{Part III. Uniformly Tight Ancient Dynamics}

\subsection*{Overview}

We study bounded ancient solutions of the three-dimensional Navier--Stokes
equations in backward self-similar variables under uniform \(L^3\)-tightness.
We first establish the analytic compactness framework: bounded mild ancient
solutions are smooth, locally compact under time translation, and the
common \(L^\infty_\tau L^3_y\) and tightness bounds used in the compactness
statements are stable under locally smooth limits.  We then prove a
minimal-element reduction for closed normalized uniformly tight collections:
any nonempty such collection contains a nonzero solution of minimal
\(L^\infty(\R^3\times\R)\) size.

For the time-translation hull of a bounded ancient solution we construct compact
invariant hulls and invariant probability measures by Krylov--Bogolyubov
averaging.  These measures satisfy a stationary statistical form of the
self-similar Navier--Stokes equation.  Passing to the barycenter produces an
explicit quadratic covariance term, isolating the obstruction to obtaining a
single stationary self-similar profile.

Finally, we close the uniformly \(L^3\)-tight class directly from its defining
hypotheses.  The tail condition, together with the boundedness of the ancient
solution, gives the
sequence-\(L^3\) bound needed after the physical pullback.  Since this pullback
preserves the critical \(L^3\)-norm on each time slice, the endpoint ancient
Liouville theorem of
Albritton--Barker then forces the pullback, and hence the centered solution, to
vanish.  Therefore the uniformly \(L^3\)-tight class of the Type I
ancient-limit reduction is excluded without any compact-hull Lyapunov
hypothesis.

\section{Introduction and main results}\label{p2:sec:introduction}

\subsection{The uniformly tight ancient class}

We consider the three-dimensional incompressible Navier--Stokes equations in
backward self-similar variables.  The centered velocity \(V\) and pressure \(P\)
solve
\begin{equation}\label{p2:eq:centered}
   \partial_\tau V+(V\cdot\nabla)V+\nabla P
   =
   \Delta V-\frac12(V+y\cdot\nabla V),
   \qquad \divergence V=0
\end{equation}
on \(\R^3\times\R\).  In \Cref{p1:thm:extraction}, admissible Type I
singularities are reduced to nonzero normalized ancient solutions of
\eqref{p2:eq:centered}.  We then study the uniformly \(L^3\)-tight class, meaning
the class in which
\begin{equation}\label{p2:eq:L3-tightness}
   \lim_{R\to\infty}
   \sup_{\tau\in\R}\int_{|y|>R}|V(y,\tau)|^3\,dy=0 .
\end{equation}
Several compactness statements below also record the common bound
\begin{equation}\label{p2:eq:L3-bound}
   \sup_{\tau\in\R}\|V(\tau)\|_{L^3(\R^3)}<\infty ,
\end{equation}
but the endpoint exclusion of the tight class uses only boundedness and the
tail condition \eqref{p2:eq:L3-tightness}.

This condition on the centered limit is distinct from the original-solution condition
\[
   u\in L^\infty(0,T;L^3(\R^3)).
\]
For finite-energy suitable weak solutions, the original-solution assumption
\(u\in L^\infty_tL^3_x\) near a putative singular time is incompatible with
singularity by the theorem of Escauriaza--Seregin--\v{S}ver\'ak
\cite{ESS2003}.  Uniform \(L^3\)-tightness in this paper is a condition on
centered ancient solutions after blow-up.

\subsection{Minimal elements and compact hulls}

The first reduction is variational:
\[
   \hbox{nonempty tight normalized class}
   \quad\Longrightarrow\quad
   \hbox{\(L^\infty\)-minimal ancient element}.
\]
The normalization is a fixed local velocity nonvanishing condition.  It
prevents minimizing sequences from losing all mass by time translation.
Once a minimal element \(V_*\) is selected, its time-translation hull is
\[
   \calH(V_*)
   =
   \overline{\{\Theta_sV_*:\ s\in\R\}}^{C^\infty_{\loc}},
   \qquad
   (\Theta_sV)(y,\tau)=V(y,\tau+s).
\]
The hull is compact because bounded mild ancient solutions are locally smooth
with estimates depending only on their common \(L^\infty\) bound.

\subsection{Invariant measures and covariance}

Every compact invariant hull \(K\) carries invariant probability measures.
Indeed, Krylov--Bogolyubov averages
\[
   \mu_T=\frac1T\int_0^T\delta_{\Theta_sW_0}\,ds
\]
have weak subsequential limits.  Such invariant measures satisfy a stationary
statistical form of \eqref{p2:eq:centered}.  If
\[
   \overline W(y)=\int_K W(y,0)\,d\mu(W),
\]
then the averaged nonlinearity decomposes as
\[
   \int_K W_iW_j\,d\mu
   =
   \overline W_i\,\overline W_j+C_{ij}.
\]
Thus the barycenter solves the stationary self-similar equation only up to the
explicit covariance term \(C\).  The invariant-measure argument alone does
not prove that a recurrent compact hull contains a single stationary profile.

\subsection{\texorpdfstring{Endpoint \(L^3\) closure of the tight class}{Endpoint L3 closure of the tight class}}

The compact-hull and invariant-measure constructions identify the covariance
obstruction in a purely local topology, but they are not needed to close the
uniformly \(L^3\)-tight class.  The tail condition already gives the
sequence-\(L^3\) information required by the endpoint ancient Liouville theorem,
because the centered solution is bounded.  Under the physical change of variables
\[
   u(x,t)=(-t)^{-1/2}
   V\left(\frac{x}{\sqrt{-t}},-\log(-t)\right),
   \qquad t<0,
\]
the \(L^3\)-norm is invariant:
\[
   \|u(t)\|_{L^3_x}
   =
   \|V(-\log(-t))\|_{L^3_y}.
\]
Consequently every tail-tight centered ancient solution produces
a physical mild ancient solution with bounded \(L^3\)-norm along a
sequence \(t_k\downarrow-\infty\).  The endpoint theorem of
Albritton--Barker \cite{AlbrittonBarker2019} then gives \(u\equiv0\), and the
change of variables gives \(V\equiv0\).

This is a class-local argument.  It uses only the boundedness and tail
tightness already present in the ancient solution class and introduces no
additional estimate outside that class, recurrence object, or Lyapunov
functional.

The logical status of the paper is summarized as follows.

\begin{center}
\small
\begin{tabular}{p{0.58\linewidth}p{0.25\linewidth}}
\hline
Step & Status in this paper\\
\hline
bounded ancient smoothing & proved\\
local smooth compactness & proved\\
closure of uniform \(L^3\)-tightness & proved\\
minimal \(L^\infty\) element & proved\\
compact trajectory hull & proved\\
invariant measures & proved\\
stationary statistical identity & proved\\
covariance obstruction & identified\\
physical pullback and critical \(L^3\) invariance & proved\\
endpoint \(L^3\) ancient Liouville closure & cited and applied\\
tight-class consequence for the Type I reduction & proved corollary\\
\hline
\end{tabular}
\end{center}

Thus the uniformly \(L^3\)-tight class is closed by its own hypotheses.  The
compact-hull material remains useful for describing possible recurrent
dynamics and the covariance obstruction, but no compact-recurrence or
Lyapunov assumption enters the tight-class exclusion proved here.

\subsection{Main results}

The first headline theorem is the analytic compactness and closure statement.

\begin{theorem}[Compactness and closure of uniformly tight ancient solutions]
\label{p2:thm:tight-compactness-intro}
Let \(\{V_n\}\) be a sequence of bounded mild ancient solutions of
\eqref{p2:eq:centered}.  Assume
\[
   \sup_n\|V_n\|_{L^\infty(\R^3\times\R)}<\infty .
\]
Then a subsequence converges locally smoothly to a bounded mild ancient
solution \(V\).

If in addition there are \(L<\infty\) and a modulus
\(\omega:[1,\infty)\to[0,\infty)\), \(\omega(R)\to0\), such that
\[
   \sup_n\sup_{\tau\in\R}\|V_n(\tau)\|_{L^3(\R^3)}\le L
\]
and
\[
   \sup_n\sup_{\tau\in\R}
   \int_{|y|>R}|V_n(y,\tau)|^3\,dy\le \omega(R)
   \qquad (R\ge1),
\]
then
\[
   \sup_{\tau\in\R}\|V(\tau)\|_{L^3(\R^3)}\le L,
   \qquad
   \sup_{\tau\in\R}\int_{|y|>R}|V(y,\tau)|^3\,dy\le\omega(R),
\]
and in particular \(V\) is uniformly \(L^3\)-tight.
\end{theorem}

Compatible pressure representatives also converge locally weakly modulo
time-dependent gauges under the additional \(L^\infty_\tau L^3_y\) assumption;
see \Cref{p2:lem:pressure-reconstruction}.

The second headline theorem is the minimal-element reduction.

\begin{theorem}[Minimal element in a normalized uniformly tight class]
\label{p2:thm:minimal-tight-intro}
Let \(A\) be a nonempty collection of bounded mild ancient solutions satisfying
a common \(L^\infty\) bound, a common \(L^\infty_\tau L^3_y\) bound, a common
uniform \(L^3\)-tightness modulus, and a fixed local nonvanishing
normalization
\[
   \sup_{\tau\in\R}\int_{B_{R_0}}|V(y,\tau)|^3\,dy\ge \eta_0>0.
\]
Assume \(A\) is invariant under time translation and closed under locally
smooth limits which retain the normalization.  Then \(A\) contains an element
\(V_*\) minimizing
\[
   \|V\|_{L^\infty(\R^3\times\R)}
\]
among all \(V\in A\).  In particular \(V_*\not\equiv0\).
\end{theorem}

The third headline theorem is the endpoint \(L^3\) tight-class Liouville
criterion.

\begin{theorem}[Uniformly tight Liouville criterion]
\label{p2:thm:tight-liouville-intro}
Let \(A\) be a collection of bounded mild ancient solutions of
\eqref{p2:eq:centered}.  Assume there is a modulus
\(\omega:[1,\infty)\to[0,\infty)\), \(\omega(R)\to0\), such that
\[
   \sup_{\tau\in\R}
   \int_{|y|>R}|V(y,\tau)|^3\,dy\le\omega(R)
   \qquad\text{for every }V\in A,\ R\ge1,
\]
and assume that \(A\) has a fixed local nonvanishing normalization:
\[
   \sup_{\tau\in\R}\int_{B_{R_0}}|V(y,\tau)|^3\,dy\ge\eta_0>0
   \qquad\text{for every }V\in A.
\]
Then
\[
   A=\varnothing .
\]
\end{theorem}

The proof of the last theorem follows the dependency diagram
\[
\begin{array}{c}
\text{uniformly \(L^3\)-tight locally nonzero class }A\\[2mm]
\Downarrow\quad\text{boundedness plus tail tightness}\\[2mm]
\sup_k\|V(\tau_k)\|_{L^3(\R^3)}<\infty
\quad\text{for a sequence }\tau_k\to-\infty\\[2mm]
\Downarrow\quad\text{physical pullback and critical \(L^3\)-invariance}\\[2mm]
\text{mild ancient Navier--Stokes solution }u\\[1mm]
\text{with }\sup_k\|u(t_k)\|_{L^3(\R^3)}<\infty
\quad\text{for some }t_k\downarrow-\infty\\[2mm]
\Downarrow\quad\text{Albritton--Barker endpoint theorem}\\[2mm]
u=0\\[2mm]
\Downarrow\\[2mm]
\text{contradiction with local velocity nonvanishing in }A .
\end{array}
\]

\subsection{Relation to the Type I reduction}

\Cref{p1:thm:extraction} reduces admissible Type I singularities
to normalized Seregin ancient limits and organizes such limits into small,
stationary \(L^3\), uniformly tight, structured/decay, and residual classes.
Here we prove the tight-class consequence without an additional
compact-dynamical hypothesis: the uniformly \(L^3\)-tight class contains no
nonzero normalized ancient limit.

\subsection{Organization of the paper}

\Cref{p2:sec:centered-topology} fixes the centered equation, mild formulation,
time translation, and local smooth topology.  \Cref{p2:sec:smoothing-compactness}
proves smoothing, pressure reconstruction, stability, and compactness.
\Cref{p2:sec:tightness-rigidity} proves closure of uniform \(L^3\)-tightness and
recalls stationary \(L^3\) rigidity.  \Cref{p2:sec:minimal-elements} proves the
minimal-element reduction.  \Cref{p2:sec:compact-hulls} constructs compact
trajectory hulls and minimal invariant subsets.  \Cref{p2:sec:measures-covariance}
constructs invariant measures and identifies the covariance obstruction.
\Cref{p2:sec:endpoint-L3-closure} proves the physical pullback identities,
applies the endpoint \(L^3\) ancient Liouville theorem, and states the
consequence used in the Type I reduction.

\section{Centered ancient solutions and the local smooth topology}
\label{p2:sec:centered-topology}

\subsection{The centered equation}

Throughout the paper \(V\) denotes a centered ancient velocity and \(P\) its
pressure.  An element of a compact hull is denoted by \(W\), and a
stationary profile by \(w(y)\).  Define
\[
   L V=\Delta V-\frac12 y\cdot\nabla V-\frac12 V .
\]
After applying the Helmholtz--Leray projection \(\mathbb P\), the centered
equation becomes
\begin{equation}\label{p2:eq:projected}
   \partial_\tau V
   =
   LV-\mathbb P\divergence(V\otimes V).
\end{equation}

\subsection{The Ornstein--Uhlenbeck--Stokes semigroup}

Let \(S(a)=e^{aL}\), \(a>0\).  It is given by
\begin{equation}\label{p2:eq:OU-semigroup}
   (S(a)f)(y)
   =
   e^{-a/2}
   \int_{\R^3}
   \frac{\exp\!\left(-\frac{|z-e^{-a/2}y|^2}{4(1-e^{-a})}\right)}
        {[4\pi(1-e^{-a})]^{3/2}}
   f(z)\,dz .
\end{equation}
The factor \(e^{-a/2}\) represents the damping caused by the term \(-V/2\).

\begin{lemma}[Centered Stokes kernel estimate]\label{p2:lem:centered-stokes-kernel}
For \(a>0\) and \(\phi\in C_c^\infty(\R^3;\R^3)\),
\begin{equation}\label{p2:eq:stokes-adjoint-estimate}
   \|(S(a)\mathbb P\divergence)^*\phi\|_{L^1(\R^3)}
   \le
   C e^{-a/2}(1-e^{-a})^{-1/2}\|\phi\|_{L^1(\R^3)}.
\end{equation}
Consequently the Duhamel term in \eqref{p2:eq:mild} is well-defined by duality
for bounded \(V\).
\end{lemma}

\begin{proof}
Set \(\lambda=e^{-a/2}\) and \(\theta=1-e^{-a}\).  From
\eqref{p2:eq:OU-semigroup},
\[
   S(a)f(y)=e^{-a/2}(e^{\theta\Delta}f)(\lambda y).
\]
Thus
\[
   S(a)^*\phi(z)
   =
   e^{-a/2}\lambda^{-3}
   (e^{\theta\Delta}\psi)(z),
   \qquad
   \psi(x)=\phi(x/\lambda),
\]
and \(\|\psi\|_{L^1}=\lambda^3\|\phi\|_{L^1}\).  The Leray projection is a
homogeneous Fourier multiplier of degree zero, so it commutes with this
dilation.  The classical heat--Stokes kernel estimate
\[
   \|(e^{\theta\Delta}\mathbb P\divergence)^*\phi\|_{L^1}
   \le C\theta^{-1/2}\|\phi\|_{L^1},
   \qquad \theta>0,
\]
follows from the Oseen kernel bound; see
\cite{FabesJonesRiviere1972,GigaMiyakawa1985,Sohr2001}.  Substituting
\(\theta=1-e^{-a}\) and using
\[
   e^{-a/2}\lambda^{-3}\|\psi\|_{L^1}=e^{-a/2}\|\phi\|_{L^1}
\]
gives \eqref{p2:eq:stokes-adjoint-estimate}.  Since
\((1-e^{-a})^{-1/2}\) is integrable at \(a=0\), the duality pairing in
\eqref{p2:eq:mild} is integrable on finite time intervals for
\(V\in L^\infty\).
\end{proof}

\subsection{Riesz transform convention}

We use the convention that \(R_i\) has Fourier symbol \(i\xi_i/|\xi|\).  Thus
\(R_iR_j\) has symbol \(-\xi_i\xi_j/|\xi|^2\), and
\[
   P=R_iR_j(F_{ij})
   \quad\Longleftrightarrow\quad
   -\Delta P=\partial_i\partial_jF_{ij}.
\]
Pressures are determined only modulo functions of time unless a whole-space
Leray gauge is explicitly chosen.

\subsection{Bounded mild ancient solutions}

\begin{definition}[Bounded mild ancient solution]\label{p2:def:mild}
A vector field \(V:\R^3\times\R\to\R^3\) is a bounded mild ancient solution of
\eqref{p2:eq:centered} if
\[
   V\in L^\infty(\R^3\times\R),\qquad \divergence V=0
\]
in distributions, and for every \(\sigma<\tau\),
\begin{equation}\label{p2:eq:mild}
   V(\tau)
   =
   S(\tau-\sigma)V(\sigma)
   -
   \int_\sigma^\tau
   S(\tau-s)\mathbb P\divergence(V\otimes V)(s)\,ds
\end{equation}
in the weak-* topology of \(L^\infty(\R^3)\).
\end{definition}

The Duhamel term in \eqref{p2:eq:mild} is interpreted by duality.  For
\(\phi\in C_c^\infty(\R^3;\R^3)\),
\[
   \left\langle
      S(\tau-s)\mathbb P\divergence(V\otimes V)(s),\phi
   \right\rangle
   =
   \left\langle
      V\otimes V(s),
      (S(\tau-s)\mathbb P\divergence)^*\phi
   \right\rangle ,
\]
and \Cref{p2:lem:centered-stokes-kernel} makes the scalar integrand absolutely
integrable on \(s\in(\sigma,\tau)\).

\subsection{Time translation}

For \(s\in\R\) define
\[
   (\Theta_sV)(y,\tau)=V(y,\tau+s).
\]
The equation, the mild formulation, the \(L^\infty\) norm, and the
\(L^\infty_\tau L^3_y\) and tail bounds used below are invariant under
\(\Theta_s\).  Indeed, replacing \((\sigma,\tau)\) by
\((\sigma+s,\tau+s)\) in \eqref{p2:eq:mild} and changing variables in the time
integral gives the mild formula for \(\Theta_sV\).

\subsection{The locally smooth topology}

\begin{definition}[Local smooth topology]\label{p2:def:local-smooth-topology}
The locally smooth topology on \(C^\infty_{\loc}(\R^3\times\R;\R^3)\) is
generated by the seminorms
\[
   \|V\|_{m,N}
   =
   \max_{|\alpha|+j\le m}
   \sup_{\substack{|y|\le N\\ |\tau|\le N}}
   |\partial_y^\alpha\partial_\tau^jV(y,\tau)|,
   \qquad m,N\in\N .
\]
One compatible metric is
\[
   d(U,V)
   =
   \sum_{m,N=1}^{\infty}
   2^{-(m+N)}
   \frac{\|U-V\|_{m,N}}{1+\|U-V\|_{m,N}}.
\]
\end{definition}

Thus \(V_n\to V\) locally smoothly if and only if the convergence is in
\(C^m(K)\) for every compact cylinder
\(K\Subset\R^3\times\R\) and every \(m\ge0\).

\section{Smoothing, pressure reconstruction, and compactness}
\label{p2:sec:smoothing-compactness}

\subsection{Local smoothing}

\begin{lemma}[Smoothing of bounded mild solutions]\label{p2:lem:mild-smoothing}
Every bounded mild ancient solution is smooth on \(\R^3\times\R\) and satisfies
the projected equation \eqref{p2:eq:projected} classically.  Moreover, for each
compact cylinder \(K\Subset\R^3\times\R\) and each integer \(m\ge0\),
\[
   \|V\|_{C^m(K)}
   \le
   C\bigl(m,K,\|V\|_{L^\infty(\R^3\times\R)}\bigr).
\]
The same estimate holds uniformly for any family with a common
\(L^\infty\) bound.
\end{lemma}

\begin{proof}
Fix a compact time interval \(J\Subset\R\) containing the time projection of
\(K\), and set \(t=-e^{-\tau}\).  On the corresponding compact interval
\(t(J)\Subset(-\infty,0)\), define
\[
   U(x,t)=(-t)^{-1/2}V(x/\sqrt{-t},-\log(-t)).
\]
The map \((x,t)\mapsto (x/\sqrt{-t},-\log(-t))\) is a smooth diffeomorphism
between \(\R^3\times t(J)\) and \(\R^3\times J\), with all derivatives bounded
on compact subcylinders.

The centered mild formula \eqref{p2:eq:mild} transforms, under this parabolic
change of variables, into the usual projected Navier--Stokes mild formula on
every compact subinterval of \(t(J)\):
\[
   U(t)=e^{(t-s)\Delta}U(s)
   -
   \int_s^t e^{(t-r)\Delta}\mathbb P\divergence(U\otimes U)(r)\,dr,
   \qquad s<t.
\]
This transformation uses only the projected mild equation: the Leray projection
is homogeneous of degree zero and commutes with the parabolic dilation, while
the Ornstein--Uhlenbeck kernel in \eqref{p2:eq:OU-semigroup} becomes the heat
kernel with time parameter \(t-s\).  No pressure formulation is used at this
stage.  The \(L^\infty\) norm of \(U\) on \(\R^3\times t(J)\) is bounded by a
constant depending only on \(J\) and \(\|V\|_{L^\infty}\).

The standard local smoothing theorem for bounded projected mild
Navier--Stokes solutions \cite{Kato1984,GigaMiyakawa1985} gives first local
derivative bounds away from the initial time in any mild representation.  In
the present ancient setting one may choose the initial time strictly before the
compact subcylinder under consideration.  Iterating the differentiated
Duhamel formula and then applying the interior parabolic bootstrap for the
Stokes system \cite{Ladyzhenskaya1968,Sohr2001} gives uniform bounds for all
derivatives of \(U\) on compact subcylinders of \(\R^3\times t(J)\).  The mild
formula may then be differentiated in time, so \(U\) solves the projected
Navier--Stokes equation classically.  Returning to \((y,\tau)\) gives the
asserted estimates for \(V\) and the classical projected equation
\eqref{p2:eq:projected}.  The argument is uniform for families with a common
\(L^\infty\) bound.
\end{proof}

\subsection{Pressure reconstruction}

\begin{lemma}[Local pressure reconstruction]\label{p2:lem:local-pressure}
Let \(V\) be a smooth bounded solution of the projected centered equation
\eqref{p2:eq:projected}.  Then for every ball \(B_R\) and compact interval
\(I\Subset\R\) there exists
\[
   P_R\in C^\infty(B_R\times I),
\]
unique up to an additive function of time, such that
\[
   \partial_\tau V+(V\cdot\nabla)V+\nabla P_R
   =
   \Delta V-\frac12(V+y\cdot\nabla V),
   \qquad \divergence V=0,
\]
on \(B_R\times I\).
\end{lemma}

\begin{proof}
Set
\[
   F=
   \Delta V-\frac12(V+y\cdot\nabla V)
   -\partial_\tau V-(V\cdot\nabla)V .
\]
Since \((V\cdot\nabla)V=\divergence(V\otimes V)\) and \(V\) satisfies
\eqref{p2:eq:projected}, self-adjointness of the Leray projection gives
\[
   \int_I\int_{B_R}F\cdot\Phi\,dy\,d\tau=0
\]
for every \(\Phi\in C_c^\infty(B_R\times I;\R^3)\) satisfying
\(\divergence_y\Phi=0\).  The local de Rham theorem, applied in the spatial
variables with \(\tau\) as a parameter, gives a scalar distribution \(P_R\) on
\(B_R\times I\), unique up to a distribution depending only on time, such that
\(\nabla_yP_R=F\).  Since \(F\) is smooth, all spatial derivatives of \(P_R\)
of positive order are smooth.  Fixing, for instance, the spatial mean of
\(P_R\) over \(B_R\) to be zero for each \(\tau\) selects a representative
which is smooth on \(B_R\times I\).  This proves the local pressure
formulation.
\end{proof}

\begin{lemma}[Leray pressure and local gauges]\label{p2:lem:pressure-reconstruction}
Let \(V\) be a bounded mild ancient solution satisfying
\[
   \sup_{\tau\in\R}\|V(\tau)\|_{L^3(\R^3)}<\infty .
\]
Then the whole-space Leray pressure
\[
   P=R_iR_j(V_iV_j)
\]
belongs to \(L^\infty(\R;L^{3/2}(\R^3))\).  Moreover, if \(V_n\to V\) locally
smoothly, the \(V_n\) have a common \(L^\infty_\tau L^3_y\) bound, and
\[
   P_n=R_iR_j(V_{n,i}V_{n,j}),
\]
then on every compact cylinder \(B_R\times[\tau_1,\tau_2]\) there are bounded
measurable gauges \(a_n(\tau)\) and \(a(\tau)\) such that, after passing to a
subsequence,
\[
   P_n-a_n(\tau)
   \rightharpoonup P-a(\tau)
   \quad\hbox{weakly in }L^{3/2}(B_R\times[\tau_1,\tau_2]).
\]
Equivalently, after replacing the limiting pressure by an equivalent
representative modulo a function of time, the weak limit may be denoted by
\(P\).
\end{lemma}

\begin{proof}
The first assertion follows from the Calderon--Zygmund boundedness of Riesz
transforms on \(L^{3/2}(\R^3)\) \cite{Stein1970}:
\[
   \|P(\tau)\|_{L^{3/2}}
   \le C\|V_i(\tau)V_j(\tau)\|_{L^{3/2}}
   \le C\|V(\tau)\|_{L^3}^2 .
\]
For the local stability statement, choose
\(\chi\in C_c^\infty(B_{4R})\) with \(\chi\equiv1\) on \(B_{3R}\).  Decompose
on each time slice
\[
   P_n=R_iR_j(\chi V_{n,i}V_{n,j})
       +R_iR_j((1-\chi)V_{n,i}V_{n,j}) .
\]
The local term converges strongly in \(L^q(B_{2R}\times[\tau_1,\tau_2])\) for
every finite \(q\).  This is because \(V_n\to V\) smoothly on
\(\overline B_{4R}\times[\tau_1,\tau_2]\), so
\(\chi V_{n,i}V_{n,j}\to\chi V_iV_j\) strongly in every finite \(L^q\), and
the Riesz transforms are Calderon--Zygmund operators on \(L^q\), \(1<q<\infty\).

It remains to control the exterior part modulo constants.  Let \(K_{ij}\) be
the convolution kernel of \(R_iR_j\).  For \(x\in B_R\) and
\(|z|>4R\),
\[
   |K_{ij}(x-z)-K_{ij}(-z)|\le C_R |z|^{-4}.
\]
Define the exterior gauges
\[
   b_n(\tau)
   =
   \int_{\R^3} K_{ij}(-z)(1-\chi(z))V_{n,i}(z,\tau)V_{n,j}(z,\tau)\,dz
\]
and \(b(\tau)\) by the same formula with \(V\) in place of \(V_n\).
The cutoff removes the singularity at the origin, and the integral is
absolutely convergent at infinity because, on the dyadic annulus
\(A_m=\{2^mR<|z|<2^{m+1}R\}\),
\[
   \int_{A_m}|z|^{-3}|V_n(z,\tau)|^2\,dz
   \le
   C(2^mR)^{-3}
   \|V_n(\tau)\|_{L^3(A_m)}^2|A_m|^{1/3}
   \le C L^2(2^mR)^{-2}.
\]
The same dyadic estimate shows that \(b_n\) and \(b\) are bounded uniformly on
\([\tau_1,\tau_2]\).  The maps \(\tau\mapsto b_n(\tau)\) and
\(\tau\mapsto b(\tau)\) are measurable, because they are limits of measurable
truncated integrals.
After subtracting \(b_n(\tau)\), the exterior contribution on \(B_R\) is
\[
   E_n(x,\tau)
   =
   \int_{\R^3}
   [K_{ij}(x-z)-K_{ij}(-z)]
   (1-\chi(z))V_{n,i}(z,\tau)V_{n,j}(z,\tau)\,dz .
\]
The same dyadic estimate, now using \(|z|^{-4}\), gives
\[
   \|E_n(\cdot,\tau)\|_{L^\infty(B_R)}
   \le C(R)L^2
\]
uniformly in \(n\) and \(\tau\).  Splitting the exterior integral into
\(\{4R<|z|<R'\}\) and \(\{|z|>R'\}\), the first part converges by local smooth
convergence, while the second part is bounded uniformly by \(C L^2/R'\).  Letting
first \(n\to\infty\) and then \(R'\to\infty\) identifies the local limit of
\(E_n\) with the corresponding exterior contribution for \(V\).  Therefore,
with \(a_n=b_n\) and \(a=b\), the decompositions of \(P_n-a_n\) and \(P-a\)
have local parts converging strongly and exterior parts converging
distributionally with a uniform local \(L^{3/2}\) bound.  Hence
\[
   P_n-a_n(\tau)\rightharpoonup P-a(\tau)
   \quad\text{weakly in }L^{3/2}(B_R\times[\tau_1,\tau_2]).
\]
\end{proof}

\subsection{Stability under locally smooth limits}

\begin{lemma}[Closure of the mild formulation]\label{p2:lem:mild-limit}
Let \(V_n\) be bounded mild ancient solutions satisfying
\[
   \sup_n\|V_n\|_{L^\infty(\R^3\times\R)}<\infty .
\]
If \(V_n\to V\) locally smoothly, then \(V\) is a bounded mild ancient solution
of \eqref{p2:eq:centered}.
\end{lemma}

\begin{proof}
The common \(L^\infty\) bound passes to \(V\) by pointwise convergence on
compact subsets.  Divergence freeness passes distributionally.  Fix
\(\sigma<\tau\) and test \eqref{p2:eq:mild} for \(V_n\) against
\(\phi\in C_c^\infty(\R^3;\R^3)\).  The endpoint terms are
\[
   \langle V_n(\tau),\phi\rangle,
   \qquad
   \langle V_n(\sigma),S(\tau-\sigma)^*\phi\rangle .
\]
The first converges by local smooth convergence.  For the second, note that
\(S(\tau-\sigma)^*\phi\in L^1(\R^3)\).  Approximate this \(L^1\) function by
a compactly supported smooth function and use the common \(L^\infty\) bound to
control the \(L^1\) approximation error uniformly in \(n\).

For the Duhamel term, write
\[
   \left\langle
      S(\tau-s)\mathbb P\divergence(V_n\otimes V_n)(s),\phi
   \right\rangle
   =
   \left\langle
      V_n(s)\otimes V_n(s),
      (S(\tau-s)\mathbb P\divergence)^*\phi
   \right\rangle .
\]
For fixed \(s\), the adjoint test belongs to \(L^1\); pairing convergence
again follows by compact approximation in \(L^1\), local smooth convergence,
and the common \(L^\infty\) bound on \(V_n\otimes V_n\).  The approximation
error is uniform in \(n\) for this fixed \(s\).  Moreover
\Cref{p2:lem:centered-stokes-kernel} gives the integrable domination
\[
   C\sup_n\|V_n\|_{L^\infty}^2
   e^{-(\tau-s)/2}(1-e^{-(\tau-s)})^{-1/2}\|\phi\|_{L^1},
   \qquad s\in(\sigma,\tau).
\]
Dominated convergence in \(s\) passes the Duhamel integral to the limit.
Thus \(V\) satisfies \eqref{p2:eq:mild}.
\end{proof}

\subsection{Compactness of bounded families}

\begin{proposition}[Compactness of bounded ancient families]
\label{p2:prop:bounded-compactness}
Let \(\{V_n\}\) be bounded mild ancient solutions satisfying
\[
   \sup_n\|V_n\|_{L^\infty(\R^3\times\R)}<\infty .
\]
Then a subsequence converges locally smoothly to a bounded mild ancient
solution \(V\).
\end{proposition}

\begin{proof}
By \Cref{p2:lem:mild-smoothing}, the sequence is bounded in every
\(C^m\)-norm on every compact cylinder.  Arzela--Ascoli on a countable
exhaustion \(B_N\times[-N,N]\), followed by a diagonal extraction over
\((m,N)\), gives a subsequence converging locally smoothly to a smooth limit.
\Cref{p2:lem:mild-limit} shows that the limit is again a bounded mild ancient
solution.
\end{proof}

\section{\texorpdfstring{Uniform \(L^3\)-tightness and stationary \(L^3\) rigidity}{Uniform L3-tightness and stationary L3 rigidity}}
\label{p2:sec:tightness-rigidity}

\subsection{\texorpdfstring{Uniform \(L^3\)-tightness}{Uniform L3-tightness}}

\begin{definition}[Uniform \(L^3\)-tightness]\label{p2:def:L3-tight}
A bounded ancient solution \(V\) is uniformly \(L^3\)-tight if
\[
   \lim_{R\to\infty}
   \sup_{\tau\in\R}
   \int_{|y|>R}|V(y,\tau)|^3\,dy=0 .
\]
\end{definition}

For the compactness statements in this section we sometimes assume, in
addition, a common \(L^\infty_\tau L^3_y\) bound.  The endpoint Liouville
argument in \Cref{p2:sec:endpoint-L3-closure} does not take this bound as a
separate hypothesis; it obtains the required sequence estimate from boundedness
and the tail condition above.

\subsection{Closedness under locally smooth convergence}

\begin{lemma}[Closure of uniform \(L^3\)-tightness]\label{p2:lem:tight-closure}
Let \(V_n\) be bounded mild ancient solutions satisfying
\[
   \sup_n\sup_{\tau\in\R}\|V_n(\tau)\|_{L^3(\R^3)}\le L<\infty
\]
and assume that they have a common tightness modulus: for some
\(\omega:[1,\infty)\to[0,\infty)\) with \(\omega(R)\to0\),
\[
   \sup_n\sup_{\tau\in\R}
   \int_{|y|>R}|V_n(y,\tau)|^3\,dy\le \omega(R),
   \qquad R\ge1.
\]
If \(V_n\to V\) locally smoothly, then
\[
   \sup_{\tau\in\R}\|V(\tau)\|_{L^3(\R^3)}\le L
\]
and
\[
   \sup_{\tau\in\R}
   \int_{|y|>R}|V(y,\tau)|^3\,dy\le \omega(R),
   \qquad R\ge1.
\]
In particular \(V\) is uniformly \(L^3\)-tight.
\end{lemma}

\begin{proof}
Fix \(\tau\).  On \(B_R\), local smooth convergence gives
\[
   \int_{B_R}|V(y,\tau)|^3\,dy
   =
   \lim_{n\to\infty}\int_{B_R}|V_n(y,\tau)|^3\,dy
   \le L^3 .
\]
Letting \(R\to\infty\) and using monotone convergence gives
\(\|V(\tau)\|_{L^3}\le L\).  Since \(\tau\) was arbitrary, the same bound is
uniform in time.

For the tail bound, fix \(1\le R<R'\).  On the bounded annulus
\(\{R<|y|<R'\}\), dominated convergence gives
\[
   \int_{R<|y|<R'}|V(y,\tau)|^3\,dy
   =
   \lim_{n\to\infty}
   \int_{R<|y|<R'}|V_n(y,\tau)|^3\,dy
   \le \omega(R).
\]
Letting \(R'\to\infty\) and using monotone convergence gives
\[
   \int_{|y|>R}|V(y,\tau)|^3\,dy\le \omega(R).
\]
The estimate is independent of \(\tau\), so taking the supremum in time gives
the asserted tail bound.  Since \(\omega(R)\to0\), \(V\) is uniformly
\(L^3\)-tight.
\end{proof}

\begin{proof}[Proof of \Cref{p2:thm:tight-compactness-intro}]
The compactness assertion is \Cref{p2:prop:bounded-compactness}.  The
\(L^\infty_\tau L^3_y\) and tightness closure assertions follow from
\Cref{p2:lem:tight-closure} with the displayed constants \(L\) and \(\omega\).
\end{proof}

\subsection{Stationary time-translate limits}

The stationary self-similar profile equation used below is
\begin{equation}\label{p2:eq:stationary-profile}
   (w\cdot\nabla)w+\nabla p
   =
   \Delta w-\frac12(w+y\cdot\nabla w),
   \qquad \divergence w=0 .
\end{equation}

\begin{theorem}[Stationary limits under uniform \(L^3\)-tightness]
\label{p2:thm:stationary-limits-tight}
Let \(V\) be uniformly \(L^3\)-tight.  Let \(\tau_n\in\R\), and suppose
\[
   \Theta_{\tau_n}V\to W
\]
locally smoothly.  If \(W\) is stationary, then \(W\equiv0\).
\end{theorem}

\begin{proof}
The sequence \(\Theta_{\tau_n}V\) has the same \(L^3\) bound and tightness
modulus as \(V\).  By \Cref{p2:lem:tight-closure}, the limit \(W\) belongs to
\(L^\infty_\tau L^3_y\).  If \(W\) is stationary, write
\(W(y,\tau)=w(y)\).  Then \(w\in L^3(\R^3)\).

By \Cref{p2:lem:mild-limit}, \(W\) is a bounded mild ancient solution.  The
smoothness from \Cref{p2:lem:mild-smoothing} and the local pressure lemma
\Cref{p2:lem:local-pressure} give smooth local pressures \(p_N\) on \(B_N\) at
time \(0\).  Since
\(\partial_\tau W=0\), each \(p_N\) satisfies
\[
   (w\cdot\nabla)w+\nabla p_N
   =
   \Delta w-\frac12(w+y\cdot\nabla w)
   \quad\text{on }B_N .
\]
On overlaps the gradients of the local pressures agree, so the pressures
differ only by constants.  Adjusting those constants and passing to the
compatible union gives a scalar distribution \(p\) on \(\R^3\) for which
\eqref{p2:eq:stationary-profile} holds in distributions.  The stationary
\(L^3\) rigidity theorem \Cref{p2:thm:NRS} applies and gives \(w=0\).
\end{proof}

\subsection{Stationary self-similar rigidity}

\begin{theorem}[Stationary self-similar rigidity]\label{p2:thm:NRS}
Let \(w\in C^\infty(\R^3;\R^3)\cap L^3(\R^3)\) solve
\eqref{p2:eq:stationary-profile} in distributions for some scalar distribution
\(p\), with \(\divergence w=0\).  Then \(w\equiv0\).
\end{theorem}

\begin{proof}
This is the theorem of Ne\v{c}as, R\r{u}\v{z}i\v{c}ka, and
\v{S}ver\'ak \cite{NRS1996}, applied to the stationary Leray self-similar
system.  The smooth distributional formulation stated here is a special case
of their weak \(L^3(\R^3)\) theorem.
\end{proof}

\section{Normalized collections and minimal ancient solutions}
\label{p2:sec:minimal-elements}

\subsection{Closed normalized collections}

\begin{definition}[Closed normalized collection]\label{p2:def:closed-normalized}
A collection \(A\) of bounded mild ancient solutions is called a closed
normalized collection if:
\begin{enumerate}[label=\textup{(\roman*)}]
\item there is \(M<\infty\) such that
\[
   \|V\|_{L^\infty(\R^3\times\R)}\le M
   \qquad\text{for all }V\in A;
\]
\item there are \(R_0>0\) and \(\eta_0>0\) such that
\begin{equation}\label{p2:eq:local-nonzero}
   \sup_{\tau\in\R}
   \int_{B_{R_0}}|V(y,\tau)|^3\,dy
   \ge \eta_0
   \qquad\text{for all }V\in A;
\end{equation}
\item \(A\) is invariant under time translation:
\[
   V\in A,\ s\in\R
   \quad\Longrightarrow\quad
   \Theta_sV\in A;
\]
\item if \(V_n\in A\), \(V_n\to V\) locally smoothly, and \(V\) satisfies the
same \((R_0,\eta_0)\)-nonvanishing condition \eqref{p2:eq:local-nonzero}, then
\(V\in A\).
\end{enumerate}
\end{definition}

\subsection{Time-centering}

\begin{lemma}[Time-centering]\label{p2:lem:time-centering}
Let \(A\) be a closed normalized collection.  If \(V\in A\), then for every
\(\delta>0\) there exists \(s\in\R\) such that
\[
   \int_{B_{R_0}}|\Theta_sV(y,0)|^3\,dy
   \ge \eta_0-\delta.
\]
\end{lemma}

\begin{proof}
By \eqref{p2:eq:local-nonzero}, the supremum in time of the local \(L^3\) mass is
at least \(\eta_0\).  Choose \(s\) within \(\delta\) of this supremum.  Since
\((\Theta_sV)(y,0)=V(y,s)\), the displayed estimate follows.  Time-translation
invariance keeps \(\Theta_sV\) in \(A\).
\end{proof}

\subsection{Stability of local nonvanishing}

\begin{lemma}[Stability of local nonvanishing]\label{p2:lem:nonzero-stability}
If \(V_n\to V\) locally smoothly and
\[
   \int_{B_{R_0}}|V_n(y,0)|^3\,dy\ge \eta_0-o(1),
\]
then
\[
   \int_{B_{R_0}}|V(y,0)|^3\,dy\ge \eta_0 .
\]
\end{lemma}

\begin{proof}
Local smooth convergence implies uniform convergence on
\(\overline B_{R_0}\times\{0\}\), hence convergence in
\(L^3(B_{R_0})\).  The map \(f\mapsto\int_{B_{R_0}}|f|^3\) is continuous on
\(L^3(B_{R_0})\), and the claim follows by passing to the limit.
\end{proof}

\subsection{Lower semicontinuity of size}

\begin{lemma}[Lower semicontinuity of \(L^\infty\) size]
\label{p2:lem:lsc-size}
If \(V_n\to V\) locally smoothly and
\[
   \sup_n\|V_n\|_{L^\infty(\R^3\times\R)}<\infty,
\]
then
\[
   \|V\|_{L^\infty(\R^3\times\R)}
   \le
   \liminf_{n\to\infty}
   \|V_n\|_{L^\infty(\R^3\times\R)}.
\]
\end{lemma}

\begin{proof}
Let \(L_0\) be the liminf on the right.  If
\(|V(y_0,\tau_0)|>L_0+\eps\), choose a subsequence with
\(\|V_n\|_{L^\infty}\to L_0\).  Local uniform convergence at
\((y_0,\tau_0)\) gives
\(|V_n(y_0,\tau_0)|>L_0+\eps/2\) for all large \(n\), contradicting
\(\|V_n\|_{L^\infty}\to L_0\).  Hence \(|V|\le L_0\) pointwise.  Since the
smooth representative has pointwise supremum equal to its essential
\(L^\infty\) norm, the result follows.
\end{proof}

\subsection{Existence of a minimal element}

\begin{theorem}[Existence of a minimal ancient solution]
\label{p2:thm:minimal-existence}
Let \(A\) be nonempty and closed normalized.  Set
\[
   m=
   \inf_{V\in A}
   \|V\|_{L^\infty(\R^3\times\R)}.
\]
Then \(0<m<\infty\), and there exists \(V_*\in A\) such that
\[
   \|V_*\|_{L^\infty(\R^3\times\R)}=m.
\]
In particular \(V_*\not\equiv0\).
\end{theorem}

\begin{proof}
The common \(L^\infty\) bound gives \(m<\infty\).  The local nonvanishing gives
a positive lower bound:
\[
   \eta_0
   \le
   \sup_{\tau\in\R}\int_{B_{R_0}}|V(y,\tau)|^3\,dy
   \le
   |B_{R_0}|\|V\|_{L^\infty(\R^3\times\R)}^3,
\]
so
\[
   m\ge \left(\eta_0/|B_{R_0}|\right)^{1/3}>0.
\]

Choose \(V_n\in A\) with
\[
   m\le \|V_n\|_{L^\infty(\R^3\times\R)}\le m+\frac1n .
\]
By \Cref{p2:lem:time-centering}, after replacing \(V_n\) by a time translate,
which stays in \(A\) and preserves the \(L^\infty\) norm, we may assume
\[
   \int_{B_{R_0}}|V_n(y,0)|^3\,dy\ge \eta_0-\frac1n .
\]
The common \(L^\infty\) bound and \Cref{p2:prop:bounded-compactness} give a
locally smooth subsequential limit \(V_*\), which is a bounded mild ancient
solution.  By \Cref{p2:lem:nonzero-stability},
\[
   \int_{B_{R_0}}|V_*(y,0)|^3\,dy\ge\eta_0,
\]
so \(V_*\) retains the fixed \((R_0,\eta_0)\)-normalization.  Closedness of
\(A\) then gives \(V_*\in A\).  Finally, \Cref{p2:lem:lsc-size} gives
\[
   \|V_*\|_{L^\infty}
   \le
   \liminf_{n\to\infty}\|V_n\|_{L^\infty}
   =m.
\]
Since \(m\) is the infimum over \(A\) and \(V_*\in A\), equality holds.  The
positive lower bound on \(m\) rules out \(V_*\equiv0\).
\end{proof}

\subsection{The uniformly tight normalized class}

\begin{definition}[Uniformly tight normalized collection]
\label{p2:def:tight-normalized-collection}
Fix \(M,L,R_0,\eta_0>0\) and a modulus
\[
   \omega:[1,\infty)\to[0,\infty),
   \qquad \omega(R)\to0.
\]
A collection \(A\) of bounded mild ancient solutions is called a uniformly
tight normalized collection with data \((M,L,\omega,R_0,\eta_0)\) if:
\begin{enumerate}[label=\textup{(\arabic*)}]
\item for every \(V\in A\),
\[
   \|V\|_{L^\infty(\R^3\times\R)}\le M;
\]
\item for every \(V\in A\),
\[
   \sup_{\tau\in\R}\|V(\tau)\|_{L^3(\R^3)}\le L;
\]
\item for every \(V\in A\) and every \(R\ge1\),
\[
   \sup_{\tau\in\R}
   \int_{|y|>R}|V(y,\tau)|^3\,dy\le \omega(R);
\]
\item for every \(V\in A\),
\[
   \sup_{\tau\in\R}
   \int_{B_{R_0}}|V(y,\tau)|^3\,dy\ge\eta_0;
\]
\item \(A\) is invariant under time translation;
\item \(A\) is closed under locally smooth limits retaining item \textup{(4)}.
\end{enumerate}
\end{definition}

\begin{proposition}[Uniformly tight normalized collections are closed normalized]
\label{p2:prop:tight-closed-normalized}
Every uniformly tight normalized collection is a closed normalized collection.
\end{proposition}

\begin{proof}
The \(L^\infty\) bound, local nonvanishing, time-translation invariance, and
closedness clause in \Cref{p2:def:closed-normalized} are part of
\Cref{p2:def:tight-normalized-collection}.  The only non-formal point is that
the stated \(L^3\) bound and fixed tightness modulus are preserved under locally smooth
limits.  This follows from \Cref{p2:lem:tight-closure}.
\end{proof}

\begin{proof}[Proof of \Cref{p2:thm:minimal-tight-intro}]
By \Cref{p2:prop:tight-closed-normalized}, the collection is closed normalized.
\Cref{p2:thm:minimal-existence} gives the asserted nonzero minimizer.
\end{proof}

\section{Compact trajectory hulls}
\label{p2:sec:compact-hulls}

\subsection{Trajectory hulls}

\begin{definition}[Trajectory hull]\label{p2:def:trajectory-hull}
For a bounded mild ancient solution \(V\), its trajectory hull is
\[
   \calH(V)
   =
   \overline{\{\Theta_sV:\ s\in\R\}}^{C^\infty_{\loc}}.
\]
\end{definition}

\subsection{Compactness of bounded hulls}

\begin{theorem}[Compact trajectory hull]\label{p2:thm:compact-hull}
If \(V\) is a bounded mild ancient solution, then \(\calH(V)\) is compact in
the locally smooth topology and invariant under every time translation
\(\Theta_s\).  Every \(W\in\calH(V)\) is a bounded mild ancient solution with
\[
   \|W\|_{L^\infty(\R^3\times\R)}
   \le
   \|V\|_{L^\infty(\R^3\times\R)}.
\]
\end{theorem}

\begin{proof}
Every translate \(\Theta_sV\) is a bounded mild ancient solution with the same
\(L^\infty\) bound as \(V\).  Hence every sequence of translates has a
locally smoothly convergent subsequence by \Cref{p2:prop:bounded-compactness}.
If \((W_n)\subset\calH(V)\), choose translates \(\Theta_{s_n}V\) with
\(d(W_n,\Theta_{s_n}V)<1/n\).  A subsequence of \(\Theta_{s_n}V\) converges
locally smoothly, and the corresponding subsequence of \(W_n\) has the same
limit.  Thus \(\calH(V)\) is sequentially compact.  Since the locally smooth
topology is metrizable, it is compact.

The \(L^\infty\) estimate and the mild equation pass to elements of the hull
by local smooth convergence and \Cref{p2:lem:mild-limit}.  For invariance, note
that \(\Theta_t\) maps the orbit of \(V\) onto itself.  If \(W_n\to W\)
locally smoothly, then \(\Theta_tW_n\to\Theta_tW\) locally smoothly because
time translation only shifts the compact time interval on which the
seminorms are evaluated.  Hence \(\Theta_t\calH(V)=\calH(V)\).

\end{proof}

\begin{corollary}[Uniform tightness on hulls]\label{p2:cor:tightness-on-hulls}
If \(V\) is uniformly \(L^3\)-tight with bound \(L\) and modulus \(\omega\),
then every \(W\in\calH(V)\) satisfies
\[
   \sup_{\tau\in\R}\|W(\tau)\|_{L^3(\R^3)}\le L
\]
and
\[
   \sup_{\tau\in\R}\int_{|y|>R}|W(y,\tau)|^3\,dy\le\omega(R),
   \qquad R\ge1.
\]
The same estimates hold uniformly on every compact invariant subset
\(K\subset\calH(V)\).
\end{corollary}

\begin{proof}
Every translate of \(V\) has the same \(L^3\) bound and tightness modulus.
Applying \Cref{p2:lem:tight-closure} to any sequence of translates converging
locally smoothly to \(W\in\calH(V)\) gives the displayed estimates for \(W\).
Taking the supremum over \(W\in K\subset\calH(V)\) gives the final assertion.
\end{proof}

\subsection{Compact invariant sets}

\begin{definition}[Compact invariant set]\label{p2:def:compact-invariant}
A nonempty compact subset \(K\) of the locally smooth space of bounded mild
ancient solutions is invariant if
\[
   \Theta_tK=K
   \qquad\text{for every }t\in\R .
\]
It is minimal if it contains no proper nonempty compact invariant subset.
\end{definition}

\begin{lemma}[Continuity of the translation action]\label{p2:lem:action-continuity}
Let \(K\) be a compact invariant set.  For each \(t\in\R\),
\(\Theta_t:K\to K\) is a homeomorphism, and the action
\[
   \R\times K\ni(t,W)\mapsto \Theta_tW\in K
\]
is jointly continuous.
\end{lemma}

\begin{proof}
For fixed \(t\), convergence \(W_n\to W\) locally smoothly implies
\[
   \|\Theta_tW_n-\Theta_tW\|_{C^m(B_R\times[-T,T])}
   =
   \|W_n-W\|_{C^m(B_R\times[-T+t,T+t])}\to0.
\]
Thus \(\Theta_t\) is continuous, and its inverse is \(\Theta_{-t}\).  If
\(t_n\to t\) and \(W_n\to W\), then on any fixed compact cylinder the
difference
\[
   \Theta_{t_n}W_n-\Theta_tW
\]
is bounded by a term involving \(W_n-W\) on a slightly larger compact cylinder
and a term involving the uniform continuity of finitely many derivatives of
\(W\) on that larger cylinder.  Both terms tend to zero, proving joint
continuity.
\end{proof}

\subsection{Minimal compact invariant subsets}

\begin{proposition}[Minimal compact invariant subsets]\label{p2:prop:minimal-subset}
Let \(K\) be a nonempty compact invariant set.  Then \(K\) contains a
nonempty compact minimal invariant subset.
\end{proposition}

\begin{proof}
Let \(\mathcal F\) be the family of nonempty compact invariant subsets of
\(K\), ordered by reverse inclusion.  For a chain \(\{K_\alpha\}\), the
intersection \(\bigcap_\alpha K_\alpha\) is nonempty by compactness and the
finite intersection property, since the chain is totally ordered by ordinary
inclusion.  The intersection is compact and invariant: invariance follows
from \(\Theta_tK_\alpha=K_\alpha\) for all \(\alpha\) and the fact that
\(\Theta_t\) is a homeomorphism.  Thus every chain has an upper bound in the
reverse-inclusion order.  Zorn's lemma gives a maximal element for that order,
equivalently a minimal compact invariant subset under ordinary inclusion.
\end{proof}

\subsection{The zero trajectory}

\begin{lemma}[The zero trajectory]\label{p2:lem:zero-trajectory}
The singleton \(\{0\}\) is compact and invariant.  If a minimal compact
invariant set contains the zero trajectory, then it equals \(\{0\}\).
\end{lemma}

\begin{proof}
The zero solution satisfies \(\Theta_t0=0\) for all \(t\).  Therefore
\(\{0\}\) is a nonempty compact invariant subset of any invariant set
containing zero.  Minimality forces equality.
\end{proof}

\section{Invariant measures and the covariance obstruction}
\label{p2:sec:measures-covariance}

\subsection{Krylov--Bogolyubov averages}

The following is the standard Krylov--Bogolyubov averaging argument
\cite{KryloffBogoliouboff1937,Billingsley1999}.

\begin{theorem}[Krylov--Bogolyubov averages]\label{p2:thm:KB}
Let \(K\) be a compact invariant set and \(W_0\in K\).  For \(T>0\), define
\[
   \mu_T
   =
   \frac1T\int_0^T\delta_{\Theta_sW_0}\,ds .
\]
Then every sequence \(T_n\to\infty\) has a subsequence such that
\[
   \mu_{T_n}\rightharpoonup\mu
\]
weakly as Borel probability measures on \(K\), and every such limit \(\mu\)
is invariant:
\[
   (\Theta_t)_\#\mu=\mu
   \qquad\text{for all }t\in\R .
\]
If \(K\) is minimal, then every invariant Borel probability measure on \(K\)
has support equal to \(K\).
\end{theorem}

\begin{proof}
The orbit map \(s\mapsto\Theta_sW_0\) is continuous by
\Cref{p2:lem:action-continuity}; thus \(\mu_T\) is the push-forward of normalized
Lebesgue measure on \([0,T]\).  Since \(K\) is compact and metrizable, the
space of Borel probability measures on \(K\) is weakly sequentially compact
\cite{Billingsley1999}.  Hence \(\mu_{T_n}\) has a weakly convergent
subsequence.

For \(F\in C(K)\) and \(t>0\),
\[
   \int_K F(\Theta_tW)\,d\mu_T(W)-\int_KF(W)\,d\mu_T(W)
   =
   \frac1T\int_T^{T+t}F(\Theta_sW_0)\,ds
   -
   \frac1T\int_0^tF(\Theta_sW_0)\,ds .
\]
The absolute value is at most \(2t\|F\|_{C(K)}/T\).  Letting
\(T=T_{n_k}\to\infty\) gives invariance for \(t>0\); negative times follow by
applying the positive-time result to \(F\circ\Theta_t\).

If \(K\) is minimal and \(\mu\) is invariant, then \(\supp\mu\) is nonempty,
closed, and invariant.  Indeed, if \(U\) is a neighborhood of
\(\Theta_tW\), then \(\Theta_{-t}U\) is a neighborhood of \(W\), and invariance
shows \(\mu(U)=\mu(\Theta_{-t}U)>0\) whenever \(W\in\supp\mu\).  Hence
\(\Theta_t(\supp\mu)\subset\supp\mu\).  Applying the same argument with
\(-t\) gives the reverse inclusion, so \(\supp\mu\) is invariant.  It is
nonempty because \(\mu(K)=1\), and it is compact because it is closed in
compact \(K\).  Minimality forces \(\supp\mu=K\).
\end{proof}

\subsection{Support on minimal hulls}

\begin{corollary}[Nonzero support]\label{p2:cor:nonzero-support}
Let \(K\) be a minimal compact invariant set with \(K\ne\{0\}\).  Then zero is
not in \(K\), and every invariant probability measure on \(K\) has support
equal to \(K\).
\end{corollary}

\begin{proof}
If zero were in \(K\), \Cref{p2:lem:zero-trajectory} would give \(K=\{0\}\),
contrary to the assumption.  The support conclusion is the minimal-support part
of \Cref{p2:thm:KB}.
\end{proof}

\subsection{Stationary statistical identity}

\begin{theorem}[Stationary statistical form]\label{p2:thm:statistical-identity}
Let \(K\) be a compact invariant set and \(\mu\) an invariant Borel
probability measure on \(K\).  Then for every solenoidal
\(\phi\in C_c^\infty(\R^3;\R^3)\),
\begin{equation}\label{p2:eq:statistical-identity}
\begin{aligned}
0
=
\int_K
\bigg[
   &\int_{\R^3}
   W(y,0)\cdot
   \left(\Delta\phi+\phi+\frac12 y\cdot\nabla\phi\right)\,dy  \\
   &+
   \int_{\R^3}
   W_i(y,0)W_j(y,0)\partial_j\phi_i(y)\,dy
\bigg]\,d\mu(W).
\end{aligned}
\end{equation}
\end{theorem}

\begin{proof}
Define
\[
   F(W)=\int_{\R^3}W(y,0)\cdot\phi(y)\,dy .
\]
Local smooth convergence makes \(F\) continuous on \(K\): if \(W_n\to W\),
then \(W_n(\cdot,0)\to W(\cdot,0)\) uniformly on the compact support of
\(\phi\).  Invariance gives
\[
   \int_K\frac{F(\Theta_hW)-F(W)}{h}\,d\mu(W)=0
   \qquad (0<|h|\le1).
\]
Choose \(R\) with \(\supp\phi\subset B_R\).  On
\(\overline B_R\times[-1,1]\), compactness of \(K\) in the locally smooth
topology gives
\[
   M_\phi
   =
   \sup_{W\in K}
   \|\partial_\tau W\|_{L^\infty(\overline B_R\times[-1,1])}
   <\infty .
\]
For \(0<|h|\le1\),
\[
   \frac{F(\Theta_hW)-F(W)}{h}
   =
   \int_{\R^3}
   \left(\frac1h\int_0^h\partial_\tau W(y,s)\,ds\right)\cdot\phi(y)\,dy,
\]
and the absolute value is bounded by
\(|B_R|\|\phi\|_{L^\infty}M_\phi\), uniformly in \(W\) and \(h\).  For each
fixed \(W\), smoothness gives uniform convergence of the inner average to
\(\partial_\tau W(y,0)\) on \(\supp\phi\).  Dominated convergence on \(K\)
yields
\[
   \int_K\int_{\R^3}\partial_\tau W(y,0)\cdot\phi(y)\,dy\,d\mu(W)=0 .
\]
For each \(W\), the smooth local equation \eqref{p2:eq:centered} at time zero,
tested against solenoidal \(\phi\), gives
\[
\begin{aligned}
\int_{\R^3}\partial_\tau W\cdot\phi\,dy
=
&\int_{\R^3} W\cdot\Delta\phi\,dy
 +\int_{\R^3} W\cdot\phi\,dy
 +\frac12\int_{\R^3}W\cdot(y\cdot\nabla\phi)\,dy\\
&+\int_{\R^3}W_iW_j\partial_j\phi_i\,dy .
\end{aligned}
\]
The pressure term vanishes because \(\divergence\phi=0\); the drift identity
uses
\[
   \int y_j\partial_jW_i\,\phi_i\,dy
   =
   -3\int W\cdot\phi\,dy-\int W\cdot(y\cdot\nabla\phi)\,dy .
\]
The right-hand side is a continuous bounded function of \(W\in K\).  Averaging
over \(K\) gives \eqref{p2:eq:statistical-identity}.
\end{proof}

\subsection{Barycenter and covariance}

\begin{theorem}[Barycenter and covariance]\label{p2:thm:barycenter-covariance}
Let \(K\) and \(\mu\) be as in \Cref{p2:thm:statistical-identity}.  Define
\[
   \overline W(y)=\int_K W(y,0)\,d\mu(W)
\]
and
\[
   C_{ij}(y)
   =
   \int_K W_i(y,0)W_j(y,0)\,d\mu(W)
   -
   \overline W_i(y)\overline W_j(y).
\]
Then \(\overline W\) and \(C\) are \(C^\infty_{\loc}\), \(\overline W\) is
solenoidal, and for every solenoidal
\(\phi\in C_c^\infty(\R^3;\R^3)\),
\begin{equation}\label{p2:eq:barycenter-covariance}
   \int_{\R^3}
   \overline W\cdot
   \left(\Delta\phi+\phi+\frac12 y\cdot\nabla\phi\right)\,dy
   +
   \int_{\R^3}
   \overline W_i\overline W_j\partial_j\phi_i\,dy
   +
   \int_{\R^3}C_{ij}\partial_j\phi_i\,dy
   =0 .
\end{equation}
In particular, \(\overline W\) is a weak stationary self-similar solution
against solenoidal test fields if and only if
\[
   \int_{\R^3}C_{ij}\partial_j\phi_i\,dy=0
\]
for every solenoidal test field \(\phi\).
\end{theorem}

\begin{proof}
Compactness of \(K\) in the locally smooth topology gives, for every compact
ball and every derivative order, a uniform bound on
\(\partial_y^\alpha W(y,0)\) for \(W\in K\).  Hence one may differentiate
\(\overline W\) under the integral sign.  More explicitly, for every
multiindex \(\alpha\) and every \(y\) in a fixed compact ball,
\[
   \partial_y^\alpha\overline W(y)
   =
   \int_K\partial_y^\alpha W(y,0)\,d\mu(W),
\]
because the difference quotients are dominated by the next derivative
seminorm, uniformly on \(K\).  The same argument applies to the second moment
\[
   M_{ij}(y)=\int_KW_i(y,0)W_j(y,0)\,d\mu(W),
\]
using Leibniz' rule and the uniform bounds on derivatives of \(W\).  This
gives \(C^\infty_{\loc}\) regularity of \(\overline W\), \(M\), and
\(C=M-\overline W\otimes\overline W\).  Since every \(W\) is divergence free,
Fubini against scalar test functions gives \(\divergence\overline W=0\):
for \(\eta\in C_c^\infty(\R^3)\),
\[
   \int_{\R^3}\overline W\cdot\nabla\eta\,dy
   =
   \int_K\int_{\R^3}W(y,0)\cdot\nabla\eta(y)\,dy\,d\mu(W)=0.
\]

For each compactly supported test field \(\phi\), the maps
\[
   W\mapsto\int_{\R^3}W(y,0)\cdot\phi(y)\,dy,
   \qquad
   W\mapsto\int_{\R^3}W_i(y,0)W_j(y,0)\partial_j\phi_i(y)\,dy
\]
are continuous and bounded on \(K\), again by compactness in
\(C^\infty_{\loc}\).  Thus all integrals below are finite, and Fubini's theorem
applies.

Applying \Cref{p2:thm:statistical-identity} and passing the linear term through
the \(K\)-integral gives the first integral in
\eqref{p2:eq:barycenter-covariance}.  Passing the nonlinear term through the
integral gives
\[
   \int_K W_iW_j\,d\mu
   =
   \overline W_i\overline W_j+C_{ij}.
\]
Substitution yields \eqref{p2:eq:barycenter-covariance}.  Comparing this identity
with the pressure-free stationary weak formulation for \(\overline W\) gives
the stated equivalence with vanishing covariance contribution.
\end{proof}

\section{\texorpdfstring{Endpoint \(L^3\) closure of the uniformly tight class}{Endpoint L3 closure of the uniformly tight class}}
\label{p2:sec:endpoint-L3-closure}

\subsection{Physical pullback}

Let \(V\) be a smooth centered ancient velocity and let \(P\) be a compatible
local pressure.  For \(t<0\) define
\begin{equation}\label{p2:eq:physical-pullback}
   u(x,t)=(-t)^{-1/2}
   V\left(\frac{x}{\sqrt{-t}},-\log(-t)\right),
\end{equation}
and
\begin{equation}\label{p2:eq:physical-pressure-pullback}
   p(x,t)=(-t)^{-1}
   P\left(\frac{x}{\sqrt{-t}},-\log(-t)\right).
\end{equation}
Equivalently,
\[
   V(y,\tau)=e^{-\tau/2}u(e^{-\tau/2}y,-e^{-\tau}),
   \qquad
   P(y,\tau)=e^{-\tau}p(e^{-\tau/2}y,-e^{-\tau}).
\]

\begin{lemma}[Pullback solves physical Navier--Stokes]
\label{p2:lem:physical-pullback}
Let \(V\) be a bounded mild ancient solution of \eqref{p2:eq:centered}.  Then
the pullback \(u\) in \eqref{p2:eq:physical-pullback} is a smooth divergence-free
solution of
\begin{equation}\label{p2:eq:physical-NS}
   \partial_tu+(u\cdot\nabla)u+\nabla p=\Delta u,
   \qquad \divergence u=0,
\end{equation}
on \(\R^3\times(-\infty,0)\), with pressure \(p\) locally given by
\eqref{p2:eq:physical-pressure-pullback}.  Moreover, for every \(T<0\),
\[
   \|u\|_{L^\infty(\R^3\times(-\infty,T])}
   \le
   (-T)^{-1/2}\|V\|_{L^\infty(\R^3\times\R)}.
\]
\end{lemma}

\begin{proof}
By \Cref{p2:lem:mild-smoothing}, \(V\) is smooth.  By
\Cref{p2:lem:local-pressure}, on every compact parabolic cylinder there is a
smooth pressure representative \(P\), unique up to a function of \(\tau\), for
which \eqref{p2:eq:centered} holds classically.  Adding a function of \(\tau\) to
\(P\) only adds a function of \(t\) to \(p\), and hence does not affect
\(\nabla p\).

Set
\[
   \lambda(t)=(-t)^{-1/2},\qquad
   y=\lambda(t)x,\qquad
   \tau=-\log(-t).
\]
Then \(\dot \tau=\lambda^2\) and
\[
   \partial_t y=\frac12\lambda^2 y,
   \qquad
   \dot\lambda=\frac12\lambda^3 .
\]
Using \(u=\lambda V(y,\tau)\), one obtains
\[
   \partial_tu
   =
   \lambda^3
   \left(
      \partial_\tau V+\frac12y\cdot\nabla V+\frac12V
   \right),
\]
while
\[
   (u\cdot\nabla_x)u=\lambda^3(V\cdot\nabla_y)V,
   \qquad
   \Delta_xu=\lambda^3\Delta_yV,
   \qquad
   \nabla_xp=\lambda^3\nabla_yP .
\]
Substituting these identities into \eqref{p2:eq:centered} gives
\eqref{p2:eq:physical-NS}.  The divergence identity follows similarly:
\[
   \divergence_xu=\lambda^2\divergence_yV=0.
\]
Finally, if \(t\le T<0\), then \((-t)^{-1/2}\le(-T)^{-1/2}\), and the stated
\(L^\infty\) bound follows directly from \eqref{p2:eq:physical-pullback}.
\end{proof}

\subsection{\texorpdfstring{Critical \(L^3\) invariance}{Critical L3 invariance}}

\begin{lemma}[Invariance of the critical norm]
\label{p2:lem:L3-pullback-invariance}
Let \(u\) be the physical pullback of \(V\).  Then, for every \(t<0\),
\begin{equation}\label{p2:eq:L3-pullback-invariance}
   \|u(t)\|_{L^3(\R^3)}
   =
   \|V(-\log(-t))\|_{L^3(\R^3)} .
\end{equation}
\end{lemma}

\begin{proof}
With \(y=x/\sqrt{-t}\), \eqref{p2:eq:physical-pullback} gives
\[
   |u(x,t)|^3=(-t)^{-3/2}|V(y,-\log(-t))|^3,
   \qquad
   dx=(-t)^{3/2}dy .
\]
Multiplying these two identities and integrating over \(\R^3\) gives
\eqref{p2:eq:L3-pullback-invariance}.
\end{proof}

\begin{lemma}[Tail tightness gives an endpoint sequence]
\label{p2:lem:tail-tight-sequence-L3}
Let \(V\) be a bounded function on \(\R^3\times\R\) satisfying
\[
   \lim_{R\to\infty}
   \sup_{\tau\in\R}\int_{|y|>R}|V(y,\tau)|^3\,dy=0 .
\]
Then, for every sequence \(\tau_k\in\R\),
\[
   \sup_k\|V(\tau_k)\|_{L^3(\R^3)}<\infty .
\]
\end{lemma}

\begin{proof}
Let \(M=\|V\|_{L^\infty(\R^3\times\R)}\).  Choose \(R\ge1\) such that
\[
   \sup_{\tau\in\R}\int_{|y|>R}|V(y,\tau)|^3\,dy\le1 .
\]
Then, for every \(k\),
\[
   \int_{\R^3}|V(y,\tau_k)|^3\,dy
   \le M^3 |B_R|
      + \sup_{\tau\in\R}\int_{|y|>R}|V(y,\tau)|^3\,dy
   \le M^3|B_R|+1 .
\]
This gives the claimed endpoint sequence bound.
\end{proof}

\subsection{Mildness after finite terminal-time restriction}

\begin{lemma}[Physical pullbacks are mild ancient solutions on finite intervals]
\label{p2:lem:physical-mild}
Let \(V\) be a bounded mild ancient solution of \eqref{p2:eq:centered}, and let
\(u\) be its physical pullback.  For every \(T<0\), the restriction of \(u\)
to \(\R^3\times(-\infty,T]\) is a bounded mild ancient solution of the
physical Navier--Stokes equations after translating the terminal time \(T\) to
\(0\).
\end{lemma}

\begin{proof}
By \Cref{p2:lem:physical-pullback}, \(u\) is a smooth bounded classical solution
on \(\R^3\times(-\infty,T]\).  Let \(s<t\le T\).  The classical equation on
the finite slab \([s,t]\) gives, for every solenoidal
\(\phi\in C_c^\infty(\R^3;\R^3)\), the weak Duhamel identity
\[
\begin{aligned}
   \langle u(t),\phi\rangle
   &=
   \langle e^{(t-s)\Delta}u(s),\phi\rangle \\
   &\quad
   -
   \int_s^t
   \left\langle
      e^{(t-r)\Delta}\mathbb P\divergence(u\otimes u)(r),
      \phi
   \right\rangle\,dr .
\end{aligned}
\]
Equivalently,
\[
   u(t)
   =
   e^{(t-s)\Delta}u(s)
   -
   \int_s^t
   e^{(t-r)\Delta}\mathbb P\divergence(u\otimes u)(r)\,dr
\]
in the weak-* topology of \(L^\infty(\R^3)\).  The heat--Stokes kernel bound
used in \Cref{p2:lem:centered-stokes-kernel} makes the integral finite on every
finite interval because \(u\) is bounded on
\(\R^3\times(-\infty,T]\).  Thus the restriction is a bounded mild ancient
solution.  The time translate \(u_T(x,\sigma)=u(x,\sigma+T)\),
\(\sigma\le0\), has the same property on \((-\infty,0]\).
\end{proof}

\subsection{\texorpdfstring{Endpoint \(L^3\) ancient Liouville theorem}{Endpoint L3 ancient Liouville theorem}}

\begin{theorem}[Albritton--Barker endpoint ancient Liouville theorem]
\label{p2:thm:AB-endpoint}
Let \(u\) be a bounded mild ancient solution of the three-dimensional
Navier--Stokes equations on \(\R^3\times(-\infty,0]\).  If there exists a
sequence \(s_k\downarrow-\infty\) such that
\[
   \sup_k\|u(s_k)\|_{L^3(\R^3)}<\infty,
\]
then \(u\equiv0\).
\end{theorem}

\begin{proof}
This is the endpoint ancient Liouville theorem of Albritton--Barker
\cite{AlbrittonBarker2019}.  The formulation above is the translated
terminal-time version needed here.
\end{proof}

\begin{theorem}[Tail-tight centered ancient solutions vanish]
\label{p2:thm:tight-liouville}
Let \(V\) be a bounded mild ancient solution of \eqref{p2:eq:centered}.  Assume
\[
   \lim_{R\to\infty}
   \sup_{\tau\in\R}\int_{|y|>R}|V(y,\tau)|^3\,dy=0 .
\]
Then \(V\equiv0\).
\end{theorem}

\begin{proof}
Let \(u\) be the physical pullback of \(V\).  Fix \(T<0\).  By
\Cref{p2:lem:physical-mild}, after translating the terminal time \(T\) to \(0\),
the restriction of \(u\) to \((-\infty,T]\) is a bounded mild ancient
solution.  Choose \(t_k=-e^k\).  For all sufficiently large \(k\), \(t_k<T\),
and \(\tau_k=-\log(-t_k)=-k\).  By
\Cref{p2:lem:tail-tight-sequence-L3},
\[
   \sup_k\|V(-k)\|_{L^3(\R^3)}<\infty .
\]
Hence \eqref{p2:eq:L3-pullback-invariance} gives
\[
   \|u(t_k)\|_{L^3(\R^3)}
   =
   \|V(-k)\|_{L^3(\R^3)}
\]
uniformly along the same sequence, after discarding the finitely many indices
for which \(t_k\ge T\).
Set \(u_T(x,\sigma)=u(x,\sigma+T)\) for \(\sigma\le0\) and
\(s_k=t_k-T\).  Then \(s_k\downarrow-\infty\) and
\[
   \|u_T(s_k)\|_{L^3(\R^3)}
   =
   \|u(t_k)\|_{L^3(\R^3)}
\]
is uniformly bounded.  Applying \Cref{p2:thm:AB-endpoint} to \(u_T\) gives
\(u_T\equiv0\) on \(\R^3\times(-\infty,0]\), equivalently \(u=0\) on
\(\R^3\times(-\infty,T]\).  Since \(T<0\) was arbitrary, \(u=0\) on
\(\R^3\times(-\infty,0)\).  The inverse change of variables then gives
\[
   V(y,\tau)=e^{-\tau/2}u(e^{-\tau/2}y,-e^{-\tau})=0
\]
for every \((y,\tau)\).
\end{proof}

\section{Tight-Class Liouville Criterion and Type I Reduction}
\label{p2:sec:tight-liouville-interface}

\subsection{Normalized tight collections}

\begin{theorem}[No nonempty normalized uniformly tight collection]
\label{p2:thm:no-tight-normalized-collection}
No uniformly tight normalized collection in the sense of
\Cref{p2:def:tight-normalized-collection} is nonempty.
\end{theorem}

\begin{proof}
Suppose \(A\) is such a collection and choose \(V\in A\).  By the tightness
modulus in item \textup{(3)} of \Cref{p2:def:tight-normalized-collection}, \(V\)
satisfies the tail condition in \Cref{p2:thm:tight-liouville}.  Hence
\Cref{p2:thm:tight-liouville} gives \(V\equiv0\).  This contradicts the local
velocity normalization in the same definition:
\[
   \sup_{\tau\in\R}
   \int_{B_{R_0}}|V(y,\tau)|^3\,dy\ge\eta_0>0 .
\]
Hence \(A\) is empty.
\end{proof}

\begin{proof}[Proof of \Cref{p2:thm:tight-liouville-intro}]
If \(A\) were nonempty, choose \(V\in A\).  The common tail modulus gives the
tail condition in \Cref{p2:thm:tight-liouville}, so \(V\equiv0\).  This contradicts
the fixed local nonvanishing normalization in \Cref{p2:thm:tight-liouville-intro}.
\end{proof}

\subsection{Use in the Type I reduction}

\begin{definition}[Tight-class Liouville consequence]
\label{p2:def:tight-class-input}
The tight-class Liouville consequence is the assertion that no normalized Seregin
ancient limit generated by an admissible Type I singularity and satisfying the
compact velocity nonvanishing normalization in \Cref{p1:prop:nonzero} is uniformly \(L^3\)-tight.
\end{definition}

\begin{lemma}[From \Cref{p1:thm:extraction} velocity nonvanishing to local velocity normalization]
\label{p2:lem:seregin-to-velocity-nonzero}
Let \(V\) be a smooth bounded normalized Seregin ancient limit.  Suppose that
the extraction theorem \Cref{p1:thm:extraction} gives a compact velocity lower bound: there are a
compact set \(K\Subset\mathbb R^3\times\mathbb R\) and \(\eta_I>0\) such that
\[
   \iint_K |V(y,\tau)|^3\,dy\,d\tau\ge\eta_I .
\]
Then there exist \(R_0>0\) and \(\eta_0>0\) such that
\[
   \sup_{\tau\in\R}\int_{B_{R_0}}|V(y,\tau)|^3\,dy\ge\eta_0 .
\]
\end{lemma}

\begin{proof}
Choose \(R_0>0\) and a compact interval \(I=[a,b]\), \(b>a\), such that
\(K\subset B_{R_0}\times I\).  Then
\[
   \eta_I
   \le \int_a^b\int_{B_{R_0}}|V(y,\tau)|^3\,dy\,d\tau
   \le (b-a)\sup_{\tau\in\R}
          \int_{B_{R_0}}|V(y,\tau)|^3\,dy .
\]
Thus the conclusion holds with \(\eta_0=\eta_I/(b-a)\).
\end{proof}

\begin{corollary}[Tight-class consequence for the Type I reduction]
\label{p2:cor:tight-input-for-paper-I}
\Cref{p2:def:tight-class-input} holds.
\end{corollary}

\begin{proof}
Suppose, contrary to \Cref{p2:def:tight-class-input}, that the Seregin
extraction produces a normalized ancient limit \(V\) in
the uniformly \(L^3\)-tight class from an admissible Type I singularity.  By the
extraction theorem \Cref{p1:thm:extraction}, \(V\) carries compact velocity nonvanishing:
there are \(K\Subset\mathbb R^3\times\mathbb R\) and \(\eta_I>0\) with
\[
   \iint_K |V(y,\tau)|^3\,dy\,d\tau\ge\eta_I .
\]
In particular \(V\not\equiv0\).  Since \(V\) is assumed uniformly
\(L^3\)-tight, it satisfies the tail condition
\[
   \lim_{R\to\infty}
   \sup_{\tau\in\R}\int_{|y|>R}|V(y,\tau)|^3\,dy=0 .
\]
\Cref{p2:thm:tight-liouville} then gives \(V\equiv0\), contradicting the compact
velocity lower bound above.  Hence no normalized Seregin ancient limit
generated by an admissible Type I singularity can belong to the uniformly
\(L^3\)-tight class.
\end{proof}

The tight-class consequence used in the Type I reduction is therefore proved
without the residual-class hypothesis.  The proof uses only boundedness, the tail-tightness condition
of the uniformly \(L^3\)-tight class, the local-in-time physical pullback,
and the endpoint ancient Liouville theorem of Albritton--Barker.



\clearpage
\section*{Part IV. Structured Ancient Solution Exclusions}

\subsection*{Overview}

We study bounded ancient solutions of the three-dimensional Navier--Stokes
equations arising from Type I blow-up arguments.  We state standard PDE
hypotheses under which such solutions are spatially constant; in the Type I
class, these constants vanish.  We also prove the perturbative weighted
vorticity Lyapunov rigidity used by the conditional Type I reduction.

For the unforced equations, known axisymmetric Liouville theorems cover the
following cases: absence of swirl, pointwise scale-invariant control, finite
\(L_t^\infty L_x^p\) control of the swirl variable \(\Gamma=ru_\theta\), and
axial periodicity with bounded \(\Gamma\).  When one of these hypotheses holds
for a nonzero Type I blow-up limit, the Galilean-trivial conclusions of the
cited Liouville theorems reduce to the zero solution and contradict the
nontriviality of the limit.

We also record the Lyapunov rigidity argument coming from Gallay--Wayne's
small-data decay theorem for the three-dimensional rescaled vorticity equation.
It yields a coercive Lyapunov functional forcing any ancient trajectory that
stays in the perturbative weighted-vorticity ball to be zero.  The section concludes
with a summary of the Liouville hypotheses and cited PDE results
used below.

\section{Introduction}

\subsection{Ancient solutions from Type I blow-up}

The Type I blow-up argument in \Cref{p1:thm:extraction} produces, after
parabolic rescaling around a Type I singular point, a nonzero ancient solution
\((U,P)\) of
\begin{equation}\label{p3:eq:NS}
   \partial_tU+(U\cdot\nabla)U+\nabla P=\Delta U,
   \qquad \nabla\cdot U=0,
   \qquad t<0,
\end{equation}
satisfying the Type I ancient bound
\begin{equation}\label{p3:eq:typeI}
   \sup_{t<0}\sqrt{-t}\,
   \|U(t)\|_{L^\infty(\mathbb R^3)}<\infty .
\end{equation}
Thus every result proving \(U\equiv0\) for an ancient solution satisfying
\eqref{p3:eq:typeI} has the same consequence in the Type I problem:
\[
   \text{PDE hypothesis}
   \quad\Longrightarrow\quad
   U\equiv0
   \quad\Longrightarrow\quad
   \text{contradiction with nonzero blow-up limit}.
\]

We treat hypotheses for which a specific Liouville mechanism is
available.  The cases considered are:
\[
   \begin{gathered}
   \text{axisymmetric Liouville hypotheses},\qquad
   \text{small weighted rescaled vorticity}.
   \end{gathered}
\]
We make no assertion for decaying or symmetry-restricted ancient solutions
outside the hypotheses listed below.  Each conclusion is tied to an explicit
PDE hypothesis and a specified analytic result.

For the Type I application in \Cref{p1:thm:extraction}, only those cases in which the stated
PDE hypothesis and the cited or proved Liouville theorem imply \(U\equiv0\)
are used as exclusions.  Ancient solutions outside these hypotheses are treated
separately.

\subsection{Axisymmetric Liouville hypotheses}

Several unforced axisymmetric cases are already covered by known Liouville
theorems.  Koch--Nadirashvili--Seregin--\v{S}ver\'ak cover the no-swirl case
and the pointwise scale-invariant condition \(|u|\le C/r\).  The
swirl-bearing finite-\(L^p\) condition
\[
   \Gamma=ru_\theta\in L_t^\infty L_x^p,\qquad 1\le p<\infty,
\]
is covered by Lei--Zhang--Zhao.  The \(z\)-periodic case with bounded
\(\Gamma\) is covered by Lei--Ren--Zhang.

The only additional step needed for Type I ancient limits is to convert
Galilean-triviality into zero in the blow-up frame.  We time-shift away from
the possible singular time \(t=0\), apply the bounded ancient Liouville
theorem, patch the resulting constants on backward intervals, and use
\eqref{p3:eq:typeI}: if \(U(x,t)\equiv c\), then
\[
   \sup_{t<0}\sqrt{-t}\,\|U(t)\|_\infty
   =
   \sup_{t<0}\sqrt{-t}\,|c|
\]
is finite only when \(c=0\).

\subsection{Lyapunov hypotheses}

The Lyapunov result below uses a forward Lyapunov functional in the
Gallay--Wayne perturbative weighted-vorticity regime.  If \(w\)
solves the rescaled three-dimensional vorticity equation and remains in the
Gallay--Wayne small ball of \(L^2_\sigma(m)\), \(m>7/2\), then the functional
\begin{equation}\label{p3:eq:intro-weighted-lyap}
   \mathcal L_m(w_0)
   =
   \sup_{\tau\ge0}e^{2\tau}\|\Phi_\tau w_0\|_m^2
\end{equation}
is coercive and contracts along the forward semiflow.  Sending the starting
time to \(-\infty\) then forces every ancient trajectory in the small ball to
be zero.

\subsection{Main results}

The first result states the unforced axisymmetric Liouville consequences
relevant to the Type I argument.

\begin{theorem}[Axisymmetric Liouville consequences for Type I ancient solutions]
\label{p3:thm:axisymmetric-liouville}
Let \(U\) be a smooth Type I ancient solution of \eqref{p3:eq:NS} satisfying
\eqref{p3:eq:typeI}.  If \(U\) satisfies any of the following axisymmetric
assumptions, then \(U\equiv0\):
\begin{enumerate}[label=\textup{(\roman*)}]
\item axisymmetric without swirl;
\item axisymmetric and satisfying a pointwise scale-invariant bound
\[
   |U(x,t)|\le \frac{C}{r};
\]
\item axisymmetric with
\[
   \Gamma=rU_\theta\in L_t^\infty L_x^p
   \bigl(\mathbb R^3\times(-\infty,0)\bigr),
   \qquad 1\le p<\infty;
\]
\item axisymmetric, periodic in the physical axial variable \(z\), and with
\[
   \Gamma=rU_\theta\in L^\infty
   \bigl(\mathbb R^3\times(-\infty,0)\bigr).
\]
\end{enumerate}
Consequently none of these assumptions can hold for a nonzero Type I blow-up
limit.
\end{theorem}

\begin{theorem}[Perturbative weighted-vorticity Lyapunov rigidity]
\label{p3:thm:weighted-vorticity-vanishing-intro}
Let \(m>\frac72\).  Let \(w(\tau)\) be an ancient mild solution of the
Gallay--Wayne rescaled vorticity equation
\begin{equation}\label{p3:eq:weighted-vorticity-intro}
   \partial_\tau w
   =
   \Lambda w-(v\cdot\nabla)w+(w\cdot\nabla)v,
   \qquad \nabla\cdot w=0,
\end{equation}
with \(v\) recovered by the three-dimensional Biot--Savart law.  Assume
\[
   \sup_{\tau\in\mathbb R}\|w(\tau)\|_{L^2_\sigma(m)}
   \le \rho_m,
\]
where \(\rho_m\) is the Gallay--Wayne small-data radius.  Then
\[
   w\equiv0.
\]
\end{theorem}

Here ``ancient mild'' means that every forward time-shift of \(w\) is the
Gallay--Wayne mild solution launched from its value at the starting time.

\begin{theorem}[Consequences for Type I blow-up limits]
\label{p3:thm:typeI-consequences}
Let \(U\) be a nonzero Type I blow-up limit.  Then \(U\) cannot satisfy any of
the axisymmetric assumptions in \Cref{p3:thm:axisymmetric-liouville}; nor can its
corresponding rescaled vorticity satisfy the hypotheses of
\Cref{p3:thm:weighted-vorticity-vanishing-intro}.
\end{theorem}

Only the explicit hypotheses stated above are used.

\subsection{Relation to the preceding parts}

\Cref{p1:thm:extraction} reduces Type I singularities to nonzero ancient
limits satisfying \eqref{p3:eq:typeI}.  \Cref{p2:thm:tight-liouville} treats
uniformly \(L^3\)-tight ancient solutions.  We prove the axisymmetric
and perturbative weighted-vorticity exclusions used by the Type I argument.
Known Liouville theorems rule out the axisymmetric cases listed
above, and the Lyapunov argument handles the perturbative weighted-vorticity
case.
\Cref{p3:thm:typeI-zero-hypotheses} does not depend on
\Cref{p2:thm:tight-liouville}.

\subsection{Summary of hypotheses}

\begin{center}
\begin{tabular}{p{0.36\textwidth}p{0.18\textwidth}p{0.34\textwidth}}
\hline
Hypothesis & Conclusion & Analytic ingredient \\
\hline
axisymmetric no swirl & \(U\equiv0\) & KNSS \\
pointwise \(|u|\le C/r\) & \(U\equiv0\) & KNSS \\
\(\Gamma\in L_t^\infty L_x^p\), \(1\le p<\infty\) & \(U\equiv0\) & LZZ \\
\(z\)-periodic bounded \(\Gamma\) & \(U\equiv0\) & LRZ \\
small weighted vorticity & \(w\equiv0\) & Gallay--Wayne \\
\hline
\end{tabular}
\end{center}

\subsection{Organization}

\Cref{p3:sec:typeI} introduces the Type I ancient setup, self-similar variables, the
invariance facts, and the time-shift lemma.
\Cref{p3:sec:axisymmetric-liouville-theorems} fixes axisymmetric notation and the
known Liouville theorems.  \Cref{p3:sec:axisymmetric-consequences} proves the
unforced axisymmetric Liouville consequences.
\Cref{p3:sec:weighted-vorticity} proves the Gallay--Wayne Lyapunov rigidity, and
\Cref{p3:sec:summary} summarizes the consequences.

\section{Ancient Type I limits and self-similar variables}
\label{p3:sec:typeI}

\subsection{Type I ancient solutions}

\begin{definition}[Type I ancient solution]
\label{p3:def:typeI-ancient}
A smooth mild solution \(U\) of \eqref{p3:eq:NS} on
\(\mathbb R^3\times(-\infty,0)\) is called Type I ancient if
\eqref{p3:eq:typeI} holds.  Mild is understood in the standard velocity
formulation \cite{Kato1984,Sohr2001}: for every \(-\infty<s<t<0\),
\begin{equation}\label{p3:eq:mild-formula}
   U(t)=e^{(t-s)\Delta}U(s)
   -\int_s^t e^{(t-\sigma)\Delta}\mathbb P\nabla\cdot
      (U(\sigma)\otimes U(\sigma))\,d\sigma
\end{equation}
in distributions, where \(\mathbb P\) is the Leray projector.
\end{definition}

\subsection{Nonzero Type I blow-up limits}

\begin{definition}[Nonzero Type I blow-up limit]
\label{p3:def:nonzero-typeI-limit}
A Type I ancient solution \(U\) is called a nonzero Type I blow-up limit if it
arises as the limit of parabolic rescalings around a Type I singular point and
the limiting solution is not identically zero.  The arguments below require
only the conclusion \(U\not\equiv0\).
\end{definition}

\begin{lemma}[Nontriviality of blow-up limits]
\label{p3:lem:nonzero-typeI-limit}
   A nonzero Type I blow-up limit is not identically zero.
\end{lemma}

\begin{proof}
This is part of the definition above.
\end{proof}

\subsection{Backward self-similar variables}

For \(t<0\), define
\begin{equation}\label{p3:eq:self-similar}
   V(y,\tau)=\sqrt{-t}\,U(x,t),\qquad
   Q(y,\tau)=(-t)P(x,t),\qquad
   y=\frac{x}{\sqrt{-t}},\qquad
   \tau=-\log(-t).
\end{equation}
Equivalently,
\[
   x=e^{-\tau/2}y,\qquad t=-e^{-\tau}.
\]

\begin{lemma}[Self-similar equation]\label{p3:lem:self-similar-equation}
Let \((U,P)\) be a smooth solution of \eqref{p3:eq:NS}, and define \((V,Q)\) by
\eqref{p3:eq:self-similar}.  Then \(V\) is defined on
\(\mathbb R^3\times\mathbb R\) and satisfies
\begin{equation}\label{p3:eq:centered}
   \partial_\tau V+(V\cdot\nabla)V+\nabla Q
   =
   \Delta V-\frac12(V+y\cdot\nabla V),
   \qquad \nabla\cdot V=0.
\end{equation}
Moreover \eqref{p3:eq:typeI} is equivalent to
\[
   \sup_{\tau\in\mathbb R}
   \|V(\tau)\|_{L^\infty(\mathbb R^3)}<\infty .
\]
\end{lemma}

\begin{proof}
Let \(a=\sqrt{-t}=e^{-\tau/2}\).  Then \(V=aU\), \(y=x/a\), and
\[
   \partial_t\tau=\frac1{-t}=\frac1{a^2},
   \qquad
   \partial_t y=\frac{y}{2a^2},
   \qquad
   \partial_t(a^{-1})=\frac1{2a^3}.
\]
Since \(U=a^{-1}V\), differentiating at fixed \(x\) gives
\[
   \partial_tU
   =
   a^{-3}
   \left(
      \partial_\tau V+\frac12 y\cdot\nabla V+\frac12V
   \right).
\]
Spatial derivatives scale as
\[
   \nabla_xU=a^{-2}\nabla_yV,\qquad
   \Delta_xU=a^{-3}\Delta_yV,
\]
and the pressure scaling gives
\[
   \nabla_xP=a^{-3}\nabla_yQ.
\]
Thus
\[
   (U\cdot\nabla_x)U=a^{-3}(V\cdot\nabla_y)V.
\]
Multiplying \eqref{p3:eq:NS} by \(a^3\) yields \eqref{p3:eq:centered}; the
divergence equation follows from
\(\nabla_x\cdot U=a^{-2}\nabla_y\cdot V\).  The spatial map
\(y\mapsto x=e^{-\tau/2}y\) is a bijection for each \(\tau\), and
\[
   \|V(\tau)\|_\infty
   =
   \sqrt{-t}\,\|U(t)\|_\infty .
\]
Taking suprema over \(t<0\), equivalently \(\tau\in\mathbb R\), gives the
last assertion.
\end{proof}

\subsection{Preservation of axisymmetry and swirl variables}

For axisymmetric flows we write
\[
   x=(r\cos\theta,r\sin\theta,z),\qquad
   U=U_re_r+U_\theta e_\theta+U_ze_z,
\]
and in centered variables
\[
   y=(\rho\cos\theta,\rho\sin\theta,Z),\qquad
   V=V_\rho e_\rho+V_\theta e_\theta+V_Ze_Z.
\]
The scalar swirl variables are
\[
   \Gamma_U(r,z,t)=rU_\theta(r,z,t),\qquad
   \Gamma_V(\rho,Z,\tau)=\rho V_\theta(\rho,Z,\tau).
\]

\begin{lemma}[Preservation of axisymmetric structure]
\label{p3:lem:structure-preserved}
Axisymmetry and absence of swirl are preserved under
\eqref{p3:eq:self-similar}.  If \(U\) is axisymmetric, then
\[
   \Gamma_V(\rho,Z,\tau)=\Gamma_U(r,z,t),
   \qquad
   (r,z)=\sqrt{-t}\,(\rho,Z).
\]
Fixed physical \(z\)-periodicity is preserved by physical time shifts, but
under the centered change of variables it becomes a \(\tau\)-dependent period
in the \(Z\)-coordinate.
\end{lemma}

\begin{proof}
The self-similar map is a scalar spatial dilation and a scalar multiplication
of the velocity, so it commutes with rotations around the symmetry axis.
The cylindrical basis is unchanged by radial dilation, and
\[
   V_\theta(\rho,Z,\tau)
   =
   \sqrt{-t}\,U_\theta(r,z,t),
   \qquad
   r=\sqrt{-t}\rho,\quad z=\sqrt{-t}Z .
\]
Thus \(U_\theta\equiv0\) if and only if \(V_\theta\equiv0\), and
\[
   \Gamma_V=\rho V_\theta
   =
   \rho\sqrt{-t}\,U_\theta
   =
   rU_\theta
   =
   \Gamma_U .
\]
If \(U\) has physical axial period \(L\), then \(V\) has period
\(L/\sqrt{-t}=Le^{\tau/2}\) in \(Z\), not a fixed period.
\end{proof}

\subsection{Time shifts away from the singular time}

\begin{lemma}[Time shifts away from the singular time]
\label{p3:lem:time-shift}
Let \(U\) be a smooth Type I ancient solution and fix \(T<0\).  Then
\[
   W_T(x,s)=U(x,T+s),\qquad \Pi_T(x,s)=P(x,T+s),\qquad s<0,
\]
is a smooth bounded mild ancient solution of \eqref{p3:eq:NS} on
\(\mathbb R^3\times(-\infty,0)\).  Moreover the axisymmetric, no-swirl,
\(L^p\)-swirl, bounded-\(\Gamma\), and physical \(z\)-periodic properties are
preserved.
\end{lemma}

\begin{proof}
Time translation preserves the differential equation because
\(\partial_sW_T(x,s)=\partial_tU(x,T+s)\) and all spatial derivatives are
unchanged.  For the mild formulation, apply \eqref{p3:eq:mild-formula} to \(U\)
between the physical times \(T+s<T+t<0\):
\[
   U(T+t)=e^{(t-s)\Delta}U(T+s)
   -\int_{T+s}^{T+t}e^{(T+t-\eta)\Delta}\mathbb P\nabla\cdot
      (U(\eta)\otimes U(\eta))\,d\eta .
\]
With \(\eta=T+\sigma\), this becomes the mild formula for \(W_T\) on the
interval \((s,t)\).  For \(s<0\), \(T+s<T<0\), so the Type I bound gives
\[
   \|W_T(s)\|_\infty
   =
   \|U(T+s)\|_\infty
   \le
   \frac1{\sqrt{-T}}
   \sup_{\eta<0}\sqrt{-\eta}\,\|U(\eta)\|_\infty .
\]
The invariance assertions follow because only the time variable has been
translated.  In particular
\[
   \Gamma_{W_T}(r,z,s)=\Gamma_U(r,z,T+s),
\]
so \(L_s^\infty L_x^p\) and \(L^\infty\) bounds for \(\Gamma\) restrict to the
shifted solution, and a fixed physical axial period remains the same.
\end{proof}

\subsection{PDE hypotheses used below}

The arguments below are stated directly in terms of hypotheses on ancient
solutions.  The axisymmetric hypotheses are:
\begin{enumerate}[label=\textup{(\roman*)}]
\item axisymmetry and \(U_\theta\equiv0\);
\item axisymmetry and the pointwise bound \(|U(x,t)|\le C/r\) for \(r>0\);
\item axisymmetry and
\[
   \Gamma=rU_\theta\in L_t^\infty L_x^p
   \bigl(\mathbb R^3\times(-\infty,0)\bigr),
   \qquad 1\le p<\infty;
\]
\item axisymmetry, periodicity in the physical axial variable \(z\), and
\(\Gamma\in L^\infty\).
\end{enumerate}
The additional hypothesis used later is the Gallay--Wayne small
weighted-vorticity condition.

\section{Axisymmetric notation and known Liouville theorems}
\label{p3:sec:axisymmetric-liouville-theorems}

\subsection{Cylindrical variables and the swirl scalar}

For an axisymmetric solution \(u\), write
\[
   u=u_re_r+u_\theta e_\theta+u_ze_z,\qquad
   \Gamma=ru_\theta,\qquad
   b=u_re_r+u_ze_z .
\]
When \(\Gamma\) is placed in \(L^p(\mathbb R^3)\), it is regarded as the
corresponding function of \(x=(r\cos\theta,r\sin\theta,z)\), independent of
\(\theta\):
\[
   \|\Gamma(t)\|_{L^p(\mathbb R^3)}^p
   =
   2\pi\int_{\mathbb R}\int_0^\infty
      |\Gamma(r,z,t)|^p r\,dr\,dz,\qquad 1\le p<\infty .
\]

\subsection{Constant and Galilean-trivial solutions}

\begin{definition}[Galilean-trivial]
A velocity field is Galilean-trivial if, after a Galilean transform, it is a
constant vector field.  In the Type I blow-up frame below we use only the
following consequence: when a cited Liouville theorem gives a constant
velocity field in that frame, the Type I bound forces the constant to be zero.
\end{definition}

\begin{lemma}[Spatially constant mild solutions]\label{p3:lem:spatial-constant-mild}
Let \(W\) be a smooth mild solution of \eqref{p3:eq:NS} on
\(\mathbb R^3\times(-\infty,0)\).  If \(W(x,t)=b(t)\) for all \(x\) and \(t\),
then \(b\) is constant in time.
\end{lemma}

\begin{proof}
Insert \(W(x,t)=b(t)\) into \eqref{p3:eq:mild-formula}.  The heat semigroup
leaves constant vector fields fixed, and
\(\nabla\cdot(b(\sigma)\otimes b(\sigma))=0\).  Thus
\(b(t)=b(s)\) for every \(s<t<0\).
\end{proof}

\begin{lemma}[Patching constants on backward intervals]
\label{p3:lem:patch-constants}
Let \(U:\mathbb R^3\times(-\infty,0)\to\mathbb R^3\) be a function.  Assume
that for every \(T<0\) there exists \(c_T\in\mathbb R^3\) such that
\[
   U(x,t)=c_T\qquad(x\in\mathbb R^3,\ t<T).
\]
Then there is a single \(c\in\mathbb R^3\) such that
\[
   U(x,t)=c\qquad(x\in\mathbb R^3,\ t<0).
\]
\end{lemma}

\begin{proof}
If \(T_1<T_2<0\), the two constants agree on the nonempty interval
\((-\infty,T_1)\).  Hence all \(c_T\) are equal.  For any \(t<0\), choose
\(T\) with \(t<T<0\).
\end{proof}

\begin{lemma}[Constant Type I ancient solutions vanish]
\label{p3:lem:galilean-typeI-zero}
Let \(U\) be a Type I ancient solution.  If \(U(x,t)\equiv c\), then
\(c=0\).
\end{lemma}

\begin{proof}
The Type I bound gives
\[
   \sqrt{-t}\,|c|
   \le
   \sup_{s<0}\sqrt{-s}\,\|U(s)\|_\infty<\infty
   \qquad(t<0).
\]
Letting \(t\to-\infty\) gives \(c=0\).
\end{proof}

\subsection{Known Liouville theorems}

We use the following Liouville theorems in the stated forms.

\begin{theorem}[KNSS no-swirl Liouville theorem]
\label{p3:thm:knss-noswirl}
Every bounded ancient mild axisymmetric no-swirl solution of the unforced
Navier--Stokes equations is Galilean-trivial.  In the form used below, the
solution is a spatially constant axial drift.
\end{theorem}

\begin{proof}
This is the no-swirl Liouville theorem of
Koch--Nadirashvili--Seregin--\v{S}ver\'ak \cite{KNSS2009}, stated in the form
needed for the Type I vanishing argument.
\end{proof}

\begin{theorem}[KNSS pointwise scale-invariant condition]
\label{p3:thm:knss-pointwise}
Every bounded ancient mild axisymmetric solution satisfying
\[
   |u(x,t)|\le \frac{C}{r}
   \qquad(r>0,\ t<0)
\]
is Galilean-trivial.  In the form used below, it is identically zero.
\end{theorem}

\begin{proof}
This is the pointwise Liouville theorem of \cite{KNSS2009}, stated in the form
used here.
\end{proof}

\begin{theorem}[Lei--Zhang--Zhao]
\label{p3:thm:lzz}
Let \(u\) be a bounded ancient mild axisymmetric solution.  If
\[
   \Gamma=ru_\theta\in L_t^\infty L_x^p
   \bigl(\mathbb R^3\times(-\infty,0)\bigr),
   \qquad 1\le p<\infty,
\]
then \(u\) is a constant vector field.
\end{theorem}

\begin{proof}
This is the ancient axisymmetric \(L^p\)-swirl Liouville theorem of
Lei--Zhang--Zhao \cite{LZZ}.
\end{proof}

\begin{theorem}[Lei--Ren--Zhang]
\label{p3:thm:lrz}
Let \(u\) be a bounded ancient mild axisymmetric solution, periodic in \(z\),
with
\[
   \Gamma=ru_\theta\in L^\infty
   \bigl(\mathbb R^3\times(-\infty,0)\bigr).
\]
Then \(u\) is a constant axial vector field.
\end{theorem}

\begin{proof}
This is the periodic bounded-swirl theorem of Lei--Ren--Zhang
\cite{LRZ}.  Their statement is on \(\mathbb R^2\times\mathbb T^1\); a
physical \(z\)-periodic whole-space solution descends to the quotient, and a
Navier--Stokes scaling normalizes the period.
\end{proof}

\section{Unforced axisymmetric Liouville consequences}
\label{p3:sec:axisymmetric-consequences}

\subsection{No swirl}

\begin{proposition}[No-swirl Type I vanishing]\label{p3:prop:noswirl-typeI}
Let \(U\) be a smooth Type I ancient solution.  Assume \(U\) is axisymmetric
and \(U_\theta\equiv0\).  Then \(U\equiv0\).
\end{proposition}

\begin{proof}
Fix \(T<0\) and set \(W_T(x,s)=U(x,T+s)\).  By \Cref{p3:lem:time-shift},
\(W_T\) is a bounded ancient mild axisymmetric no-swirl solution.  Applying
\Cref{p3:thm:knss-noswirl} gives that \(W_T\) is a spatially constant axial
drift, say \(W_T(x,s)=b_T(s)e_z\).  By
\Cref{p3:lem:spatial-constant-mild}, \(b_T(s)\) is constant in \(s\).  Thus
\(U(x,t)=c_T\) for \(t<T\).  Since \(T<0\) is arbitrary,
\Cref{p3:lem:patch-constants} gives \(U(x,t)=c\) for all \(t<0\).  Finally,
\Cref{p3:lem:galilean-typeI-zero} gives \(c=0\).
\end{proof}

\subsection{Pointwise scale-invariant control}

\begin{proposition}[Pointwise \(1/r\) Type I vanishing]\label{p3:prop:pointwise-typeI}
Let \(U\) be a smooth Type I ancient solution.  Assume \(U\) is axisymmetric
and satisfies
\[
   |U(x,t)|\le \frac{C}{r}
   \qquad(r>0,\ t<0).
\]
Then \(U\equiv0\).
\end{proposition}

\begin{proof}
Fix \(T<0\).  The shifted solution \(W_T(x,s)=U(x,T+s)\) is bounded, ancient,
axisymmetric, and satisfies the same pointwise bound \(|W_T|\le C/r\).
\Cref{p3:thm:knss-pointwise} gives \(W_T\equiv0\).  Hence \(U\equiv0\) on
\((-\infty,T)\).  Since \(T<0\) is arbitrary, \(U\equiv0\).
\end{proof}

\subsection{\texorpdfstring{Finite \(L^p\) control of \(\Gamma\)}
{Finite Lp control of Gamma}}

\begin{proposition}[Finite-\(L^p\) swirl Type I vanishing]
\label{p3:prop:lp-swirl-typeI}
Let \(U\) be a smooth Type I ancient solution.  Assume \(U\) is axisymmetric
and, for some \(1\le p<\infty\),
\[
   \Gamma=rU_\theta\in L_t^\infty L_x^p
   \bigl(\mathbb R^3\times(-\infty,0)\bigr).
\]
Then \(U\equiv0\).
\end{proposition}

\begin{proof}
Fix \(T<0\).  By \Cref{p3:lem:time-shift}, \(W_T(x,s)=U(x,T+s)\) is a bounded
ancient mild axisymmetric solution and
\[
   \Gamma_{W_T}\in L_s^\infty L_x^p
   \bigl(\mathbb R^3\times(-\infty,0)\bigr).
\]
\Cref{p3:thm:lzz} gives that \(W_T\) is a constant vector field.  Patch these
constants in \(T\) by \Cref{p3:lem:patch-constants} and then use
\Cref{p3:lem:galilean-typeI-zero}.
\end{proof}

\subsection{Periodic bounded swirl}

\begin{proposition}[Periodic bounded-swirl Type I vanishing]
\label{p3:prop:periodic-swirl-typeI}
Let \(U\) be a smooth Type I ancient solution.  Assume \(U\) is axisymmetric,
periodic in the physical axial variable \(z\), and satisfies
\[
   \Gamma=rU_\theta\in L^\infty
   \bigl(\mathbb R^3\times(-\infty,0)\bigr).
\]
Then \(U\equiv0\).
\end{proposition}

\begin{proof}
Fix \(T<0\).  By \Cref{p3:lem:time-shift}, the shifted solution \(W_T\) is
bounded, ancient, axisymmetric, physically \(z\)-periodic, and has bounded
\(\Gamma\).  \Cref{p3:thm:lrz} gives that \(W_T\) is a constant axial vector
field.  The constants patch as \(T\) varies, and the Type I bound eliminates
the resulting constant by \Cref{p3:lem:galilean-typeI-zero}.
\end{proof}

\subsection{Proof of the axisymmetric Liouville theorem}

\begin{proof}[Proof of \Cref{p3:thm:axisymmetric-liouville}]
The four cases are precisely the hypotheses of
\Cref{p3:prop:noswirl-typeI,p3:prop:pointwise-typeI,p3:prop:lp-swirl-typeI,p3:prop:periodic-swirl-typeI}.
Each proposition gives \(U\equiv0\).  This contradicts
\Cref{p3:lem:nonzero-typeI-limit} when \(U\) is a nonzero Type I blow-up limit.
\end{proof}

\section{Perturbative weighted-vorticity Lyapunov rigidity}
\label{p3:sec:weighted-vorticity}

\subsection{The rescaled vorticity equation}

In this section \(\tau\) is the forward self-similar time used in the
Gallay--Wayne vorticity formulation.  The ancient trajectories are indexed by
all \(\tau\in\mathbb R\), and the forward semiflow acts by increasing
\(\tau\).

Let \(w=\nabla\times v\).  The rescaled vorticity equation is
\begin{equation}\label{p3:eq:weighted-vorticity}
   \partial_\tau w
   =
   \Lambda w-(v\cdot\nabla)w+(w\cdot\nabla)v,
   \qquad \nabla\cdot w=0,
\end{equation}
where
\begin{equation}\label{p3:eq:lambda}
   \Lambda w=\Delta w+\frac12\xi\cdot\nabla w+w,
\end{equation}
and \(v\) is recovered from \(w\) by the three-dimensional Biot--Savart law.
For \(m\ge0\), set
\[
   L^2_\sigma(m)
   =
   \left\{
      f\in L^2(\mathbb R^3;\mathbb R^3):
      \nabla\cdot f=0,\quad \|f\|_m<\infty
   \right\},
\]
where
\[
   \|f\|_m^2
   =
   \int_{\mathbb R^3}(1+|\xi|)^{2m}|f(\xi)|^2\,d\xi .
\]

\subsection{Gallay--Wayne small-data decay}

\begin{theorem}[Gallay--Wayne small-data decay]\label{p3:thm:gallay}
Let \(0<\mu\le1\) and \(m>2\mu+\frac12\).  There exist constants
\(\rho_m>0\) and \(C_m\ge1\) such that if \(\|w_0\|_m\le\rho_m\), then
\eqref{p3:eq:weighted-vorticity} has a unique global mild solution
\(w(\tau)=\Phi_\tau w_0\) in \(C([0,\infty);L^2_\sigma(m))\) and
\[
   \|\Phi_\tau w_0\|_m
   \le
   C_m e^{-\mu\tau}\|w_0\|_m,\qquad \tau\ge0.
\]
In particular, if \(m>\frac72\), one may use \(\mu=1\).
\end{theorem}

\begin{proof}
This is the Gallay--Wayne weighted small-data theorem for the rescaled
three-dimensional vorticity equation \cite{GallayWayne2002}.  The only
conclusions used below are existence, uniqueness, the semiflow property, and
the displayed exponential decay.
\end{proof}

\subsection{The Lyapunov functional}

\begin{definition}[Perturbative weighted-vorticity condition]
\label{p3:def:weighted-vorticity-condition}
Fix \(m>\frac72\), and let \(\rho_m\) be the radius in \Cref{p3:thm:gallay}.
An ancient vorticity trajectory satisfies the perturbative
weighted-vorticity condition with exponent \(m\) if it is a function
\[
   w\in C(\mathbb R;L^2_\sigma(m))
\]
that solves \eqref{p3:eq:weighted-vorticity} for all \(\tau\in\mathbb R\) and
satisfies
\[
   \sup_{\tau\in\mathbb R}\|w(\tau)\|_m\le\rho_m .
\]
Solving for all \(\tau\) means that every time shift is compatible with the
Gallay--Wayne forward mild semiflow.
\end{definition}

\begin{proposition}[Semiflow Lyapunov functional]\label{p3:prop:weighted-lyap}
Fix \(m>\frac72\).  For \(\|w_0\|_m\le\rho_m\), define
\begin{equation}\label{p3:eq:weighted-lyap}
   \mathcal L_m(w_0)
   =
   \sup_{\tau\ge0}e^{2\tau}\|\Phi_\tau w_0\|_m^2 .
\end{equation}
Then \(\mathcal L_m\) is well-defined and satisfies
\begin{equation}\label{p3:eq:weighted-coercive}
   \|w_0\|_m^2
   \le
   \mathcal L_m(w_0)
   \le
   C_m^2\|w_0\|_m^2 .
\end{equation}
Moreover, for every \(s\ge0\),
\begin{equation}\label{p3:eq:weighted-contract}
   \mathcal L_m(\Phi_s w_0)
   \le
   e^{-2s}\mathcal L_m(w_0).
\end{equation}
\end{proposition}

\begin{proof}
The Gallay--Wayne estimate with \(\mu=1\) gives
\[
   e^{2\tau}\|\Phi_\tau w_0\|_m^2
   \le C_m^2\|w_0\|_m^2,
\]
which proves finiteness and the upper bound.  The lower bound is the value at
\(\tau=0\).  By uniqueness,
\[
   \Phi_\tau(\Phi_s w_0)=\Phi_{\tau+s}w_0 .
\]
Hence
\[
   \mathcal L_m(\Phi_s w_0)
   =
   \sup_{\tau\ge0}e^{2\tau}\|\Phi_{\tau+s}w_0\|_m^2
   =
   e^{-2s}\sup_{\sigma\ge s}e^{2\sigma}\|\Phi_\sigma w_0\|_m^2
   \le e^{-2s}\mathcal L_m(w_0).
\]
\end{proof}

\subsection{Ancient rigidity in the small weighted ball}

\begin{proof}[Proof of \Cref{p3:thm:weighted-vorticity-vanishing-intro}]
Let \(w\) be an ancient solution in the small ball and fix
\(\tau_0\in\mathbb R\).  For \(s<\tau_0\), the shifted function
\[
   W_s(\theta)=w(s+\theta),\qquad \theta\ge0,
\]
is the Gallay--Wayne mild solution with initial datum \(w(s)\).  The small-ball
hypothesis gives \(\|w(s)\|_m\le\rho_m\), so the forward solution is in the
domain of \(\Phi_\theta\), and uniqueness gives
\[
   W_s(\theta)=\Phi_\theta w(s),\qquad \theta\ge0.
\]
Choosing \(\theta=\tau_0-s\) gives
\[
   w(\tau_0)=\Phi_{\tau_0-s}w(s).
\]
Since \(\|w(\tau_0)\|_m\le\rho_m\), the Lyapunov functional is also defined at
\(w(\tau_0)\).  Applying \Cref{p3:prop:weighted-lyap} with forward time
\(\tau_0-s\),
\[
   \mathcal L_m(w(\tau_0))
   \le
   e^{-2(\tau_0-s)}\mathcal L_m(w(s))
   \le
   C_m^2\rho_m^2e^{-2(\tau_0-s)} .
\]
Letting \(s\to-\infty\) gives \(\mathcal L_m(w(\tau_0))=0\).  The coercive
lower bound \eqref{p3:eq:weighted-coercive} gives \(w(\tau_0)=0\).  Since
\(\tau_0\) was arbitrary, \(w\equiv0\).
\end{proof}

\section{Summary of PDE consequences}
\label{p3:sec:summary}

\subsection{Zero conclusions}

\begin{theorem}[Zero conclusions for Type I blow-up limits]
\label{p3:thm:typeI-zero-hypotheses}
The following hypotheses cannot hold for a nonzero Type I blow-up limit:
\begin{enumerate}[label=\textup{(\roman*)}]
\item axisymmetry and \(U_\theta\equiv0\);
\item axisymmetry and the pointwise bound \(|U(x,t)|\le C/r\);
\item axisymmetry and finite-\(L^p\) control of the swirl scalar,
\[
   \Gamma\in L_t^\infty L_x^p,\qquad 1\le p<\infty;
\]
\item axisymmetry, \(z\)-periodicity, and bounded \(\Gamma\);
\item the Gallay--Wayne small weighted-vorticity condition, whenever the
velocity is recovered from the corresponding vorticity by the Biot--Savart law.
\end{enumerate}
\end{theorem}

\begin{proof}
Items (i)--(iv) are \Cref{p3:thm:axisymmetric-liouville}.  Item (v) is
\Cref{p3:thm:weighted-vorticity-vanishing-intro}; since the velocity in that
case is recovered from \(w\) by the Biot--Savart law, \(w\equiv0\) gives the
zero velocity.  In each Type I case, the zero conclusion contradicts
\Cref{p3:lem:nonzero-typeI-limit}.
\end{proof}

\subsection{Consequences for Type I blow-up arguments}

\begin{proof}[Proof of \Cref{p3:thm:typeI-consequences}]
The axisymmetric cases are ruled out by \Cref{p3:thm:axisymmetric-liouville}.  The
perturbative weighted-vorticity case is ruled out by
\Cref{p3:thm:weighted-vorticity-vanishing-intro} once the corresponding rescaled
vorticity satisfies the small weighted condition; the Biot--Savart
reconstruction then gives zero velocity.  These zero conclusions contradict
the nontriviality of a Type I blow-up limit.
\end{proof}

\begin{theorem}[Combined zero conclusion]
\label{p3:thm:stated-hypotheses-zero}
Let \(U\) be a Type I ancient solution obtained as a nonzero Type I blow-up
limit.  If \(U\) satisfies one of the hypotheses in
\Cref{p3:thm:typeI-zero-hypotheses}, then \(U\equiv0\), a contradiction.
\end{theorem}

\begin{proof}
This is \Cref{p3:thm:typeI-zero-hypotheses}.
\end{proof}

\subsection{Summary table}

\begin{center}
\begin{tabular}{p{0.36\textwidth}p{0.18\textwidth}p{0.34\textwidth}}
\hline
Hypothesis & Conclusion & Analytic ingredient \\
\hline
no swirl & \(U\equiv0\) & KNSS plus Type I constant removal \\
pointwise \(1/r\) & \(U\equiv0\) & KNSS \\
\(\Gamma\in L_t^\infty L_x^p\) & \(U\equiv0\) & LZZ plus Type I constant removal \\
periodic bounded \(\Gamma\) & \(U\equiv0\) & LRZ plus Type I constant removal \\
small weighted vorticity & \(w\equiv0\) & Gallay--Wayne Lyapunov functional \\
\hline
\end{tabular}
\end{center}

The no-swirl, pointwise scale-invariant, finite-\(L^p\)-swirl, and periodic
bounded-swirl cases are covered by known axisymmetric Liouville theorems.  The
perturbative weighted-vorticity case is covered by a Lyapunov functional
derived from Gallay--Wayne decay.
Thus every zero conclusion used in the Type I application rests on an explicit
PDE hypothesis stated in this paper.



\clearpage
\begin{thebibliography}{99}

\bibitem{AlbrittonBarker2019}
D.~Albritton and T.~Barker,
\emph{On local Type I singularities of the Navier--Stokes equations and
Liouville theorems},
J. Math. Fluid Mech. \textbf{21} (2019), Article 43.

\bibitem{CKN1982}
L.~Caffarelli, R.~Kohn, and L.~Nirenberg,
\emph{Partial regularity of suitable weak solutions of the Navier--Stokes
 equations},
Comm. Pure Appl. Math. \textbf{35} (1982), no. 6, 771--831,
doi:\ \nolinkurl{10.1002/cpa.3160350604}.

\bibitem{ESS2003}
L.~Escauriaza, G.~A.~Seregin, and V.~{\v S}verak,
\emph{\(L_{3,\infty}\)-solutions of the Navier--Stokes equations and
 backward uniqueness},
Russian Math. Surveys \textbf{58} (2003), no. 2, 211--250,
doi:\ \nolinkurl{10.1070/RM2003v058n02ABEH000609}.

\bibitem{Giga1986}
Y.~Giga,
\emph{Solutions for semilinear parabolic equations in \(L^p\) and regularity
 of weak solutions of the Navier--Stokes system},
J. Differential Equations \textbf{62} (1986), no. 2, 186--212,
doi:\ \nolinkurl{10.1016/0022-0396(86)90096-3}.

\bibitem{Leray1934}
J.~Leray,
\emph{Sur le mouvement d'un liquide visqueux emplissant l'espace},
Acta Math. \textbf{63} (1934), no. 1, 193--248,
doi:\ \nolinkurl{10.1007/BF02547354}.

\bibitem{Lin1998}
F.-H.~Lin,
\emph{A new proof of the Caffarelli--Kohn--Nirenberg theorem},
Comm. Pure Appl. Math. \textbf{51} (1998), no. 3, 241--257,
doi:\ \nolinkurl{10.1002/(SICI)1097-0312(199803)51:3<241::AID-CPA2>3.0.CO;2-A}.

\bibitem{Scheffer1976}
V.~Scheffer,
\emph{Partial regularity of solutions to the Navier--Stokes equations},
Pacific J. Math. \textbf{66} (1976), no. 2, 535--552,
doi:\ \nolinkurl{10.2140/pjm.1976.66.535}.

\bibitem{Seregin2012}
G.~Seregin,
\emph{A certain necessary condition of potential blow up for Navier--Stokes
 equations},
Comm. Math. Phys. \textbf{312} (2012), no. 3, 833--845,
doi:\ \nolinkurl{10.1007/s00220-011-1391-x}.

\bibitem{Serrin1962}
J.~Serrin,
\emph{On the interior regularity of weak solutions of the Navier--Stokes
 equations},
Arch. Rational Mech. Anal. \textbf{9} (1962), 187--195,
doi:\ \nolinkurl{10.1007/BF00253344}.

\bibitem{Stein1970}
E.~M. Stein,
\emph{Singular Integrals and Differentiability Properties of Functions},
Princeton Mathematical Series, No.~30, Princeton University Press, 1970.

\bibitem{Aubin1963}
J.-P. Aubin,
\emph{Un theoreme de compacite},
C. R. Acad. Sci. Paris \textbf{256} (1963), 5042--5044.

\bibitem{Kato1984}
T. Kato,
\emph{Strong \(L^p\)-solutions of the Navier--Stokes equation in
\(\mathbb R^m\), with applications to weak solutions},
Math. Z. \textbf{187} (1984), no. 4, 471--480,
doi:\url{https://doi.org/10.1007/BF01174182}.

\bibitem{KNSS2009}
G. Koch, N. Nadirashvili, G. Seregin, and V. \v{S}verak,
\emph{Liouville theorems for the Navier--Stokes equations and applications},
Acta Math. \textbf{203} (2009), no. 1, 83--105.

\bibitem{LionsMagenes}
J.-L. Lions and E. Magenes,
\emph{Non-homogeneous boundary value problems and applications. Vol. I},
Springer-Verlag, 1972.

\bibitem{LRZ2022}
Z. Lei, X. Ren, and Q. S. Zhang,
\emph{A Liouville theorem for the axially symmetric Navier--Stokes equations},
J. Funct. Anal. \textbf{283} (2022), no. 3, 109522.

\bibitem{LZZ2017}
Z. Lei, Q. Zhang, and N. Zhao,
\emph{Improved Liouville theorems for axially symmetric Navier--Stokes equations},
Sci. Sin. Math. \textbf{47} (2017), no. 10, 1183--1198.

\bibitem{NRS1996}
J. Necas, M. Ruzicka, and V. \v{S}verak,
\emph{On Leray's self-similar solutions of the Navier--Stokes equations},
Acta Math. \textbf{176} (1996), no. 2, 283--294.

\bibitem{Simon1987}
J. Simon,
\emph{Compact sets in the space \(L^p(0,T;B)\)},
Ann. Mat. Pura Appl. (4) \textbf{146} (1987), 65--96,
doi:\url{https://doi.org/10.1007/BF01762360}.

\bibitem{Sohr2001}
H. Sohr,
\emph{The Navier--Stokes equations. An elementary functional analytic approach},
Birkhauser, 2001.

\bibitem{Billingsley1999}
P. Billingsley,
\emph{Convergence of Probability Measures}, second ed.,
Wiley, New York, 1999.

\bibitem{FabesJonesRiviere1972}
E. B. Fabes, B. F. Jones, and N. M. Riviere,
\emph{The initial value problem for the Navier--Stokes equations with data in
\(L^p\)},
Arch. Rational Mech. Anal. 45 (1972), 222--240.

\bibitem{GigaMiyakawa1985}
Y. Giga and T. Miyakawa,
\emph{Solutions in \(L_r\) of the Navier--Stokes initial value problem},
Arch. Rational Mech. Anal. 89 (1985), 267--281.

\bibitem{KryloffBogoliouboff1937}
N. Kryloff and N. Bogoliouboff,
\emph{La theorie generale de la mesure dans son application a l'etude des
systemes dynamiques de la mecanique non lineaire},
Ann. of Math. (2) 38 (1937), 65--113.

\bibitem{Ladyzhenskaya1968}
O. A. Ladyzhenskaya, V. A. Solonnikov, and N. N. Ural'tseva,
\emph{Linear and Quasilinear Equations of Parabolic Type},
American Mathematical Society, Providence, 1968.

\bibitem{CSTY2008}
C.-C. Chen, R. M. Strain, T.-P. Tsai, and H.-T. Yau,
\emph{Lower bound on the blow-up rate of the axisymmetric Navier--Stokes
equations},
Int. Math. Res. Not. IMRN \textbf{2008}, Art. ID rnn016, 31 pp.,
doi:10.1093/imrn/rnn016.

\bibitem{CSTY2009}
C.-C. Chen, R. M. Strain, T.-P. Tsai, and H.-T. Yau,
\emph{Lower bounds on the blow-up rate of the axisymmetric Navier--Stokes
equations II},
Comm. Partial Differential Equations \textbf{34} (2009), no. 3, 203--232,
doi:10.1080/03605300902793956.

\bibitem{GallayWayne2002}
T. Gallay and C. E. Wayne,
\emph{Long-time asymptotics of the Navier--Stokes and vorticity equations on
\(\mathbb R^3\)},
Philos. Trans. Roy. Soc. London Ser. A \textbf{360} (2002), no. 1799,
2155--2188, doi:10.1098/rsta.2002.1068.

\bibitem{LeiZhang2011}
Z. Lei and Q. S. Zhang,
\emph{A Liouville theorem for the axially-symmetric Navier--Stokes equations},
J. Funct. Anal. \textbf{261} (2011), no. 8, 2323--2345,
doi:10.1016/j.jfa.2011.06.016.

\bibitem{LRZ}
Z. Lei, X. Ren, and Q. S. Zhang,
\emph{A Liouville theorem for axially symmetric Navier--Stokes equations on
\(\mathbb R^2\times\mathbb T^1\)},
Math. Ann. \textbf{383} (2022), no. 1--2, 415--431,
doi:10.1007/s00208-020-02128-9.

\bibitem{LZZ}
Z. Lei, Q. Zhang, and N. Zhao,
\emph{Improved Liouville theorems for axially symmetric Navier--Stokes
equations},
Sci. Sin. Math. \textbf{47} (2017), no. 10, 1183--1198,
doi:10.1360/N012016-00149.

\bibitem{SereginSverak2009}
G. Seregin and V. \v{S}ver\'ak,
\emph{On Type I singularities of the local axi-symmetric solutions of the
Navier--Stokes equations},
Comm. Partial Differential Equations \textbf{34} (2009), no. 2, 171--201,
doi:10.1080/03605300802683687.

\end{thebibliography}

\end{document}
